\documentclass[11pt]{article}
\usepackage[margin=1in]{geometry}
\usepackage{amsmath,amssymb,amsthm,mathtools}
\usepackage[T1]{fontenc}
\usepackage{lmodern}
\usepackage{microtype}
\usepackage{algorithm,algpseudocode}
\usepackage[section]{placeins}
\makeatletter
\setlength{\@fptop}{0pt}
\makeatother
\newcommand{\wedkey}[2]{\mathsf{key}_{#1}(#2)}
\newenvironment{wedalgorithm}[2]{%
  \begin{algorithm}[!htbp]%
  \caption{#1}\label{#2}%
  \small\begin{algorithmic}[1]%
}{%
  \end{algorithmic}\end{algorithm}%
}
\setlength{\emergencystretch}{2em}
\usepackage{xurl}
\usepackage[colorlinks=true,linkcolor=blue,citecolor=blue,urlcolor=blue]{hyperref}
\newtheorem{theorem}{Theorem}[section]
\newtheorem{lemma}[theorem]{Lemma}
\newtheorem{proposition}[theorem]{Proposition}
\newtheorem{corollary}[theorem]{Corollary}
\theoremstyle{definition}
\newtheorem{definition}[theorem]{Definition}
\newtheorem{remark}[theorem]{Remark}
\newcommand{\WED}{\operatorname{WED}}
\newcommand{\conn}{\operatorname{conn}}
\newcommand{\Lip}{\operatorname{Lip}}
\newcommand{\ep}{\varepsilon}
\newcommand{\E}{\mathbb E}
\newcommand{\Prb}{\mathbb P}
\newcommand{\pos}[1]{\left(#1\right)_+}
\newcommand{\norm}[1]{\left\lVert#1\right\rVert}
\newcommand{\abs}[1]{\left|#1\right|}
\newcommand{\dd}{\mathrm d}
\newcommand{\cB}{\mathcal B}
\newcommand{\cA}{\mathcal A}
\newcommand{\cC}{\mathcal C}
\newcommand{\cQ}{\mathcal Q}

\newtheorem{observation}[theorem]{Observation}
\newcommand{\gap}{\bot}
\newcommand{\poly}{\operatorname{poly}}
\newcommand{\polylog}{\operatorname{polylog}}
\newcommand{\eps}{\varepsilon}
\newcommand{\1}{\mathbf 1}
\newcommand{\HH}{\mathcal H}
\newcommand{\QQ}{\mathcal Q}
\newcommand{\KL}{D_{\mathrm{KL}}}
\newcommand{\ip}[2]{\langle #1,#2\rangle}
\title{An Almost-Linear Approximation Scheme\\for Metric-Weighted Edit Distance}
\author{Jihan Wang}
\date{October 7, 2026}
\hypersetup{pdftitle={An Almost-Linear Approximation Scheme for Metric-Weighted Edit Distance},pdfauthor={Jihan Wang}}
\begin{document}
\maketitle
\begin{abstract}
We give an almost-linear randomized approximation scheme for edit
distance under arbitrary metric weights.
For every fixed rational $\varepsilon>0$, the algorithm returns an
estimate in $[D,(1+\varepsilon)D]$ with probability at least $5/6$,
where $D$ is the weighted edit distance.
Its worst-case expected running time is
$N^{1+O(\log\log\log N/\log\log N)}=N^{1+o(1)}$
for total string length $N$, with no dependence on the ratio between
positive costs.
The same running time gives factor $1+O((\log\log N)^{-3})$
when the accuracy is chosen from the input length.
The model permits constant-time exact metric queries and exact real
arithmetic. Finite binary computation and random sampling are charged.
We adapt the unit-cost refinement of OpenAI to metric weights.
Gap-mass coordinates control the cost of moving alignment endpoints.
A coarse estimator answers a substring query with source $X$ in
$(1+|X|)N^{o(1)}$ expected time, without a factor for the target
fragment's length. Approximate singleton comparisons contribute error
once, at the leaves. A single capped and rounded integer metric reduces
arbitrary real costs to the bounded-cost setting.
\end{abstract}

\section{Introduction}\label{paper:sec:introduction}

Edit distance is the minimum cost of changing one string into another by
insertions, deletions, and substitutions. It is a basic measure of string
similarity, with a long history in coding and sequence
alignment~\cite{Lev66,NW70,WF74,sellers1974}.
In the unit-cost setting, each nontrivial edit costs one.
Weighted edit distance allows the cost to depend on the edited symbols.
We study the metric setting: a metric $w$ is given on the alphabet
together with a gap symbol $\gap$. Deleting $a$ costs $w(a,\gap)$,
inserting $b$ costs $w(b,\gap)$, and substituting $a$ by $b$ costs $w(a,b)$.
The triangle inequality makes the edit-script definition equivalent to
the minimum cost of a monotone alignment path.

For strings $x,y$, let $N=|x|+|y|$ and $D=\WED_w(x,y)$.
Our goal is a $(1+\varepsilon)$ approximation in almost-linear time.
The weights may have arbitrary positive ratios.
This setting differs substantially from unit costs.
Moving a target endpoint by one position can cross an arbitrarily
expensive symbol. A long target fragment can also consist mostly of
symbols whose gap costs are very small.
An efficient algorithm must measure alignment errors by their cost,
while charging its computation by the number of input symbols.

Recent approximation schemes for unit-cost edit distance make this
question particularly timely.
Mao and Rubinstein~\cite{MR26} obtained a $(1+\varepsilon)$ approximation
in quasi-strongly subquadratic time.
OpenAI~\cite{openai} subsequently gave an almost-linear randomized
approximation scheme.
For arbitrary metric weights, the recent algorithms of
Mader, Tavasoli, and Wang~\cite{mtw} and
Das, Kipouridis, and Kociumaka~\cite{dkk} achieve factor
$3+\varepsilon$ in strongly subquadratic time.
We extend the unit-cost refinement to arbitrary metric weights and
obtain a $(1+\varepsilon)$ approximation in almost-linear expected time.

The metric assumption is essential for general approximation with
oracle access to costs. For growing alphabets, Gorbachev~\cite{gor2026}
proved an unconditional $\Omega(N^2)$ expected oracle-query lower bound
for approximating nonmetric minimum-cost alignment within any factor
depending only on $N$. The obstruction persists under symmetry alone
or the triangle inequality alone. Particular nonmetric classes can
still admit fast algorithms, including symmetric costs over a binary
alphabet~\cite{gor2026}.

\subsection{Our result}

The algorithm uses a metric oracle and exact real arithmetic.
Metric queries and real addition, subtraction, multiplication, division,
and comparison take constant time. We charge the finite binary
computation and random sampling used by the algorithm.
The recent arbitrary-metric algorithms also assume constant-time
metric queries and exact cost arithmetic~\cite{mtw,dkk}.
For a string $Z$, write
$M_w(Z)=\sum_{a\text{ in }Z}w(a,\gap)$, where the sum counts occurrences.

\begin{theorem}\label{wrapper:thm:full}
There is one randomized algorithm in the metric-oracle model of
Section~\ref{paper:sec:model} that takes strings $x,y$ and a rational
accuracy $\varepsilon>0$ as input. It returns an estimate $\widehat D$
satisfying
\[
 \Prb\{D\le\widehat D\le(1+\varepsilon)D\}\ge5/6.
\]
For every fixed accuracy, its worst-case expected running time is
\[
 N^{1+O(\log\log\log N/\log\log N)}=N^{1+o(1)},
\]
uniformly over the input strings and metric weights.
It uses polynomial workspace and terminates on every sequence of random
bits. It returns zero exactly when $x=y$, and its output always belongs
to $[0,M_w(x)+M_w(y)]$.
\end{theorem}

The running time has no dependence on the ratio between positive edit
costs. It includes executions on which an intermediate estimate fails.
The running-time bound is asymptotic. The large-input analysis requires
$\chi=\log_2H\ge2^{1000}$, where $H=\Theta(\log_2N)$.
Thus the cutoff is roughly $\log_2\log_2N\ge2^{1000}$.
Additional cutoff conditions depend on the accuracy. Below the cutoffs,
the algorithm uses exact dynamic programming. Its finite accuracy
encoding is specified in Section~\ref{paper:sec:model}.

The accuracy can also be chosen from the input length.

\begin{corollary}\label{wrapper:cor:vanishing}
There is a randomized algorithm that, for strings of total length
$N\ge4$, returns an estimate $\widehat D$ satisfying
\[
 \Prb\left\{D\le\widehat D\le
 \left(1+\frac{30}{(\log_2\log_2N)^3}\right)D\right\}\ge5/6.
\]
Its worst-case expected running time is
\[
 N^{1+O(\log\log\log N/\log\log N)},
\]
uniformly over the input strings and metric weights. The constant in
the running-time bound is absolute.
\end{corollary}

We construct a coarse table for polynomially bounded integer metric costs.
After shared preprocessing, a query with source $X$ has expected
evaluation cost $(1+|X|)N^{o(1)}$, uniformly in the target fragment.
The same bound holds conditional on the history before the entry's
first request. Gap-mass geometry and heavy-symbol methods from the
metric-weighted literature~\cite{mtw,dkk} provide the starting tools.
Nested nets of actual target boundaries control endpoint rounding.
The weighted connection inequality bounds the parent distance by the
sum of the child distances and their connection cost.

The refinement uses approximate singleton estimates that are
$3$-Lipschitz in endpoint distance.
Their total error is charged once at the comparison-tree frontier.
A single capped and rounded integer metric transfers the construction
to arbitrary real costs and keeps the work independent of positive
cost ratios.
These constructions implement the interfaces needed by OpenAI's
unit-cost refinement~\cite{openai}.\footnote{We cite the September 24,
2026 manuscript at the repository revision listed in the bibliography.
The catalogue describes a \href{https://github.com/openai/math/blob/adc7f1241b42e322a6451854ab7e4b4c146bf78a/lean/docs/121.md}{Lean formalization}
whose stated scope covers unit-cost edit distance on explicitly stored
integer strings with polynomially bounded symbols. For fixed rational
$\varepsilon\in(0,1)$, that scope gives success probability at least
$2/3$ and expected work at most $(N+2)^{1+\eta}$ for each fixed
$\eta>0$ and sufficiently large $N$. This is the qualitative
$N^{1+o(1)}$ bound. The scope does not cover our explicit rate or the
weighted construction in the exact-real oracle model.}
That refinement improves a coarse table through a hierarchy of local
comparisons. We adapt its curvature, monotone-loss, and first-crossing
arguments to the weighted setting and give proofs of the statements
used here. The weighted theorem does not use the unit-cost theorem as
a black box.

We retain the parameter forms of the unit-cost refinement~\cite{openai},
with $H=\Theta(\log N)$ and $\chi=\log_2H$.
The choices $p_*=\chi$ and $\ell=\Theta(\chi^3)$ give the explicit
subpolynomial overhead in Theorem~\ref{wrapper:thm:full}.
The cited unit-cost theorem states $N^{1+o(1)}$ without this rate.
The same explicit expected-time bound also applies to unit-cost edit
distance on explicitly stored strings with $O(\log N)$-bit integer
symbols on the word RAM with $O(\log N)$-bit words
(Remark~\ref{integer:rem:unitcost}).

The weighted scheme improves both approximation and asymptotic running
time for arbitrary metric costs.
For equal-length strings, Kuszmaul~\cite{kuszmaul2019} gave an
$O(N^\delta)$ approximation in $\widetilde O(N^{2-\delta})$ time
for each fixed $0<\delta<1$.
Mader, Tavasoli, and Wang~\cite{mtw} obtained factor $3+\varepsilon$ in
$\widetilde O(N^{7/4}/\varepsilon^8)$ time.
Das, Kipouridis, and Kociumaka~\cite{dkk} improved this to
$\widetilde O(N^{8/5}/\varepsilon^{16/5})$.
Gorbachev~\cite{gor2026} subsequently achieved factor $3$ in
$\widetilde O(N^{40/21})$ time.

\subsection{Related work}

\paragraph{Exact alignment and small distance.}
The classical dynamic program computes weighted edit distance in
quadratic time~\cite{NW70,WF74,sellers1974}.
Linear-space methods show how to recover alignments without storing the
whole dynamic-programming table~\cite{Hir75,MM88}.
Masek and Paterson~\cite{MP80} obtained a logarithmic improvement for
a fixed alphabet and commensurate edit costs.
Crochemore, Landau, and Ziv-Ukelson~\cite{clz2003} gave a
subquadratic alignment algorithm for unrestricted scoring matrices
over a constant-size alphabet.
For unit costs, small-distance algorithms exploit the fact that a
low-cost alignment stays close to the diagonal.
This principle underlies Ukkonen's algorithm~\cite{Ukk85}, the
Landau--Vishkin algorithms~\cite{LV88,LV89}, and incremental string
comparison~\cite{LMS98}.
Myers's difference algorithm~\cite{Mye86} gives another classical
distance-sensitive approach for insertion and deletion distance.
Related distance-sensitive methods have also been used for DNA
alignment under specific scoring rules~\cite{ZSSM00}.
These algorithms supply exact answers when the relevant distance
parameter is small. Their running times do not give an almost-linear
bound on unrestricted inputs.

Backurs and Indyk~\cite{BI18} showed that exact unit-cost edit distance
has no strongly subquadratic algorithm under the Strong Exponential Time
Hypothesis. Bringmann and K{\"u}nnemann~\cite{BK15} established related
conditional lower bounds for string similarity and dynamic time warping.
Abboud, Hansen, Vassilevska Williams, and Williams~\cite{AHWW16}
connected sufficiently large logarithmic improvements for exact edit
distance and LCS with new circuit lower bounds.
These results motivate approximation. Exact-distance hardness does not
itself exclude a fixed-factor or a $(1+\varepsilon)$ approximation.

\paragraph{Approximation for unit costs.}
The early subquadratic and near-linear algorithms traded approximation
quality for running time~\cite{BJKK04,BES06}.
Metric embeddings provided another route: Ostrovsky and
Rabani~\cite{OR07} embedded edit distance into $\ell_1$ with
subpolynomial distortion. Andoni and Onak~\cite{AO12} obtained a
near-linear algorithm with a comparable approximation factor.
For every fixed $\delta>0$, Andoni, Krauthgamer, and
Onak~\cite{AKO10} achieved a polylogarithmic factor in
$N^{1+\delta}$ time. Their source decomposition and precision sampling are
important predecessors of the coarse estimator used here.

Chakraborty, Das, Goldenberg, Kouck{\'y}, and Saks~\cite{CDGKS20}
gave the first classical constant-factor approximation in strongly subquadratic
time. Goldenberg, Rubinstein, and Saha~\cite{GRS20} developed stronger
tradeoffs, including a $(3+\varepsilon)$ approximation.
Fox, Mader, Quanrud, Tavasoli, and Wang~\cite{FMQTW26} gave a
distance-sensitive $(3+\varepsilon)$ approximation in
$\widetilde O_\varepsilon(N^{6/5}k^{2/5})$ time, where $k$ is the
unit-cost edit distance.
Brakensiek and Rubinstein~\cite{BR20} and Kouck{\'y} and
Saks~\cite{KS20} obtained near-linear constant-factor approximations
for sufficiently distant pairs.
Andoni and Nosatzki~\cite{AN26} gave a constant-factor approximation
on all input pairs in $N^{1+\delta}$ time for every fixed $\delta>0$.
The factor depends on $\delta$.
Approximation schemes require more than these constant-factor
guarantees. The recent results of Mao and Rubinstein~\cite{MR26}
and OpenAI~\cite{openai} reach factor $1+\varepsilon$ for unit costs.
The latter refinement is the direct starting point of this paper.

\paragraph{Weighted and bounded edit distance.}
Weighted edit distance also admits exact algorithms parameterized by
the distance or the weight range.
Here the weights are normalized so that every nontrivial edit costs at
least one, and $k$ is a bound on the distance.
Das, Gilbert, Hajiaghayi, Kociumaka, and Saha~\cite{dghks2023}
gave an $O(N+k^5)$ algorithm.
Cassis, Kociumaka, and Wellnitz~\cite{ckw2023} obtained
$\widetilde O(N+\sqrt{Nk^3})$ time.
Gorbachev and Kociumaka~\cite{gk2025} obtained improved algorithms
for small integer weights. They also gave a $(1+\varepsilon)$
approximation for normalized real weights in
$O(N+(k/\varepsilon)^{5/2}\log^3 N)$ time.
This distance-sensitive bound differs from the length-only bound sought
here.

Dynamic versions reuse work across updates~\cite{ckm2020,bgk2025}.
Weighted pattern matching asks where a pattern approximately occurs
inside a longer text~\cite{ckw2025}.
Structured geometric costs permit additional exact speedups, as shown
by Gold and Sharir~\cite{gold_sharir2018} for geometric edit distance
and dynamic time warping.
Our metric is given by an oracle and need not have such a geometric
representation.

For arbitrary metrics, Kuszmaul's simplification deletes inexpensive
symbols at a sampled cost scale~\cite{kuszmaul2019}.
The recent factor-$3+\varepsilon$ algorithms~\cite{mtw,dkk}
combine this idea with weighted local comparisons and cost-based
decompositions.
Our normalization also uses this simplification.
The main weighted refinement keeps the original input occurrences and
measures endpoint movement by their gap mass in the normalized metric.
This is needed when a large gap cost occurs beside a very cheap
substitution.

\paragraph{Queries, sketches, and related similarity problems.}
Sublinear algorithms distinguish promised distance gaps using random
access to the input strings~\cite{BEKMRRS03,GKS19,KocSaha20,BCR20}.
Recent adaptive algorithms narrow the gap to
$(k,k^{1+o(1)})$~\cite{BCFN22,BCFK24}, whereas non-adaptive queries
have a separate optimal tradeoff~\cite{GKKS22}.
Sketching algorithms encode each string separately and recover the
exact distance, and sometimes an optimal edit script, when the distance
is at most a prescribed threshold~\cite{BZ16,KPS21,BhatKoucky23,KS24}.
Streaming algorithms restrict input access to forward scans.
Randomized embeddings and low-distance algorithms also work with this
access restriction~\cite{CGK16,BZ16,KPS21}.
Independent preprocessing can reduce the cost of repeated
comparisons~\cite{GRS20}.
Our algorithm receives both strings explicitly and accesses metric costs
through an oracle.
Related questions arise for LCS~\cite{BK15,MR26} and dynamic time
warping~\cite{BK15,kuszmaul2019}.
Reductions between edit-distance variants require care with alphabets
and additive baselines~\cite{BDGK24}.

\section{Technical overview}\label{paper:sec:overview}

The proof has one central task: make the unit-cost refinement work when
alignment cost and string length describe different scales.
The refinement needs geometric control of candidate alignments and fast
local distance estimates.
It then combines those estimates through a source hierarchy.
With arbitrary weights, a position-based geometric estimate can be
wrong by an unbounded factor, and a local query can take time
proportional to a much longer target fragment.
We explain how to restore these two properties and how they support the
refinement. The exact parameter choices are deferred to the formal
construction.

\subsection{One bounded metric for the whole proof}

The numerical range must first be separated from the number of input
symbols. Searching successive cost scales between the smallest and
largest positive weights would make the running time depend on their
ratio. Instead, we find one interval containing the distance, with a
width controlled by $N$, and normalize the entire metric once.

Let $g(a)=w(a,\gap)$ be the gap cost of a symbol.
Deleting all occurrences with gap cost at most a threshold gives
simplified strings.
Kuszmaul's simplification~\cite{kuszmaul2019}, as used by
Das, Kipouridis, and Kociumaka~\cite{dkk}, has two useful features.
Removing cheap symbols costs at most the number of removed occurrences
times the threshold.
For a threshold sampled in a constant-width interval of scales, the
expected distance of the surviving strings is controlled by $D$.
If all surviving gap costs exceed a tested scale, any alignment below
that scale has no gap operations and can be tested by pairing equal
positions.

These tests search a finite set of input gap costs and their halves.
The resulting bracket $[L_{\rm br},U_{\rm br}]$ satisfies
\[
 0<L_{\rm br}\le U_{\rm br},\qquad
 U_{\rm br}/L_{\rm br}\le60N,\qquad D\le U_{\rm br}.
\]
The upper endpoint is valid on every outcome.
With probability at least $99/100$, the lower endpoint also satisfies
$L_{\rm br}\le D$.
The upper endpoint follows from a feasible alignment of the original
strings. It remains valid when a randomized scale test misses a close
pair. This allows every later computation to use a bounded metric.

Set $\eta_{\rm norm}=\varepsilon L_{\rm br}/(3N)$ and define
\[
 \bar w(a,b)=
 \left\lceil
 \frac{\min\{w(a,b),2U_{\rm br}\}}{\eta_{\rm norm}}
 \right\rceil.
\]
Both capping and upward rounding preserve the triangle inequality.
Capping at $2U_{\rm br}$ preserves the distance exactly.
An optimal original alignment uses no edge above $D\le U_{\rm br}$.
Any alignment affected by the cap has an edge of capped cost
$2U_{\rm br}>D$, so it cannot improve the optimum.
An alignment has at most $N$ edges, and rounding adds at most
$N\eta_{\rm norm}$. Consequently,
\[
 D\le\eta_{\rm norm}\WED_{\bar w}(x,y)
 \le D+\varepsilon L_{\rm br}/3.
\]
On the good bracket event, the final term is at most
$(1+\varepsilon/3)D$.

Every positive normalized cost is bounded by
$\lceil360N^2/\varepsilon\rceil$.
For each fixed accuracy and sufficiently large $N$, this is at most
$N^3$, and the total gap mass is at most $N^4$.
We evaluate a normalized query by binary search over bounded integers,
without taking the floor of an arbitrary real cost.
The coarse estimator and every refinement pass use this same
normalized metric. Their geometry and their distance estimates refer
to the same edit costs.

\subsection{Cost coordinates and fast singleton comparisons}

For the rest of the overview, $w$ denotes the bounded integer metric,
and $D$ denotes the distance in that metric.
With $g(a)=w(a,\gap)$, define its prefix gap masses by
\[
 M_x(i)=\sum_{k\le i}g(x_k),\qquad
 M_y(j)=\sum_{k\le j}g(y_k),\qquad
 R=M_x(|x|)+M_y(|y|).
\]
An alignment vertex $(i,j)$ has shift $M_y(j)-M_x(i)$.
The reverse triangle inequality gives
$|g(x_i)-g(y_j)|\le w(x_i,y_j)$.
Thus a diagonal step changes the shift by at most its substitution cost.
A gap step changes it by exactly that step's cost.
The total variation of shift along any alignment segment is bounded
by the segment's cost~\cite{dkk}.
This is the weighted replacement for the position-based shift used in
the unit-cost refinement.

String positions are numbered from one, and boundaries from zero.
A state compares a source interval $I=(l,u]$ with a target endpoint pair
$q=(q^-,q^+)$, where $q^-$ and $q^+$ are boundary indices.
Write $x[I]=x_{l+1}\cdots x_u$,
$y[q]=y_{q^-+1}\cdots y_{q^+}$, and
$c(I,q)=\WED_w(x[I],y[q])$.
The endpoint distance between two such pairs is
\[
 \dd(q,r)=
 |M_y(q^-)-M_y(r^-)|+|M_y(q^+)-M_y(r^+)|.
\]
Adding or deleting the target pieces between corresponding endpoints
shows that $c(I,\cdot)$ is $1$-Lipschitz in $\dd$.
We also need to combine child comparisons whose target intervals do
not meet. Their connection cost is the gap mass crossed between
successive endpoints. A connection inequality bounds the parent
distance by the child distances plus this charge.
It applies even when the child intervals overlap.
Each added source-deletion cost is assigned to a distinct backward
target crossing. The added cost is bounded by the connection charge.

Rounding must respect this geometry.
A grid of target indices can jump across a very expensive symbol.
We instead use nets of actual target boundaries, nested across dyadic
mass scales. Nearby net points are separated in gap mass, and each
boundary rounds downward with a controlled mass displacement.
A large gap between consecutive prefix masses contributes a boundary
to the net. We never split a target symbol into fractional positions.
The same separation gives short local candidate lists, while nesting
lets different stages round to compatible boundaries.

The second local difficulty is comparison against a singleton source.
For a source symbol $a$ and a target fragment $V$, the exact cost is
\[
 g(a)+M_w(V)
 -\max\bigl(
 \{0\}\cup
 \{g(a)+g(b)-w(a,b): b\text{ occurs in }V\}
 \bigr).
\]
The maximum is the saving obtained by pairing $a$ with one target
occurrence rather than deleting and inserting both symbols.
For unit costs, an occurrence query finds this saving.
For arbitrary metrics, it may require many different substitution
queries. Scanning $V$ at every source leaf would lose the desired
running time.

The heavy-symbol reduction~\cite{dkk} avoids that scan.
Keep the occurrences of $V$ with the largest gap costs and query their
substitution costs. If a potentially better occurrence is omitted,
its possible saving is bounded by twice its gap cost.
The retained occurrences then force enough remaining insertion cost
to absorb the missed saving.
For a chosen leaf accuracy $\varepsilon_{\rm leaf}\in(0,1]$, the resulting
estimate $Q$ satisfies
\[
 c\le Q\le(1+\varepsilon_{\rm leaf})c,
 \qquad
 |Q(a,V)-Q(a,V')|\le3\dd(V,V').
\]
It uses only $O(1/\varepsilon_{\rm leaf})$ substitution queries,
after range-reporting preprocessing.
The Lipschitz bound is as important as the approximation bound:
later comparisons move target endpoints.
We retain these actual approximate leaf values throughout the proof.

\subsection{A coarse table charged by source length}

Refinement starts from a coarse estimate at every state it may request.
Its requests are adaptive, so an estimate only for the full strings
would not suffice. We need one random table that is accurate
simultaneously on the predetermined state family.
We also need to evaluate a requested entry quickly when its target
fragment is much longer than its source fragment.

For a source fragment $X$, build a balanced tree of its positions. A tree node records a mass shift.
A leaf uses that shift to predict the mass interval opposite its source
symbol. It chooses either a deletion or a substitution against an
allowed target occurrence, and pays for the discrepancy of the chosen
endpoints. Internal nodes connect the shifts of their children.
This gives a min-plus recurrence with a penalty for changing shift.

The recurrence is a geometric relaxation of alignment.
Any assignment of shifts and leaf operations can be converted to an
alignment after paying its connection charges and the difference
between the total source and target gap masses.
Conversely, an optimal alignment supplies a shift at every source cut.
The mass-variation bound controls the total penalty of those shifts
over the tree.
The resulting comparison factor depends on tree branching and depth,
rather than on the weights or on the target length.
This is the weighted analogue of the source hierarchy used in the
unit-cost coarse estimator~\cite{AKO10,openai}.

Evaluating the recurrence on every integer shift would still be too
expensive. Precision sampling gives each node an additive error
allowance. A larger allowance permits a coarser shift grid and cheaper
leaf evaluation. A smaller allowance reveals more detail.
At a singleton, a target symbol whose gap cost is small relative to
the allowance can be replaced by an empty-target candidate with a
controlled additive error. Only the heavier occurrences need actual
substitution queries.

The deterministic restriction on allowed leaf starts confines these
occurrences to a mass interval of length proportional to the distance
guess. It contains only $O(1+k/a)$ occurrences with gap cost above a fixed
fraction of the allowance $a$, where $k$ is the guess.
This count does not depend on the number of cheap symbols in the
fragment.
Range reporting finds the heavy occurrences.
Piecewise-linear envelopes evaluate their leaf costs on the grid,
and two sweeps evaluate the min-plus shift transforms.
These operations produce costs for many shifts without scanning the
target fragment.

Child estimates are reconstructed through the tree using these random
allowances.
The analysis fixes a node's incoming allowance before exposing its
children's fresh draws. It then proves the same approximation
invariant on every relevant grid point.
Parent and child grids need not be nested.
Their transform error is charged to the chosen mesh sizes.
A union bound over the predetermined entry family gives a simultaneous
coarse table with factor $A_{\rm init}$.

Each entry uses its own random stream and completed shared
preprocessing.
Before an entry's first request, its stream is unexposed.
Thus its expected evaluation cost remains bounded after conditioning
on the request history.
The expected inverse allowance grows by a polylogarithmic factor at
each source level. The growing branching factor shortens the tree
enough to keep the accumulated overhead subpolynomial.
These bounds give expected work $(1+|X|)N^{o(1)}$ for an entry with source $X$,
uniformly in the target fragment.
This source-length bound is the computational counterpart of the
gap-mass geometry.
It lets the refinement ask for many different target comparisons while
charging the total time through the source tree.

\subsection{Why refinement can concentrate on narrow bands}

Fix a complete good coarse table $U$, with
$c\le U\le Ac$ at every state.
One refinement pass reduces this factor by a growing parameter $F$.
Put $a=\max\{1,A/F\}$.
Scaled coarse estimates and endpoint cones supply initial values
between $c/a$ and $O(Fc)$.
A cone adds the gap-mass cost of moving its center's endpoints.
The minimum of these cones provides controlled estimates at nearby
states as well.

Repeated child comparisons decrease these values.
Their average provides a stable cost profile for the next stage.
The refinement uses this profile to predict how an alignment's shift
evolves as cost accumulates.
Searching all child endpoint combinations would be expensive.
We instead search near this prediction. To justify that restriction,
we must bound both the prediction error and the cost of nodes where
the prediction fails.

To explain this, fix an optimal alignment for analysis and split it
recursively at the source cuts.
Round its target anchors to the mass nets.
The child comparisons form a comparison tree.
Their costs sum to the parent cost, apart from a controlled endpoint
rounding allowance.
At one nonterminal, let $m_i$ be the mean cost estimates of its
comparison children, let $z_i$ be their mass imbalances, and put
$W=\sum_i m_i$.
For $m_i>0$, the ratio $z_i/m_i$ describes shift change per unit of
estimated cost. If $m_i=0$, the imbalance bound gives $z_i=0$;
we take $0/0$ to be zero. The comparison tree stops at physical leaves
and at states of true cost zero. The comparison bounds ensure
$W>0$ at a nonterminal.

The main obstruction is variation in these ratios.
It is measured by
\[
 K=\sum_i\frac{z_i^2}{m_i}
   -\frac{(\sum_i z_i)^2}{W}.
\]
This quantity is nonnegative by Cauchy--Schwarz.
It vanishes when the child ratios agree.
We call $K$ the curvature. The average shift change per unit of
estimated cost is the velocity $v=(\sum_i z_i)/W$.
The same inequality gives
\[
 \sum_i|z_i-vm_i|\le\sqrt{KW}.
\]
Thus small $K$ makes every prefix shift close to a common velocity
times its accumulated mean cost.
There is also a discrepancy $G$ between the comparison-child sum and
the cheapest wide comparison. A wide comparison can use different
target endpoints at a source cut and pays their connection cost.
The computed prefix minimum lies between the comparison prefix's mean
cost minus $G$ and that mean cost.
Small $G$ lets the algorithm use this computed profile in place of the
unknown comparison profile.

The differences between a node's energy $z^2/m$ and the sum of its
children's energies telescope over the comparison tree.
Together with the averaged child comparisons, this bounds the sum of
$K+G$ over all nonterminals.
The algorithm need not know which nodes have low curvature.
The analysis uses the bound to charge the nodes where the narrow-band
prediction fails. If a node instead has nearly extreme velocity,
little decrement remains available and a separate bound charges it.
At the other nodes, the comparison action lies in the algorithm's
narrow band.

A band consists of child comparisons with shared target endpoints.
Its defining constraint uses the computed prefix-cost profile.
Away from extreme velocity, that constraint is contractive, so the
allowed shifts at each cut have a small diameter.
This small diameter lets the algorithm search a structured region even
though the full set of child actions is large.
This curvature-and-band mechanism is inherited from the unit-cost
refinement~\cite{openai}.
The weighted geometry makes its endpoint and connection estimates
valid for arbitrary metric costs.

Approximate leaves require a change in the telescope.
At a terminal, the table value is $Q$, while the true cost is $c$.
The energy is still bounded by $c$, since $Q\ge c\ge|z|$.
The telescoping sum of mean costs must use the actual terminal value
$Q$. The terminal true costs sum to $O(D)$, so replacing them by
$(1+\varepsilon_{\rm leaf})$ approximations adds only
$O(\varepsilon_{\rm leaf}D)$.
The leaf error is charged once at the frontier, rather than at every
level of the source hierarchy.

\subsection{Choosing comparisons and reducing their error}

Two sampling steps have different roles.
Shared samples compare canonical rounded actions for nearby
representatives in one mass cell.
They group children according to dyadic upper bounds on their costs.
Larger cost bounds receive larger sample budgets.
The winning representative is then evaluated at its actual current
cost. This gives a candidate that is reliable throughout the cell.

The online step chooses one whole band action from a distribution
guided by earlier child values.
The entropy objective balances those values against uncertainty in
the choice of action.
Sampling a whole action preserves its shared-endpoint constraints.

The regret analysis uses probability vectors with fixed mass in each
child block. It normalizes a child's cost by its dyadic bound.
The normalized costs decrease with the refinement index.
The monotone-loss inequality from the unit-cost refinement~\cite{openai}
controls the excess
cost of the online choices in terms of these decrements.
The fixed block masses cancel a common decrement within a block.
Only variation between its candidate child states remains.
This is what makes the bound fit the narrow-band geometry.
We prove the fixed-block version used here.

Intermediate passes need a further argument.
When $a>1$, the refinement's lower bound is $c/a$.
A child estimate can be below its true cost, so a comparison action
cannot always be treated as an upper estimate.
The baseline is the minimum of endpoint cones for the current group
candidates and the wide comparison. The update clips the online
candidate between this baseline and a slightly smaller fraction of it.
It preserves temporal monotonicity and the lower bound $c/a$.

For analysis, replace each current value $d_t$ by
$\overline d_t=\max\{c,d_t\}$.
Fix the complete child tables and consider the first index at which
a band action, including its endpoint cone, falls below the true
parent cost.
Before this crossing, the monotone-loss bound controls the excess
cost of the sampled actions. The extra cost from clipping an update
above its sampled action is charged to the action's child-value
decrements.
At the crossing, successful shared sampling leaves only a small
clipping error.
After it, child monotonicity keeps the crossing action's
earlier-index Bellman cost below the true cost.
The clipped analysis value is then exactly $c$ and contributes no
further positive regret.
In the final pass, $a=1$ and every child value is an upper estimate;
no crossing is needed.

Summing these bounds over the comparison tree gives one pass with
\[
 d_S\ge c/a
 \quad\text{on every outcome},\qquad
 \E(d_S-c)_+\le c/\ell,
\]
where $\ell$ grows with input length.
The expectation conditions on the complete fixed good seed table.
The curvature budget handles the excluded nodes, the online bound
handles covered nodes, and the terminal estimate supplies the leaf
bias. Together they reduce the approximation factor.

Independent copies and medians turn this expected error bound into
another simultaneous good table.
The process repeats until the final pass has $a=1$.
At that point Markov's inequality gives the required root
approximation.
Complete-table conditioning justifies each transition even though
the implementation requests only selected entries.
It is distinct from the exposed-history conditioning used to charge
their evaluation cost.

\subsection{Why the total time stays almost linear}

The mathematical tables have a polynomial state domain, but evaluating
all their entries would be too expensive.
The implementation evaluates only requested values and caches each
under one complete key.
Requests for the same table entry share that key, even when they
come from different representatives or sampled actions.
Private random streams give these cached values the same joint
distribution as the complete tables used in the proof.

The decisive count distinguishes two kinds of recursive request.
A request at an earlier refinement index can have a relatively large
number of children.
A request at the same index must descend one physical source level
and has a much smaller number of children.
Action lists and entropy calculations use values already gathered
from earlier indices.
Current child values are requested only for sampled or exact group
terms and selected actions.
Candidate enumeration uses these gathered arrays, so it does not
multiply the recursive cost by the size of the full action list.

The evaluator orders records by rank, physical height, and phase.
Same-node auxiliary reads at equal rank decrease phase. For the dependency count,
expand auxiliary records into their scalar requests while retaining
their owner decisions. Every resulting edge decreases rank or descends
one physical source level at equal rank.
These two decreases bound every scalar dependency chain.
At a fixed physical depth, the branching of the source tree is
balanced by the shrinking source lengths.
This balance lets us sum the coarse table's $(1+|X|)N^{o(1)}$ entry
bound over requests with a leading factor of $N$.
All request counts hold without conditioning on a good seed, so
the expected-time bound also includes inaccurate intermediate tables.

For the quantitative estimate, let
$H=\Theta(\log N)$ and $\chi=\log_2H$.
The parameters stay fixed through a bounded-integer core run.
We use coarse branching factor $B=H^\chi$.
Its source tree has height $O(H/\chi^2)$, so precision allowances
accumulate over fewer levels. Accounting for both the allowances and
the branching gives overhead factor $2^{O(H/\chi+\chi^2)}$.

For refinement, we choose $\ell=\Theta(\chi^3)$.
This choice supports $p_*=\chi$ while keeping the error budget within
$O(1/\ell)$. It also reduces the number of refinement steps to
$S=H/\ell^2$, which controls the number of earlier-index dependencies.
A request at the same index descends the physical source tree, whose
height is $J=O(H/\chi)$. Each level contributes a factor
$\ell^{O(1)}$ beyond the source branching already charged to string
length. The combined overhead from these levels is
$2^{O(J\log\ell)}=2^{O(H\log\chi/\chi)}$.
Earlier-index dependencies across all passes, local action enumeration,
finite entropy optimization, and sampling fit within this bound.
Thus the core takes expected time
\[
 N\,2^{O(H\log\chi/\chi)}
 =N^{1+O(\log\log\log N/\log\log N)}.
\]
The finite optimizer uses rational approximations and samples from
a finite distribution with fair bits.
Its mixture selects an action; it does not average recursive costs
into the returned table values.
This keeps the bit lengths of those values under control.

Finally, condition on the history that selected the normalization
bracket. Every such history fixes a valid bounded integer metric.
Fresh core runs have the same approximation and expected-time bounds
uniformly for that metric.
On a good bracket, the normalization and core losses multiply to
at most $(1+\varepsilon/3)^2\le1+\varepsilon$.
A constant number of runs raises the conditional success probability,
and combining it with the good-bracket probability gives $5/6$.
A failed bracket still supplies a polynomial-range metric, so it does
not introduce dependence on the original weight ratios into the
expected running time.

\section{Preliminaries}\label{paper:sec:model}

\subsection{Input and computational model}\label{wrapper:sec:model}

The input consists of strings $x,y$ of total length $N=|x|+|y|$
and a deterministic metric oracle $w$ on their alphabet together with
a gap symbol $\gap$. The strings do not contain $\gap$.
The metric has nonnegative real values, with $w(a,b)=0$ exactly when $a=b$.
Deleting $a$, inserting $b$, and substituting $a$ by $b$ cost
$w(a,\gap)$, $w(\gap,b)$, and $w(a,b)$, respectively.
We write $\WED_w(x,y)$ for the minimum cost of an edit script.

A metric query takes constant time. The model also permits exact real
addition, subtraction, multiplication, division, and comparison.
Integer indices and finite rational control data are represented in
binary, and their actual bit operations are charged.
Random sampling uses independent fair bits.
The model provides no floor operation or digit extraction on an
arbitrary real register and no ideal random real draw.
Space used by real values is measured in registers.

Each input occurrence has an $O(\log(N+2))$-bit positional handle.
The oracle accepts two handles, or a handle and $\gap$.
It returns the distance between the underlying symbols.
Different handles can name the same symbol, so equality is tested
by a zero metric value.
The algorithm outputs a distance estimate.
The running-time bound assumes oracle access to costs; it does not
include reading an explicitly supplied quadratic distance table.

The accuracy input is a rational
$\varepsilon=p_{\rm acc}/q_{\rm acc}>0$, with coprime positive integers
$p_{\rm acc},q_{\rm acc}$.
For a positive integer $u$ with $k$ binary digits, define its code as
$0^{k-1}\operatorname{bin}(u)$.
The accuracy encoding concatenates the numerator and denominator codes.
For the fixed-accuracy running-time guarantee, constants and thresholds
can depend on $\varepsilon$.

\subsection{Intervals, masses, and alignments}

Characters are indexed from one and boundaries from zero.
For boundaries $a\le b$ of a string $Z$, write
$Z(a,b]=Z_{a+1}\cdots Z_b$, with $Z(a,a]$ empty.
For a source interval $I=(l,u]$ and a target endpoint pair
$q=(q^-,q^+)$, use the shorthands $x[I]=x(l,u]$ and
$y[q]=y(q^-,q^+]$.
These boundary conventions follow~\cite{openai}.

Following the gap-mass notation of~\cite{mtw}, put
\[
 g(a)=w(a,\gap),\qquad
 M(Z)=\sum_{a\text{ in }Z}g(a),\qquad
 M_Z(j)=M(Z(0,j]).
\]
The sum counts occurrences. We write $M_w$ when the metric must be
specified. Unqualified masses use the metric fixed in the current
construction. In a local query on fragments $X,V$, their prefix masses
$M_X(i),M_V(j)$ start at their own first boundaries.
For arbitrary intervals, put
\[
 c(I,q)=\WED_w(x[I],y[q]),\qquad
 D=\WED_w(x,y),\qquad R=M(x)+M(y).
\]

An alignment is a monotone grid path with deletion, insertion, and
substitution edges. Its cost is the sum of the corresponding edit costs.
An alignment gives an edit script of the same cost.
Conversely, trace the occurrences through an edit script.
Each surviving original occurrence contributes a diagonal.
Deleted original occurrences and surviving inserted occurrences contribute
gap edges. The triangle inequality replaces successive changes to an
occurrence by one direct operation of no greater cost.
An inserted occurrence that is later deleted can be omitted.
The surviving original occurrences retain their order, so these operations
form an alignment. Thus minimum alignment cost equals edit-script distance.

We use logarithms to base two unless $\ln$ is written.
For an integer $k\ge1$, $[k]=\{1,\ldots,k\}$.
Empty minima are $+\infty$.

\subsection{Exact dynamic programming}

For strings of lengths $n,m$, let $D(i,j)$ be the minimum alignment cost
for $x(0,i]$ and $y(0,j]$. Then
\begin{align*}
 D(0,0)&=0,\qquad D(i,0)=\sum_{h=1}^i g(x_h),\qquad
 D(0,j)=\sum_{h=1}^j g(y_h),\\
 D(i,j)&=\min\bigl\{
 D(i-1,j)+g(x_i),\ D(i,j-1)+g(y_j),\\
 &\hspace{38mm}D(i-1,j-1)+w(x_i,y_j)\bigr\}.
\end{align*}
Every nonempty alignment ends in one of these three operations.
Removing it gives the corresponding prefix alignment.
Conversely, each prefix alignment extends by that operation.
Induction on $i+j$ proves the recurrence.

The computation uses $O((n+1)(m+1))$ metric queries and exact arithmetic
operations, with polynomial-in-$\log(N+2)$ indexing overhead per cell.
Two rows suffice for storage.
Empty inputs and a source of length at most one take $O(N+1)$ operations.
We use this routine for the exact branches of the algorithm.

\section{Normalization}\label{paper:sec:normalization}

The exact routine handles inputs with $N<2$. In the constructions below,
assume $N\ge2$. We construct one bounded integer metric on the original
alphabet. Its distance approximates the original distance on the good
bracket event. The bounds on its costs hold on every outcome.

\subsection{A bounded-randomness bracket with a certified upper endpoint}

\begin{lemma}\label{wrapper:lem:nogap}
Suppose every character of strings $S,T$ has gap cost at least $b>0$.
Then $\min\{\WED_w(S,T),b\}$ can be computed in $O(1+|S|+|T|)$ operations.
If the lengths differ, return $b$. Otherwise, return
\[
\min\left\{b,\sum_jw(S[j],T[j])\right\}.
\]
\end{lemma}
\begin{proof}
Any alignment of cost strictly below $b$ contains no insertion or deletion.
It must pair equal positions. If no such alignment exists, both formulas
return the cap $b$.
\end{proof}

For $r>0$, write $Z_{>r}$ for the string obtained from $Z$ by deleting
every occurrence whose gap cost is at most $r$.

\begin{lemma}\label{wrapper:lem:simplify}
For every scale $b>0$ and an ideal uniform $r\in[b,2b)$, used only for analysis,
\[
\E_r\WED_w(x_{>r},y_{>r})\le3D.
\]
\end{lemma}
\begin{proof}
Project an optimal original alignment to the surviving characters.
Original gap edges do not increase in cost. For a diagonal pairing $a_1,a_2$,
put $\alpha=\min\{g(a_1),g(a_2)\}$ and $\beta=\max\{g(a_1),g(a_2)\}$.
Let $s=w(a_1,a_2)$. If both symbols survive, the projected cost is $s$.
If neither survives, it is zero. If only one survives, it is the heavier
symbol and the projected cost is $\beta$. The reverse triangle inequality
gives $\beta-s\le\alpha$. A mixed outcome has probability at most
$(\beta-\alpha)/b\le s/b$. If such an outcome has positive probability,
then $\alpha<2b$. Hence the expected projected diagonal cost is at most
\[
 s+\Prb(\text{mixed})\max\{\beta-s,0\}
 \le s+(s/b)\alpha\le3s.
\]
If mixing is impossible, the upper bound $s$ suffices. Sum over the fixed
original alignment. Its projection is a feasible filtered alignment, so
the expected filtered optimum is at most $3D$.
\end{proof}

\begin{lemma}\label{wrapper:lem:test}
Fix a scale $b>0$. There is a fair-bit test with the following properties:
\begin{enumerate}
\item Every \textup{\textsc{Close}} output certifies $D<3Nb$.
\item If $D<b/10$, then the probability of a \textup{\textsc{Far}} output is at most
      $1/(200N)$.
\item It uses deterministically $O(N\log N)$ metric/real operations and $N\operatorname{poly}(\log N)$ finite-control work. Its number
      of fair bits is $O(\log^2N)$.
\end{enumerate}
\end{lemma}
\begin{proof}
Put $q=\lceil\log_2N\rceil+2$. Draw a uniform integer
$J\in\{0,\ldots,2^q-1\}$ using exactly $q$ fair bits. Its cell is
\[
I_J=\left[b+\frac{Jb}{2^q},b+\frac{(J+1)b}{2^q}\right).
\]
If no input symbol has gap cost strictly inside $I_J$, take its left endpoint
as the threshold $r$. Otherwise take $r=b$. Sort the input gap costs
once and use predecessor searches to implement this test.
A gap cost at a cell's left endpoint has already been removed there;
one at its right endpoint is outside the half-open cell. Thus filtering
is constant on every cell declared constant.

At most $N$ cells contain a symbol's gap cost in their interior.
The probability of choosing such a cell is at most $N2^{-q}\le1/4$.
Couple the uniform cell with an ideal uniform point inside it.
On every other cell, the finite threshold and the ideal threshold
produce the same filtered strings.
If $D<b/10$, Lemma~\ref{wrapper:lem:simplify} gives expected filtered
distance less than $b/2$ under the ideal threshold.
Markov's inequality gives distance below $b$ with probability greater
than $1/2$. Removing the bad-cell probability leaves probability greater
than $1/4$ for the finite sampler.

For each sampled threshold compute
$h=\min\{\WED_w(x_{>r},y_{>r}),b\}$ by Lemma~\ref{wrapper:lem:nogap}.
Return \textup{\textsc{Close}} if any sample has $h<b$. Such a sample gives
\[
D\le h+M_w(x\setminus x_{>r})+M_w(y\setminus y_{>r})
<b+2Nb\le3Nb.
\]
Run $t=4\lceil\log_2(200N)\rceil$ independent samples. Since
$(3/4)^4<1/2$, their common failure probability is at most $1/(200N)$.
All logarithmic counts here are computed from bounded integer bit lengths.
Each sample costs $O(N+\log N)$ operations and exactly $q$ fair bits.
\end{proof}

\begin{lemma}\label{wrapper:lem:bracket}
For distinct nonempty strings, a fair-bit algorithm uses
$O(N\log^2N)$ metric/real operations and $N\operatorname{poly}(\log N)$
finite-control work, and returns $0<L_{\rm br}\le U_{\rm br}$ such that
\begin{align*}
&U_{\rm br}/L_{\rm br}\le60N\quad\hbox{on every outcome},\\
&D\le U_{\rm br}\quad\hbox{on every outcome},\\
&\Prb\{L_{\rm br}\le D\le U_{\rm br}\}\ge99/100.
\end{align*}
It uses $O(N)$ real registers, $\operatorname{poly}(\log N)$-bit control
data, and $O(\log^3N)$ fair bits. It uses no floor on real registers.
\end{lemma}
\begin{proof}
Sort the distinct candidate scales
\[
\mathcal S=\{g(a)/2,g(a):a\text{ occurs in }x\text{ or }y\}.
\]
Let $A=\min\mathcal S$ and $B=\max\mathcal S$.
Every original gap costs at least $2A$. Use Lemma~\ref{wrapper:lem:nogap} to
compute $e_0=\min\{D,2A\}$. If $e_0<2A$, return $[e_0,e_0]$.
This is a positive exact singleton because the strings are distinct.
Otherwise $2A\le D\le NB$.

While $A,B$ are not consecutive in $\mathcal S$, test a median candidate
strictly between them. Replace $A$ by that candidate on \textup{\textsc{Far}},
and replace $B$ by it on \textup{\textsc{Close}}. Every \textup{\textsc{Close}} answer is
certified. Each update preserves $D<3NB$, so this bound holds
throughout the search on every outcome.
For analysis, assign an independent random packet to every candidate
scale. The algorithm draws a packet only when it tests that scale.
There are at most $2N$ scales. By Lemma~\ref{wrapper:lem:test}, with
probability at least $99/100$, every possible \textup{\textsc{Far}} answer also
has its valid implication $A/10\le D$. On this event the search maintains
\[
A/10\le D<3NB.
\]
This union is over fixed candidate scales, so adaptive median choices
do not affect it.

If the final endpoints have $B\le2A$, return $[A/10,3NB]$.
The width is at most $60N$, its upper endpoint is valid by
the bound $D<3NB$, and its lower endpoint is valid on the stated event.

Otherwise $B>2A$. No input symbol has gap cost in $(A,2B)$: such a cost in
$(A,B)$ would itself be an intervening candidate, while a gap cost in
$[B,2B)$ would have an intervening half. The simplification at $A$ has
all surviving gap costs at least $2B$. Compute
\[
e=\min\{\WED_w(x_{>A},y_{>A}),B\},\qquad
L_{\rm br}=\max\{A/10,e/5\},\qquad U_{\rm br}=15NL_{\rm br}.
\]
The width is $15N$ on every outcome.
If $e=B$, then the bound $D<3NB$ gives $D<3NB=3Ne\le15NL_{\rm br}$.
If $e<B$, restoring deleted characters costs at most $NA$, so
\[
D\le e+NA\le5L_{\rm br}+10NL_{\rm br}\le15NL_{\rm br}.
\]
These two upper-bound arguments do not use the \textup{\textsc{Far}} implications.
They prove the asserted upper bound on every outcome.
For the lower bound, all thresholds in $[A,2A)$ give the same filtered
strings. Lemma~\ref{wrapper:lem:simplify} consequently gives $e\le5D$.
Together with the event $A/10\le D$, this proves $L_{\rm br}\le D$.

The number of scale tests is $O(\log N)$. Lemma~\ref{wrapper:lem:test} gives the
claimed deterministic running time and bit count. Thresholds and bracket
endpoints use only exact arithmetic, original real values, and bounded
integer multipliers or divisors.
\end{proof}

The bracket construction follows Das, Kipouridis, and
Kociumaka~\cite{dkk}. We retain the upper-endpoint guarantee from its
proof on every outcome. This guarantee allows normalization even when
the lower endpoint is inaccurate.

\subsection{One capped and rounded metric}\label{wrapper:sec:normalization}

\begin{lemma}\label{wrapper:lem:normalize}
Let $[L_{\rm br},U_{\rm br}]$ be any output of Lemma~\ref{wrapper:lem:bracket}. Define
\[
\eta_{\rm norm}=\frac{\eps L_{\rm br}}{3N},\qquad
\bar w(a,b)=\left\lceil\frac{\min\{w(a,b),2U_{\rm br}\}}{\eta_{\rm norm}}\right\rceil,
\qquad \bar D=\WED_{\bar w}(x,y).
\]
Then $\bar w$ is an integer metric with
\[
1\le\bar w(a,b)\le
\left\lceil\frac{360N^2}{\eps}\right\rceil\quad(a\ne b).
\]
On every bracket outcome,
\[
D\le\eta_{\rm norm}\bar D\le D+\eps L_{\rm br}/3.
\]
On the good bracket event $L_{\rm br}\le D$, this becomes
\[
D\le\eta_{\rm norm}\bar D\le(1+\eps/3)D.
\]
\end{lemma}
\begin{proof}
The function $x\mapsto\min\{x,2U_{\rm br}\}$ is nondecreasing and subadditive
on nonnegative reals, so capping preserves the triangle inequality.
Upward rounding after scaling preserves it because
$\lceil x+y\rceil\le\lceil x\rceil+\lceil y\rceil$.
Positive costs remain positive, and zero costs remain zero.
The deterministic width bound gives
\[
2U_{\rm br}/\eta_{\rm norm}=6NU_{\rm br}/(\eps L_{\rm br})\le360N^2/\eps.
\]

By Lemma~\ref{wrapper:lem:bracket}, $D\le U_{\rm br}$ on every outcome. An original optimal
alignment uses no edge more expensive than $D$, and is unchanged by
capping at $2U_{\rm br}$. Any alignment changed by capping has an edge of capped
cost $2U_{\rm br}>D$, so it cannot improve the optimum.
The capped distance equals $D$. A monotone alignment has at most $N$ edges.
Rounding increases the cost of an optimal capped alignment by at most
$N\eta_{\rm norm}=\eps L_{\rm br}/3$, while it never decreases any capped cost. This proves
the asserted transfer bound. The good-event claim follows from $L_{\rm br}\le D$.
\end{proof}

\paragraph{A fixed integer search range.}
Use exact dynamic programming whenever $\eps N<361$. In the remaining
branch,
\[
\frac{360N^2}{\eps}\le\frac{360}{361}N^3<N^3.
\]
Since $N^3$ is an integer, every normalized cost is at most $N^3$.
There is no need to calculate a real ceiling to find the search bound.
For one normalized query, obtain $c=\min\{w(a,b),2U_{\rm br}\}$ and binary-search
for the least integer $j\in\{0,\ldots,N^3\}$ with $c\le j\eta_{\rm norm}$.
This returns exactly $\bar w(a,b)$ using $O(\log N)$ exact real
comparisons and arithmetic operations. Its integer control data have
$O(\log N)$ bits. The construction never extracts digits of $c/\eta_{\rm norm}$.
To avoid an implicit integer-to-real conversion at each search step,
maintain exact real copies of the integer lower and upper bounds.
If their integer sum is odd, subtract $1/2$ from half their real sum;
otherwise use that half-sum directly. This is the real copy of the
binary integer midpoint. Update both copies together, adding one when
the lower bound advances. Initial real copies are obtained by scanning
the binary encoding of $N$ with repeated doubling, then cubing it.
The integer parity test uses finite binary control.
Each search step uses a constant number of permitted real operations.
The total normalized gap mass $\bar R$ is at most $N^4$ and at least $N$.


\section{Weighted geometry and local queries}\label{paper:sec:geometry}

Fix an integer metric with total input gap mass $R\le N^C$ for a fixed
constant $C$. Positive gap costs are at least one.
Endpoint coordinates count gap mass, but each endpoint is an original
character boundary. We use the gap-mass geometry to compare arbitrary
source and target fragments.

\FloatBarrier

\subsection{Two weighted geometric facts}

In the weighted alignment grid, a horizontal step inserts a target
symbol $b$ at cost $g(b)$. A vertical step deletes a source symbol $a$
at cost $g(a)$. A diagonal step pairs $a$ with $b$ at cost $w(a,b)$.
The metric assumptions imply that shortest alignment cost equals weighted
edit distance. The following gap-mass bound is also used by
Das, Kipouridis, and Kociumaka~\cite{dkk}.

\begin{lemma}\label{seed:lem:potential}
Along an alignment, the change in target prefix mass minus source prefix
mass has absolute value at most the cost of the corresponding subpath.
\end{lemma}
\begin{proof}
A gap edge changes this potential by its gap cost in absolute value. A
diagonal pairing $a,b$ changes it by $g(b)-g(a)$, whose absolute value is
at most $w(a,b)$ by the reverse triangle inequality. Sum along the subpath.
\end{proof}

\begin{lemma}\label{seed:lem:connection}
Let $X_1,\dots,X_r$ be consecutive source fragments, where $r\ge1$.
Fix an ambient target string $Y$ and a parent fragment $Y(a,b]$.
For each $i$, choose a fragment $V_i=Y(a_i,b_i]$, possibly empty,
and put $u_i=M_Y(a_i)$ and $v_i=M_Y(b_i)$.
The fragments may overlap or leave the parent fragment.
Then
\[
 \WED_w(X_1\cdots X_r,Y(a,b])
 \le \sum_i\WED_w(X_i,V_i)
       +|u_1-M_Y(a)|+\sum_{i<r}|u_{i+1}-v_i|
       +|M_Y(b)-v_r|.
\]
\end{lemma}
\begin{proof}
Place optimal child alignments in source order.
Join their target endpoints by monotone walks along target boundaries,
starting at $a$ and ending at $b$.
Each connecting walk pays the gap masses of the occurrences it crosses.
Its cost equals the absolute prefix-mass difference between its endpoints.

Select the first forward crossing of each occurrence in $Y(a,b]$.
These selected crossings occur in target order.
Keep every child deletion.
Keep a diagonal only when its target crossing is selected; otherwise
replace it by deletion of the same source occurrence.
Keep selected child insertions and omit the others.
For each selected connecting crossing, insert that target occurrence.
The resulting alignment consumes the source in order and emits exactly
$Y(a,b]$.

The concatenated walk crosses an occurrence inside the parent once
more forward than backward. For an occurrence outside the parent,
the two counts are equal.
After removing the selected crossings, the remaining forward crossings
can be paired bijectively with backward connecting crossings.

Replacing a diagonal $(x_i,y_j)$ by deletion increases cost by at most
$g(y_j)$, since $g(x_i)\le w(x_i,y_j)+g(y_j)$.
Charge this increase to the paired backward crossing.
Charge each selected connecting insertion to its own forward crossing.
These charges use distinct crossings and total at most the displayed
connection mass. Omitted child insertions only save cost.
\end{proof}

\subsection{Mass geometry and endpoint nets}

Define
{\small
\[
 z(I,q)=(M_y(q^+)-M_y(q^-))-(M_x(u)-M_x(l)),\qquad
 \dd(q,r)=|M_y(q^-)-M_y(r^-)|+|M_y(q^+)-M_y(r^+)|.
\]
\par}
The shift at a path point $(i,j)$ is $M_y(j)-M_x(i)$.

\begin{lemma}\label{refine:lem:geometry}
We have $c\ge|z|$. Both $c(I,\cdot)$ and $|z(I,\cdot)|$ are
$1$-Lipschitz in $\dd$.
Partition a source interval $I$ into $k\ge1$ consecutive intervals
$I_1,\ldots,I_k$, allowing empty intervals.
For a target endpoint pair $r$ and target states $s_1,\ldots,s_k$, define
\[
 \conn_r(s)=|M_y(s_1^-)-M_y(r^-)|+
 \sum_{i<k}|M_y(s_{i+1}^-)-M_y(s_i^+)|+
 |M_y(s_k^+)-M_y(r^+)|.
\]
Then
\[
 c(I,r)\le\sum_i c(I_i,s_i)+\conn_r(s),\qquad
 |z(I,r)|\le\sum_i|z(I_i,s_i)|+\conn_r(s).
\]
The variation of shift along any alignment segment is at most that segment's cost.
\end{lemma}
\begin{proof}
Lemma~\ref{seed:lem:potential} proves the shift-variation statement.
Applying it to an optimal alignment gives $c\ge|z|$.

An endpoint change inserts or removes the target symbols crossed and costs at most their total gap mass. The triangle inequality for edit distance proves the Lipschitz bound for $c$. The displayed formula for $z$ proves its bound.

The first connection inequality is Lemma~\ref{seed:lem:connection},
with ambient target $y$ and parent fragment $y[r]$.
For the second, the signed target-mass changes telescope.
The source children partition the parent, so
\[
 z(I,r)-\sum_i z(I_i,s_i)
 =M_y(r^+)-M_y(s_k^+)
  +\sum_{i<k}\bigl(M_y(s_{i+1}^-)-M_y(s_i^+)\bigr)
  +M_y(s_1^-)-M_y(r^-).
\]
The scalar triangle inequality bounds this difference by $\conn_r(s)$.
\end{proof}

\begin{lemma}\label{refine:lem:ordered-splitting}
Let $I=(l,u]$ and $k\ge1$. Fix prescribed source cuts
$l=l_0\le l_1\le\cdots\le l_k=u$ and put $I_i=(l_{i-1},l_i]$.
For every target pair $q$, there are boundaries
$q^-=v_0\le v_1\le\cdots\le v_k=q^+$ such that
\[
 c(I,q)=\sum_{i=1}^k c\bigl(I_i,(v_{i-1},v_i)\bigr).
\]
The cuts may repeat, so some source intervals may be empty.
\end{lemma}
\begin{proof}
Take an optimal parent alignment as a chronological sequence of vertices.
Fix its start and end times for cut occurrences zero and $k$.
For every internal cut occurrence, choose the first arrival at its source
row. These times are weakly ordered even when numeric cuts repeat.
Consecutive chosen times partition every path edge into valid child
alignments with original occurrence identities. Each subpath is optimal
for its child endpoints, since a cheaper replacement would reduce the
parent cost. Their costs sum to the parent optimum.
Only the outside cuts are pinned to the path's endpoints. Internal cuts
at the same source row share its first-arrival vertex. Horizontal edges
after that vertex belong to the child after the last internal cut at
that row. If the source is empty, all internal cuts use the start
vertex and the final child receives the entire horizontal path.
\end{proof}

\begin{lemma}\label{refine:lem:nets}
There are nested sets $\mathcal N_h$ of target boundaries, for dyadic integers $h\ge1$, such that $\mathcal N_1$ is the complete boundary set, consecutive points of $\mathcal N_h$ are separated by mass at least $h$, and downward rounding $\operatorname{pred}_h$ satisfies
\[
 0\le M_y(q)-M_y(\operatorname{pred}_h(q))<2h.
\]
At $h=1$ the error is zero.
After prefix preprocessing, all levels through mass $R$ require
$O(N\log R)$ comparisons and stored boundary records, with
$O_C(N\log^2(N+2))$ binary work and stored bits.
A mass interval of length $b$ contains $O(1+b/h)$ net points.
Reporting $k$ points uses $O(\log N)$ mass comparisons and $k$ record
accesses, with $O_C((\log N+k)\log(N+2))$ binary work.
\end{lemma}
\begin{proof}
Starting with $\mathcal N_1$, construct each $\mathcal N_h$ from
$\mathcal N_{h/2}$.
Retain boundary zero, then greedily select the first old-net point whose
mass is at least $h$ beyond the previous selected point.
This gives nesting and separation.

Let $p=\operatorname{pred}_{h/2}(q)$.
Induction gives $M_y(q)-M_y(p)<h$.
Let $r$ be the selected predecessor of $p$.
The greedy rule gives $M_y(p)-M_y(r)<h$; the difference is zero if $r=p$.
No old-net point lies between $p$ and $q$, so
$r=\operatorname{pred}_h(q)$.
Adding the two bounds proves the rounding error below $2h$.
This argument also covers an arbitrarily large gap between consecutive
prefix masses.
Each level is one scan. Binary search in the sorted net array gives
range reporting.
\end{proof}

\subsection{Integer leaf queries}

An empty source interval has exact cost equal to target gap mass. At a singleton source $a$, write $u=g(a)$ and $G=\sum_{b\in V}g(b)$ for the target interval $V$. Define the saving
\[
 B(a,b)=u+g(b)-w(a,b),\qquad 0\le B(a,b)\le2g(b).
\]
The exact cost is $u+G-\max(\{0\}\cup\{B(a,b):b\in V\})$.

Fix a positive integer $k$.
Retain the $k$ largest-gap positions in $V$, or all positions if fewer exist.
Among positions with equal gap cost, retain the earlier original
position first.
Define
\[
 Q_k(a,V)=u+G-\max\bigl(\{0\}\cup\{B(a,b):b\text{ is retained}\}\bigr).
\]
For an empty source, let $Q_k$ be the exact answer.
For target fragments $V,V^{\prime}$, $\dd(V,V^{\prime})$ denotes the distance between their endpoint pairs.

\begin{lemma}\label{refine:lem:leaves}
The value $Q_k$ is an integer in $[0,R]$, is zero exactly at zero-cost states, and satisfies
\[
 c\le Q_k\le(1+2/k)c,\qquad
 |Q_k(a,V)-Q_k(a,V')|\le3\dd(V,V').
\]
Preprocessing uses $O(N)$ comparisons and records, with
$O_C(N\log(N+2))$ binary work and stored bits, in addition to gap queries.
A query uses $O(k\log N)$ bounded-word operations and $O(k)$ metric calls.
Its binary control and arithmetic cost is
$O_C(k\log^2(N+2))$, in addition to the metric-query cost.
\end{lemma}
\begin{proof}
If the maximum saving is attained by zero or a retained position, the
value is exact. Otherwise, let $b$ be a maximizing position that is not retained.
All $k$ retained positions have gap cost at least $g(b)$,
so $g(b)\le G/(k+1)$. The exact cost satisfies
\[
 c=u+G-B(a,b)\ge G-g(b)\ge kG/(k+1),
\]
where the first inequality uses $B(a,b)\le u+g(b)$.
The missed saving is at most $2g(b)$.
Its ratio to $c$ is at most $2/k$.
A zero-cost singleton state has one matching target symbol, so it is exact.

Insert an endpoint symbol of gap cost $t$.
If it enters the retained set, it displaces at most one old symbol,
whose gap cost is at most $t$.
The maximum retained saving can then increase or decrease by at most
$2t$. Since $G$ increases by $t$, the change in $Q_k$ lies in $[-t,3t]$.
Deletion reverses this operation.
Telescoping the endpoint operations proves the regularity bound.

A range-maximum tree on target gap costs is built in linear time.
Its maximum records choose the earliest position among ties.
A priority queue reports an interval's maximum and replaces the
interval by its two remaining sides.
This gives the $k$ heaviest positions in $O(k\log N)$ time.
Das, Kipouridis, and Kociumaka~\cite{dkk} use this reporting structure.
Evaluate the reported positions' savings with $k$ metric calls.
\end{proof}

\begin{remark}
Use the line metric with coordinates $\mathtt c=-1$, $\bot=0$,
$\mathtt b=1$, and $\mathtt a=2$.
Let the ambient target be $Y=\mathtt c^k\mathtt b$.
Let $V$ be its suffix of length $k$, and let $V'=Y$.
For source $\mathtt a$, the estimator retains every position of $V$
and returns $Q_k(\mathtt a,V)=k$.
In $V'$, the earlier-position tie rule retains its first $k$ symbols,
all equal to $\mathtt c$, giving $Q_k(\mathtt a,V')=k+3$.
The endpoint distance is $\dd(V,V')=g(\mathtt c)=1$,
while the true distances are $k$ and $k+1$.
Thus the Lipschitz constant $3$ is sharp.
\end{remark}


\section{The weighted coarse estimator}\label{paper:sec:seed}

We construct a random coarse table for a fixed family of fragment pairs.
The table must be accurate simultaneously, although the evaluator can
request its entries adaptively. A separate work bound must remain valid
conditional on the history before an entry's first request. We do not
condition that work bound on approximation success.

The source hierarchy, precision reconstruction, and conditional
first-request accounting are adapted from~\cite{openai}.
The mass potential and heavy-character reporting follow
\cite{dkk}. The weighted implementation uses
allowance-dependent grids and leaf envelopes that do not scan the target
fragment.

\subsection{Model, parameters, and shared preprocessing}\label{seed:sec:statement}

Fix constants $C,s>0$. The strings $X_0,Y_0$ have total length at most
$N\ge2$. The metric $w$ is integer-valued on their alphabet together with
$\gap$, and every positive metric value belongs to $[1,N^C]$. A metric
query takes constant time and its bounded integer answer is used in binary.
If the answer is supplied in an exact real register, choose a fixed
integer $d\ge C$ and binary-search $\{0,\ldots,N^d\}$ for its value.
This recovers its binary representation with polylogarithmic overhead.
Maintain real copies of the integer search bounds, as in the query
implementation of Lemma~\ref{wrapper:lem:normalize}.
Each occurrence has an $O(\log N)$-bit handle
to its original input position. Different handles may name the same
symbol. Test equality by a zero metric query. Local prefix-mass offsets
do not change the handles supplied to substitution queries.

Write
\[
 g(a)=w(a,\gap),\qquad M(Z)=\sum_{a\text{ in }Z}g(a),\qquad
 R_0=M(X_0)+M(Y_0).
\]
The sum counts occurrences. For a positive integer $z$, let
$\operatorname{leastpow}_2(z)$ be the least power of two at least $z$.
Set
\[
 H=\operatorname{leastpow}_2\bigl(\lceil\log_2(N+R_0+2)\rceil\bigr),
 \qquad \chi=\log_2H,
\]
and
\[
 p_*=\max\{1,\chi\},\qquad B=H^{p_*},\qquad
 A_{\mathrm{init}}=BH^4,\qquad Q_0=H^4,\qquad
 \lambda=\frac1{64H}.
\]
These parameters use the full input, even for a shorter queried source.
They are computable by bit lengths and integer comparisons. Both $B$
and $Q_0$ are powers of two. Since $R_0\le N^{C+1}$, we have
$H=\Theta_C(\log N)$ and $B=N^{o(1)}$.
We analyze the fixed large-input branch $\chi\ge2^{1000}$, where
$p_*=\chi$ and $\log_2B=\chi^2$. In this branch,
$128H^5(1+\ln B)\le H^9$. Exact dynamic programming handles the remaining
finite input-size range.

Before exposing random bits, fix an indexed mathematical family
$\mathcal E$ of at most $N^s$ fragment pairs $(X,V)$ of $X_0,Y_0$.
Distinct indices may name the same pair. Assign each index its own
private stream of independent fair bits. Streams for different indices
and table instances are independent. The evaluator does not enumerate
the family. It allocates entries and samples their bits only as needed.
An entry may read only its full index, its private random bits,
its local workspace, and the completed deterministic shared data.
It does not call another entry or read another entry's auxiliary state.

Shared preprocessing queries every gap cost and stores both strings'
prefix masses. It also builds a segment tree on the target positions.
Each tree node stores a position of maximum gap mass in its interval,
with ties resolved by original target position. These structures require
$O(N\log N)$ bounded-word operations and $O(N)$ metric queries; their
charged binary work and storage are $N\operatorname{poly}(\log N)$.
The structures are read-only during entry evaluation.

For an entry, put $m=|X|$. Obtain a local target prefix mass by subtracting
the fragment's initial global prefix mass. This takes constant-time
accesses and does not scan $V$. Local source prefix masses require
$O(1+m)$ source work. If $X$ is nonempty, set
$K=\operatorname{leastpow}_2(M(X)+M(V))$ and use all guesses
$k=1,2,4,\ldots,K$. Thus $K\le2N^{C+1}$, and the guess set is fixed
before any computed estimate is used.

\subsection{A mass-shift tree with restricted leaf candidates}

Fix a nonempty source $X=x_1\cdots x_m$ and target $V=y_1\cdots y_n$.
Write
\[
 M_X(i)=\sum_{h=1}^i g(x_h),\qquad M_V(j)=\sum_{h=1}^j g(y_h),
 \qquad M_X(0)=M_V(0)=0,\qquad \Delta=|M_X(m)-M_V(n)|.
\]
Let $k\ge1$ be an integer power of two and put
$\mathcal D_k=[-k,k]\cap\mathbb Z$.

For character $x_i$, put $g_i=g(x_i)$. At shift $v\in\mathcal D_k$ its
predicted target mass interval is
\[
 [M_X(i-1)+v,M_X(i)+v].
\]
The following leaf envelope uses actual target operations:
{\small
\begin{align*}
 F_i(v)=\min\bigg\{&
 \min_{0\le j\le n}\left[g_i+|M_V(j)-M_X(i-1)-v|+|M_V(j)-M_X(i)-v|\right],
 \\[-2pt]
 &\min_{\substack{1\le j\le n\\ |M_V(j-1)-M_X(i-1)|\le k}}
 \left[w(x_i,y_j)+|M_V(j-1)-M_X(i-1)-v|+|M_V(j)-M_X(i)-v|\right]
 \bigg\}.
\end{align*}
\par}
The first minimum deletes $x_i$ and selects an actual target boundary.
The second pairs $x_i$ with an actual target occurrence.
The two absolute-value terms compare its fragment's mass endpoints
with the predicted interval.
If no singleton satisfies the start restriction, its minimum is
$+\infty$; the deletion minimum is always finite.
For fixed $(X,V,k)$, the candidate set does not depend on $v$.
Each candidate is $2$-Lipschitz in $v$, so the minimum $F_i$ has
the same bound.

For a source node $(l,u]$, set $c_i=l+\lfloor i(u-l)/B\rfloor$
for $i=0,\ldots,B$. Its children $(c_{i-1},c_i]$ are consecutive
and have widths differing by at most one. Apply these cuts recursively,
stopping at width zero or one. The maximum
depth is
\[
 D_0=\lceil\log_B\max\{1,m\}\rceil,
\]
and there are $O(Bm)$ nodes. Set $T_w=0$ at empty nodes and $T_w=F_i$ at
a singleton node for $x_i$. Otherwise define
\[
 T_w(v)=\sum_{u\text{ child of }w}\min_{z\in\mathcal D_k}
                   \bigl[T_u(z)+2|v-z|\bigr].
\]
\begin{lemma}\label{seed:lem:tree}
Let $d=\WED_w(X,V)$ and let $o$ be the root. For every guess $k$,
\[
 d\le T_o(0)+\Delta.
\]
If $k\ge d$, then
\[
 T_o(0)\le (2+2BD_0)d.
\]
Every node table $T_w$ is $2B$-Lipschitz on $\mathcal D_k$.
\end{lemma}
\begin{proof}
Each minimum transform is $2$-Lipschitz. An internal node sums at most
$B$ such transforms, a singleton table is $2$-Lipschitz, and an empty
table is constant. This proves the regularity assertion.

For the lower bound, expand the recurrence as a minimum over assignments
of shifts to tree nodes, with root shift zero, and choices of leaf
operations. Fix such an assignment. For leaf $i$, let $(u_i,v_i')$ be the
mass endpoints of its actual chosen target fragment, and let $z_i$ be its
shift. The leaf penalty is
\[
 e_i=|u_i-(M_X(i-1)+z_i)|+|v_i'-(M_X(i)+z_i)|.
\]
The connection expression from Lemma~\ref{seed:lem:connection} is at most
\[
 \sum_i e_i+|z_1|+\sum_{i<m}|z_{i+1}-z_i|+|z_m|+\Delta.
\]
Traverse the source tree from the root through its nonempty leaves in
left-to-right order and back to the root. Triangle inequality bounds the
shift variation in this display by twice the sum of absolute shift
differences on tree edges. Empty nodes and their leaves can be skipped.
The connection lemma converts the chosen operations into an alignment
of cost at most the recurrence's value plus $\Delta$.
Minimize over the assignment to obtain the bound on $d$.

For the upper bound, fix an optimal alignment.
Each node will receive its shift at the node's left source boundary.
At every source cut $i$
choose an alignment vertex $(i,b_i)$, in order, with $b_0=0$ and $b_m=n$.
For example, use the first vertex on each internal source-cut row. These
points partition the alignment into leaf subpaths of costs $d_i$ whose
sum is $d$. Assign shift
\[
 z_i=M_V(b_i)-M_X(i)
\]
to every tree node starting at cut $i$. Lemma~\ref{seed:lem:potential} gives
$|z_i|\le d\le k$.

If the subpath for $x_i$ deletes it, the empty candidate at boundary
$b_{i-1}$ has value $2g_i\le2d_i$ at shift $z_{i-1}$. Otherwise suppose
the subpath pairs $x_i$ with $y_j$. Every vertex of the full alignment
has mass discrepancy at most $d$, so
$|M_V(j-1)-M_X(i-1)|\le k$ and this singleton is allowed in
the leaf envelope. Put $a=M_V(j-1)-M_V(b_{i-1})\ge0$. Its value is
\[
 w(x_i,y_j)+a+|a+g(y_j)-g_i|
 \le 2w(x_i,y_j)+2a\le2d_i.
\]
Here $a$ is the cost of insertions before the diagonal, and the reverse
triangle inequality bounds $|g(y_j)-g_i|$.

For a parent and child, the absolute change of their assigned shifts is
at most the cost of the alignment segment between their chosen starting
vertices. That segment lies inside the parent's segment, whose cost
also bounds the shift change.
There are at most $B$ children per parent. At each depth, parent segments
partition the alignment, so the sum of edge differences at that depth is
at most $Bd$. Sum over the $D_0$ depths and add the leaf bound $2d$.
\end{proof}

\subsection{Fast approximate leaf arrays}

Use the shared target-prefix array and maximum tree. Subtracting the
queried fragment's initial prefix mass supplies all local coordinates
$M_V(j)$ without scanning the fragment.

\begin{lemma}\label{seed:lem:leaf}
Fix a source character $x_i$, an allowance $a>0$, a guess $k$, and a
sorted list $G\subseteq\mathcal D_k$. In
\[
 O\left(\left(1+\frac{k}{a}+|G|\right)\log N\right)
\]
operations, up to finite-arithmetic overhead, one can compute values
$\widehat F_i(v)$ for $v\in G$ satisfying
\[
 F_i(v)\le\widehat F_i(v)\le F_i(v)+a/2.
\]
The bound has no dependence on the target fragment's character count.
\end{lemma}
\begin{proof}
Set $t=a/4$. Retain only allowed singleton candidates with $g(y_j)>t$;
retain all empty-boundary candidates. Compare a discarded singleton with
the empty candidate at its starting boundary. The latter minus the former
is at most
\[
 g_i-w(x_i,y_j)+g(y_j)\le2g(y_j)\le2t=a/2.
\]
This proves the approximation bound.

Allowed retained starts lie in an interval of mass width $2k$. For any two
successive retained starts, their difference exceeds $t$, since the first
retained character itself has gap cost greater than $t$. There are
$O(1+k/t)$ such characters. Binary searches locate the allowed target
position interval, and the maximum tree reports its characters heavier
than $t$ in $O((1+k/t)\log N)$ time.

For each reported character, let
\[
 u=M_V(j-1)-M_X(i-1),\quad z=M_V(j)-M_X(i),\quad c_j=w(x_i,y_j).
\]
Its value as a function of $v$ is $c_j+|v-u|+|v-z|$. Let
$l_j=\min\{u,z\}$ and $r_j=\max\{u,z\}$. This function has the three
pieces
\[
 \begin{cases}
 c_j+l_j+r_j-2v,&v\le l_j,\\
 c_j+r_j-l_j,&l_j\le v\le r_j,\\
 c_j-l_j-r_j+2v,&v\ge r_j.
 \end{cases}
\]
Sort the breakpoints. Suffix and prefix minima evaluate the minimum left
and right pieces at every queried shift. A sweep with an active minimum
heap evaluates the plateau pieces. The work is linear in the number of
reported candidates plus queried shifts, with a logarithmic factor.

For an empty candidate, put $s=M_X(i-1)+v$. Its minimum value is
\[
 g_i+\min_j\bigl(|M_V(j)-s|+|M_V(j)-s-g_i|\bigr)
 =2g_i+2\operatorname{dist}([s,s+g_i],\{M_V(0),\dots,M_V(n)\}).
\]
A binary search tests whether the interval contains a target boundary.
Otherwise its predecessor and successor give the distance. This costs
$O(\log N)$ per queried shift, including when the predicted interval
extends beyond the target.
\end{proof}

The procedure \texttt{ADDITIVE-LEAF-ARRAY} implements the construction
in Lemma~\ref{seed:lem:leaf} and returns its array in the supplied grid
order. All substitution queries use the original occurrence handles.
Reporting and heap ties use original target positions.

\subsection{Adaptive integer grids and precision reconstruction}

Set the root allowance $a_o=k$. Independently for each child $u$ of a
nonempty internal node $w$, draw $R_u$ uniformly from $\{1,\dots,B\}$ and set
\[
 t_u=\frac{a_wR_u}{BQ_0},\qquad
 a_u=\frac{\lambda}{2}t_u
     =\frac{a_wR_u}{128BH^5}.
\]
Use independent draws for different entry indices, guesses, and tree
edges. All draws can be defined in advance, including unused values at
empty nodes. A draw uses exactly $\log_2B$ fair bits.

At a nonempty node $w$, let $h_w$ be the largest power of two at most
$a_w/(16B)$ if that power is at least one, and otherwise let $h_w=1$.
Equivalently,
\[
 h_w=\max\{1,2^{\lfloor\log_2(a_w/(16B))\rfloor}\}.
\]
The implementation uses bit lengths and integer comparisons, not a
real-arithmetic logarithm. Define the integer grid
\[
 G_w=(h_w\mathbb Z\cap[-k,k])\cup\{-k,0,k\}.
\]
An empty node is represented by the identically zero function, with no
allocated grid. At a nonempty leaf, evaluate Lemma~\ref{seed:lem:leaf} with
allowance $a_w$ on $G_w$, obtaining $\widetilde T_w$.

For an internal node and $v\in G_w$, form
\[
 \widehat s_u(v)=\min_{z\in G_u}
          \bigl(\widetilde T_u(z)+2|v-z|\bigr)
\]
for each nonempty child; for an empty child set $\widehat s_u(v)=0$.
Then define
\[
 \widetilde T_w(v)=\sum_{u\text{ child of }w}
   \1_{\widehat s_u(v)\ge t_u}
   \max\{\widehat s_u(v),a_w/Q_0\}.
\]
This algorithm computes all nonempty nodes. It does not prune by source
mass, source width, or a claimed upper bound on true node costs.

For a sorted child grid, compute prefix minima of
$\widetilde T_u(z)-2z$ and suffix minima of $\widetilde T_u(z)+2z$.
Merge the parent coordinates with these two sweeps. Their minimum gives
$\widehat s_u(v)$ at every parent coordinate in
$O(|G_w|+|G_u|)$ operations. The grids need not be nested.

\begin{lemma}\label{seed:lem:grid}
Suppose a nonempty child has, on its grid,
\[
 T_u(z)/f-a_u\le\widetilde T_u(z)\le fT_u(z)+a_u,
 \qquad 1\le f\le2.
\]
Write $s_u(v)=\min_{z\in\mathcal D_k}(T_u(z)+2|v-z|)$. Then
\[
 s_u(v)/f-\lambda t_u\le\widehat s_u(v)
                  \le f s_u(v)+\lambda t_u.
\]
\end{lemma}
\begin{proof}
The lower bound with error $a_u$ follows directly: restriction of the
minimum can only increase it, and a nonnegative penalty dominates its
value divided by $f$.

If $h_u=1$, the child grid contains the full integer domain, so the upper
bound with error $a_u$ also follows directly. This case is essential when
$a_u$ is smaller than one.

If $h_u>1$, round an integer minimizer to a child grid point at distance
at most $h_u\le a_u/(16B)$. By Lemma~\ref{seed:lem:tree}, its true
value increases by at most $2Bh_u$. Its penalty increases by at most
$2h_u$. Hence
\[
 \widehat s_u(v)\le f s_u(v)+a_u+(2Bf+2)h_u
                  \le f s_u(v)+\tfrac{11}{8}a_u
                  \le f s_u(v)+\lambda t_u.
\]
The penultimate bound uses $B\ge1$ and $f\le2$; the last uses
$a_u=\lambda t_u/2$. The lower bound is weaker than the already established
error-$a_u$ bound. Empty children satisfy both inequalities exactly.
\end{proof}

\begin{lemma}\label{seed:lem:concentration}
For every fixed $c_0>0$, when $N$ is sufficiently large as a function
of $C,s,c_0$, with probability at least $1-N^{-c_0}$ the following holds
simultaneously over $\mathcal E$, all guesses, all nodes, and their grid
points:
\[
 T_w(v)/f_w-a_w\le\widetilde T_w(v)\le f_wT_w(v)+a_w,
 \qquad
 f_w=\left(\frac{1+\lambda}{1-\lambda}\right)^{D_0-\operatorname{depth}(w)}
 \le2.
\]
\end{lemma}
\begin{proof}
Fix an internal node and condition only on the root-to-node allowance
draws, including its incoming edge.
Then $a=a_w$, the node's grid, and each true child transform
$s_i=s_i(v)$ at a fixed grid point are determined.
The comparison lists below depend only on these fixed data.
They are independent of the fresh child draws $R_i$.
We apply concentration to them without conditioning on descendant
success. Separately, the child invariants will place the computed
transforms between these lists.
Let $f'$ be the factor for the children's structural depth.
Write $\bar a=a/Q_0$ for this node's allowance scale, and put
\[
 t_i=\bar aR_i/B,\quad
 h_i^- =\frac{s_i}{f'(1+\lambda)},\quad
 h_i^+=\frac{f's_i}{1-\lambda},\quad
 \phi_i(h)=\1_{h\ge t_i}\max\{h,\bar a\}.
\]
If the child invariants hold, Lemma~\ref{seed:lem:grid} gives
$s_i/f'-\lambda t_i\le\widehat s_i\le f's_i+\lambda t_i$.
Consequently,
\[
 \phi_i(h_i^-)\le\phi_i(\widehat s_i)\le\phi_i(h_i^+).
\]
Indeed, $h_i^-\ge t_i$ implies $\widehat s_i\ge h_i^-$; and
$\widehat s_i\ge t_i$ implies $\widehat s_i\le h_i^+$.

Concentrate each of the two fixed lists $h_i=h_i^-$ and $h_i=h_i^+$
without conditioning on descendant success. When $0\le h_i<\bar a$,
$\phi_i(h_i)$ equals $\bar a$ with probability
$\lfloor Bh_i/\bar a\rfloor/B$ and zero otherwise. When $h_i\ge \bar a$, it equals
$h_i$ deterministically. Thus
\[
 0\le h_i-\E\phi_i(h_i)\le \bar a/B,\qquad
 \operatorname{Var}(\phi_i(h_i))\le \bar a h_i,\qquad
 |\phi_i(h_i)-\E\phi_i(h_i)|\le \bar a.
\]
Write $V_h=\sum_i h_i$ and $Z=\sum_i\phi_i(h_i)$. The total bias is at most
$\bar a$. Bernstein's inequality, applied with deviation
$\lambda V_h+a/2$, gives
\[
 \Prb\bigl(|Z-V_h|>\lambda V_h+a\mid\text{ancestors}\bigr)
 \le 2\exp(-\lambda Q_0/12).
\]
If $V_h\ge a$, the variance denominator is at most $3\bar a V_h$ and the
squared deviation is at least $\lambda aV_h$.
If $V_h\le a$, the denominator is at most $3\bar a a$ and the squared
deviation is at least $a^2/4$. The bias is absorbed by the remaining $a/2$ because $Q_0\ge2$.

On these two concentration events, the displayed sandwich inequality gives the factor
update $(1+\lambda)/(1-\lambda)$ and additive error $a$. Induct from the
leaves, where Lemma~\ref{seed:lem:leaf} gives an even smaller error. Early
singleton leaves can use the displayed larger structural factor. Empty
nodes are exact. Since $D_0\le H$ for sufficiently large $N$, the product
factor is less than two.

It remains to justify the union bound for random grids. A grid is fixed
after conditioning on the node's ancestors. Its cardinality is at most
$2k+1$, since it contains only integer shifts. The largest guess is at
most $2(M(X)+M(V))\le2N^{C+1}$. There are at most $N^s$ entries,
$O_C(H)$ guesses per entry, and $O(BN)$ nodes per guess. Thus the number
of tested lists is bounded by $N^{s+C+2+o(1)}$. We may union-bound the
uniform conditional bound at each fixed node over its polynomial maximum
grid count, then remove the conditioning. No enumeration of all ancestor
histories is required. Finally,
$\lambda Q_0/12=H^3/768$ dominates every fixed multiple of $\log N$.
This proves the claimed simultaneous failure probability for sufficiently
large $N$.

An explicit smaller event count is also available. A nonempty internal
source node has at least two nonempty children, and the nonempty leaves
are the source occurrences. Thus a source of width $m$ has at most $m-1$
internal nodes. Empty nodes and singleton leaves need no concentration
event. There are at most $(C+1)H+2$ guesses per entry and at most
$4N^{C+1}+1$ integer shifts per grid. If $\mathcal E_{\rm rec}$ is the
intersection of both fixed-list events at every tested parent coordinate,
then
\[
 \Prb(\mathcal E_{\rm rec}^{\,c})
 \le4N^s\bigl((C+1)H+2\bigr)(N-1)(4N^{C+1}+1)
       \exp(-H^3/768).
\]
This counts the random grids only after the stated conditioning and then
integrates the uniform conditional bound. For $C=3,s=6$ it is at most
$100HN^{11}\exp(-H^3/768)$.
\end{proof}

\subsection{Standalone entry evaluation}

The following procedures evaluate exactly the arrays just defined.
Use the canonical balanced cuts, left-to-right child order, and increasing
dyadic guesses. Index each allowance draw by the entry identifier, its
guess, and its tree-edge path. The identifier includes the table-instance
identifier and the complete family index. When a family is indexed by
source-tree nodes and target endpoints, retain the full source-node label
and original endpoint indices. Equal fragments with different labels
retain different keys.

For a nonempty unequal pair, define the candidate at guess $k$ by
\[
 U_k=\left\lceil2\bigl(\widetilde T_o(0)+k\bigr)+\Delta\right\rceil,
 \qquad \Delta=|M(X)-M(V)|.
\]
The entry answer is the minimum of all these candidates and
$M(X)+M(V)$. Empty sources and equal strings are handled exactly before
sampling.

\begin{wedalgorithm}{Adaptive coarse-node arrays}{alg:coarse-node}
\Require A private coarse-tree node $w$, allowance $a_w>0$, and fixed guess $k$.
\Ensure The sorted grid array $(G_w,\widetilde T_w)$, or the zero function for an empty node.
\Procedure{COARSE-NODE}{$w,a_w,k$}
  \If{$w$ is empty} \State \Return the identically zero function \EndIf
  \State $h_w\gets\max\{1,2^{\lfloor\log_2(a_w/(16B))\rfloor}\}$, using bounded binary arithmetic.
  \State $G_w\gets$ the sorted list $(h_w\mathbb Z\cap[-k,k])\cup\{-k,0,k\}$.
  \If{$w$ is a singleton}
    \State $\widetilde T_w\gets\Call{ADDITIVE-LEAF-ARRAY}{w,k,a_w,G_w}$.
    \State \Return $(G_w,\widetilde T_w)$.
  \EndIf
  \For{each child $u$ in left-to-right order}
    \State Draw $R_u$ uniformly from $\{1,\ldots,B\}$ using the entry's private stream.
    \State $t_u\gets a_wR_u/(BQ_0)$; $a_u\gets a_wR_u/(128BH^5)$.
    \State Obtain the child array from \Call{COARSE-NODE}{$u,a_u,k$}.
  \EndFor
  \For{each child $u$}
    \If{$u$ is empty}
      \State Set $\widehat s_u(v)\gets0$ for every $v\in G_w$.
    \Else
      \State Use two merged sweeps to compute, for every $v\in G_w$,
      \Statex \hspace{\algorithmicindent}$\widehat s_u(v)\gets\min_{z\in G_u}\{\widetilde T_u(z)+2|v-z|\}$.
    \EndIf
  \EndFor
  \For{each $v\in G_w$}
    \State $\widetilde T_w(v)\gets\sum_{u\text{ child of }w}\1_{\widehat s_u(v)\ge t_u}\max\{\widehat s_u(v),a_w/Q_0\}$.
  \EndFor
  \State \Return $(G_w,\widetilde T_w)$.
\EndProcedure
\end{wedalgorithm}

\begin{wedalgorithm}{A standalone initial-seed entry}{alg:initial-entry}
\Require Fragments $X,V$, completed deterministic preprocessing, and the entry's private random stream.
\Ensure A nonnegative integer at most $M(X)+M(V)$; no other initial entry is requested.
\Procedure{INITIAL-ENTRY}{$X,V,\omega$}
  \If{$X$ is empty} \State \Return $M(V)$ \EndIf
  \If{$|X|=|V|$ and corresponding symbols have oracle distance zero}
    \State \Return $0$
  \EndIf
  \If{$\chi<2^{1000}$}
    \State \Return the exact alignment dynamic-programming value.
  \EndIf
  \State Form the local source prefix masses.
  \State $\Delta\gets|M(X)-M(V)|$; $U_{\rm gap}\gets M(X)+M(V)$.
  \State $K\gets\operatorname{leastpow}_2(U_{\rm gap})$; $U_{\rm best}\gets U_{\rm gap}$.
  \For{$k=1,2,4,\ldots,K$}
    \State Build the private balanced $B$-ary source tree, with root $o$.
    \State $(G_o,\widetilde T_o)\gets\Call{COARSE-NODE}{o,k,k}$.
    \State $U_k\gets\lceil2(\widetilde T_o(0)+k)+\Delta\rceil$.
    \State $U_{\rm best}\gets\min\{U_{\rm best},U_k\}$.
  \EndFor
  \State \Return $U_{\rm best}$.
\EndProcedure
\end{wedalgorithm}

Empty child transforms are zero and need no allocated grid. Draws for
unused empty nodes may be included in the private stream; they are
independent of all used draws. No grid point outside
$[-k,k]\cap\mathbb Z$ is used. The endpoint $K$ of the guess loop is
fixed before $U_{\rm best}$ is updated.

The entry uses only its private workspace and the completed read-only
prefix arrays and reporting tree. It never calls another initial entry
or reads another entry's auxiliary state. On completion, the evaluator
caches its integer answer. Repeated requests use that cached answer.

\FloatBarrier

\subsection{Finite arithmetic and adaptive work}

\begin{lemma}\label{seed:lem:finite}
Every entry evaluation uses polynomial space, random bits, and finite
binary numbers of $O_C(H\polylog H)$ bits. Each arithmetic operation
can be implemented with $\poly(H)$ bit operations.
\end{lemma}
\begin{proof}
The exact small-input branch has polynomial workspace and finite binary
arithmetic by the alignment dynamic program. Consider the large-input
branch.

All masses and candidate breakpoints are integers of polynomial magnitude.
The guesses are at most $2N^{C+1}$, and every grid is a subset of the
integer interval $[-k,k]$. The complete tree and all its grids
have polynomial size.

All allowances and thresholds are dyadic. A common denominator for depth
at most $D_0$ needs
\[
 O\bigl(D_0(\log B+5\log H+7)+\log H\bigr)=O(H+\chi^2)=O(H)
\]
bits. Products of the $R_i$ have at most $D_0\log B=O(H+\log B)$ bits.
An empty-boundary candidate at zero bounds every true leaf value by
$2M_X(m)+2k$, which is polynomial in $N$. The computed values also have
polynomial magnitude: each reconstruction adds at most $a_w/Q_0$ per
child beyond a chosen transformed value, and a transform can always
choose a grid endpoint at penalty at most $4k$. Summing these polynomial
contributions over the polynomial-size tree gives a polynomial bound.
Minima select existing dyadic values, and sums preserve the common
denominator. A draw of $R_i$ uses exactly $\log_2B$ unbiased bits.
Ordinary binary arithmetic proves the assertion.
\end{proof}

\begin{lemma}\label{seed:lem:work}
For one entry $(X,V)$, all guesses can be evaluated in expected
$(1+m)2^{O_C(H/\chi+\chi^2)}=(1+m)N^{o(1)}$ work, uniformly over the
target fragment and without conditioning on any accuracy event.
\end{lemma}
\begin{proof}
From the allowance definition,
\[
 \E\left[\frac{a_w}{a_u}\,\middle|\,a_w\right]
 =128H^5\sum_{j=1}^B\frac1j\le H^9
\]
for all sufficiently large $H$. At depth $r$ this implies
\[
 \E[k/a_w]\le H^{9r}.
\]
The dyadic grid convention gives
\[
 |G_w|=O(1+Bk/a_w).
\]
When $h_w>1$, its mesh is greater than $a_w/(32B)$; when $h_w=1$,
$a_w<32B$, so the same cardinality bound follows directly.

Internal work is $O(B|G_w|+\sum_u|G_u|)$, apart from arithmetic factors.
Lemma~\ref{seed:lem:leaf} bounds leaf work by
$O((1+k/a_w+|G_w|)\log N)$. Summing over $O(Bm)$ nodes gives the
conservative expected bound
\[
 (1+m)B^{O(1)}H^{9D_0+O(1)}
\]
per guess, including logarithmic reporting factors.
Lemma~\ref{seed:lem:finite} absorbs finite arithmetic and sampling into
a factor $H^{O_C(1)}$. There are $O_C(H)$ guesses.
Moreover,
\[
 D_0\le1+\frac{H}{p_*\log_2 H},\qquad
 \log_2\bigl(B^{O(1)}H^{9D_0+O(1)}\bigr)
 =O_C(H/p_*+p_*\log H)
 =O_C(H/\chi+\chi^2)=o(H).
\]
Reading source gap costs, checking equality, and creating its local
prefix array costs $O(1+m)$. Target access uses the shared structures.
All estimates depend only on the independent allowance draws for their
work bounds, so the calculation does not require successful reconstruction.
The exact branch lies in a fixed finite input-size range and is absorbed
into the multiplicative constant.
\end{proof}

Before an entry's first request, the exposed computation has not
inspected its private random stream. The request can depend on previous
entry answers, but those answers depend only on their other streams.
Conditional on the complete preceding history and the first requested
entry, its stream retains its original distribution.
Lemma~\ref{seed:lem:work} gives the same expected bound under this
conditioning. Repeated requests need only the cached-answer lookup.

The event in Lemma~\ref{seed:lem:concentration} concerns the entire fixed
family $\mathcal E$, including entries that are never requested. This
accuracy event and the conditional work argument use separate
conditionings. No accuracy event is assumed in the work argument.

\subsection{The coarse-table guarantee}

\begin{theorem}\label{seed:thm:seed}
After $N\operatorname{poly}(\log N)$ charged binary preprocessing work
and $O(N)$ metric queries, the construction gives a lazily evaluated
random table $U$ on any predetermined indexed family $\mathcal E$ of at
most $N^s$ fragment pairs. Distinct indices may name the same pair and
have separate keys and private random streams.

For every fixed $c_0>0$, when $N$ is sufficiently large as a function of
$C,s,c_0$, with probability at least $1-N^{-c_0}$,
\[
 \WED_w(X,V)\le U(X,V)\le A_{\mathrm{init}}\WED_w(X,V)
\]
simultaneously on $\mathcal E$. Every zero-distance entry returns zero
on every outcome. Every returned value is a nonnegative integer at most
$M(X)+M(V)$.

Conditional on the history before any entry's first adaptive request,
its expected evaluation work, including finite binary computation and
random sampling, is
\[
 (1+|X|)2^{O_C(H/\chi+\chi^2)}=(1+|X|)N^{o(1)},
\]
uniformly over the entry and its target fragment. This bound includes
inaccurate executions. Each entry evaluation terminates on every
sequence of random bits, uses polynomial workspace, and performs finite
binary arithmetic with $\operatorname{poly}(\log N)$ work per operation.
It reads only the completed deterministic shared data and its private
state and does not request another entry.
\end{theorem}
\begin{proof}
If $X$ is empty, its exact answer is $M(V)$. Otherwise the equality test
examines corresponding symbols only when $|X|=|V|$ and costs $O(1+|X|)$.
It returns zero on every outcome when the strings are equal. For unequal
strings, put $d=\WED_w(X,V)\ge1$.

On the simultaneous event of Lemma~\ref{seed:lem:concentration}, the
root estimate satisfies $T_o(0)/2-k\le\widetilde T_o(0)\le2T_o(0)+k$.
Every candidate satisfies
\[
 U_k\ge T_o(0)+\Delta\ge d,
\]
by Lemma~\ref{seed:lem:tree}. The fixed guess set includes a power of two
with $d\le k<2d$, since $d\le M(X)+M(V)$. For this guess,
\begin{align*}
 U_k
 &\le4T_o(0)+4k+\Delta+1\\
 &\le(18+8BD_0)d\\
 &\le BH^4d=A_{\mathrm{init}}d
\end{align*}
for sufficiently large $H$. Here Lemma~\ref{seed:lem:tree} bounds
$T_o(0)$, Lemma~\ref{seed:lem:potential} gives $\Delta\le d$, and
$1\le d$. The trivial gap candidate $M(X)+M(V)$ is always at least $d$,
so taking the minimum preserves the event's lower bound and gives the
claimed upper bound. The simultaneous probability follows from
Lemma~\ref{seed:lem:concentration}.

On every outcome, the arrays are nonnegative and the answer is clamped
by the trivial gap candidate. The output is a nonnegative integer at
most $M(X)+M(V)$. The shared preprocessing has the bound
stated at the start of the section. Lemma~\ref{seed:lem:finite} gives
finite execution and polynomial workspace. Lemma~\ref{seed:lem:work}
and the private-stream argument give the conditional first-request work
bound without conditioning on accuracy. Caching completed answers
realizes the same random table for every adaptive request order.
\end{proof}


\section{The weighted refinement algorithm}\label{paper:sec:refinement}

One refinement pass takes a complete integer table $U$ with values in
$[0,R]$. Its accuracy analysis assumes that $U$ is good.
The algorithm never tests this assumption and applies every local guard.
The mathematical tables have complete values before an adaptive
evaluation schedule is chosen.

\subsection{Global parameters and the physical source tree}

Use the coarse-estimator parameter definitions with $X_0=x$ and $Y_0=y$,
so $R_0=R$. Reuse the
parameters $H,\chi,p_*,B,A_{\rm init}$ from the coarse estimator:
$H$ is the least power of two at least
$\lceil\log_2(N+R+2)\rceil$, $\chi=\log_2H$, $p_*=\chi$,
$B=H^{p_*}$, and $A_{\rm init}=H^{p_*+4}$. No fragment request or
refinement pass recomputes these parameters. We describe the branch
$\chi\ge2^{1000}$.

The parameter forms below follow the unit-cost refinement~\cite{openai}.
We use $p_*=\chi$ and $\ell=\Theta(\chi^3)$ to obtain the explicit
overhead bound.

\begin{definition}\label{refine:def:parameters}
Let $\ell$ be the least power of two at least $\chi^3$.
The source branching factor, factor reduction, warmup length, and
refinement length are
\begin{align*}
 M&=2^{\lfloor\chi/20\rfloor},& F&=\ell,&
 T&=2^{\lceil3\chi/4\rceil},& S&=H/\ell^2.
\end{align*}
For endpoint nets and whole-cell grouping, use
\[
 P=H^{p_*+20},\qquad \theta=\ell^{-40},\qquad
 L_0=32F,\qquad L=128F.
\]
Here $P$ sets the net resolution, $\theta$ sets the rounding and cell
widths, and $L_0,L$ are the cone slopes.
The clipping slack, velocity margin, entropy parameters, sampling
tolerance, and sample budget are
\[
 \delta=\ell^{-10},\qquad \tau=\kappa=\ell^{-4},\qquad
 \eta=\ell^5,\qquad \gamma=\ell^{-20},\qquad Q_1=\ell^{80}.
\]
Finally, set $\varepsilon_{\rm leaf}=\ell^{-3}$.

Build a complete $M$-ary source tree of physical height
\[
 J=\left\lceil\log_M\max(1,|x|)\right\rceil.
\]
A node $I=(l,u]$ at positive physical height has cuts
$l_i=l+\lfloor i(u-l)/M\rfloor$ and physical children
$I_i=(l_{i-1},l_i]$. Node labels distinguish repeated numeric intervals.
Only nodes at physical height zero use the leaf rule. An internal node
of width zero or one still uses the internal-node rule.

A state is a physical node $I$ and an original target-boundary pair
$q=(q^-,q^+)$ with $0\le q^-\le q^+\le|y|$. Its true cost is
$c(I,q)=\WED_w(x[I],y[q])$. Prefix masses, $z$, the endpoint distance
$\dd$, and connection charges use this fixed metric.
\end{definition}

\begin{lemma}\label{refine:lem:parameters}
The parameters satisfy
\[
 H=2^\chi,\qquad \chi^3\le\ell<2\chi^3,\qquad
 \ell^{200}\le H^{1/100},\qquad
 p_*+21\le\sqrt\ell,\qquad p_*\chi\le H/16.
\]
Moreover,
\[
 J\le100H/\chi,\qquad M^J\le M\max(1,|x|),\qquad
 \theta P\ge M,\qquad 4MJ/P\le\tfrac12,
 \qquad64F\theta<1.
\]
Every physical leaf has source width zero or one.
The number $S$ is a positive integer, and $T=o(S)$.
For every power of two $1\le A\le A_{\rm init}$ and
$a=\max(1,A/F)$, we have $2a/P<1$.
There are fewer than $e^{10H}$ states.

The clipping, entropy, and rounded-score error budgets satisfy
\[
 0<\kappa<\tfrac12,\qquad
 2\delta+8L\theta\le3\delta\le\kappa/\eta,
\]
\[
 64F(6L+2)+2L\le512LF,\qquad
 2\gamma+512LF\theta\le\kappa\delta.
\]
The label, group-failure, and action-coordinate bounds are
\[
 \log_2(64FM)+2\le2\ell,\qquad
 4\ell(P+1)^2e^{-\ell^{37}/3000}\le H^{-4},
\]
\[
 \ln\!\bigl(M(321P)^2\bigr)\le4(p_*+21)\chi.
\]
\end{lemma}
\begin{proof}
The definitions give the first two bounds. The third follows from
$200+600\log_2\chi\le\chi/100$. This inequality holds at the cutoff,
and $\chi/100-600\log_2\chi$ increases beyond it. The remaining
polynomial-versus-exponential comparisons follow from $H=2^\chi$.
In particular, $\chi-2\log_2\ell$ is a nonnegative integer, so $S$ is
a positive integer. Also $T/S\le2\ell^2H^{-1/4}\to0$.

For the depth, $\log_2M=\lfloor\chi/20\rfloor\ge\chi/40$ and
$\log_2N\le H$. The ceiling in $J$ gives the displayed depth and leaf
count. At physical depth $j$, each interval has width at most
$\lceil|x|/M^j\rceil$: each split gives child width at most the ceiling
of its parent width divided by $M$. Since $M^J\ge|x|$, leaf widths
are at most one. The other comparisons follow from
\[
 2a/P\le2A_{\rm init}/P=2H^{-16},\qquad
 \frac{\theta P}{M}\ge
 \frac{H^{\chi+20-1/20}}{\ell^{40}},\qquad
 \frac{MJ}{P}\le\frac{100H^{1+1/20}}{\chi P}.
\]
The tree has fewer than $2M^J$ nodes and fewer than $(N+1)^2$ target
pairs per node. The preceding bounds imply the state-count claim.

For the error budgets, substitution reduces the two nontrivial checks to
\[
 1024\ell^{-39}\le\ell^{-10},\qquad
 2\ell^{-20}+65536\ell^{-38}\le\ell^{-14}.
\]
Also $3\delta\le\kappa/\eta$ follows from $3\le\ell$.
The label and coordinate logarithms follow from
$\log_2M\le\chi/20$, $\log_2P=(p_*+20)\chi$, and
$\ell\ge\chi^3$.
For the group-failure inequality, take logarithms. Its required bound is
\[
 \ln(4\ell)+2\ln(P+1)+4\ln H\le\ell^{37}/3000.
\]
For $\chi\ge2^{1000}$, the left side is at most $4\chi^2$, while the
right side is at least $\chi^{111}/3000$.
\end{proof}

\subsection{The input table and leaf rule}

\begin{definition}\label{refine:def:input}
A refinement invocation receives a complete integer table
$U(I,q)\in[0,R]$ and a power of two $1\le A\le A_{\rm init}$.
Put
\[
 a=\max(1,A/F),\qquad
 \widetilde U(I,q)=\max\{\lceil U(I,q)/a\rceil,|z(I,q)|\}.
\]
The invocation has a \emph{good input} if
$c(I,q)\le U(I,q)\le Ac(I,q)$ at every state.

At a physical leaf use the generic leaf estimator of
Lemma~\ref{refine:lem:leaves} with $k=2\ell^3$, and denote it by $Q(I,q)$.
Thus $Q$ is integer-valued in $[0,R]$, is zero exactly at zero-cost states,
and satisfies
\[
 c\le Q\le(1+\varepsilon_{\rm leaf})c,
 \qquad |Q(I,q)-Q(I,q')|\le3\dd(q,q').
\]
A supplied input table need not equal $Q$ at a physical leaf.
\end{definition}

For a good input, integrality of $Fc$ gives
$\max(c/a,|z|)\le\widetilde U\le Fc$.

\begin{definition}\label{refine:def:rounding}
Let $\mathcal D_R=\{2^j:j\in\mathbb Z_{\ge0},\ 2^j\le R\}$.
Every scale range in the construction is intersected with $\mathcal D_R$,
including on inaccurate inputs. For $b\in\mathcal D_R$, put
\[
 g_b=\max(1,b/(16P)),\qquad
 r_b(q)=(\operatorname{pred}_{g_b}(q^-),
         \operatorname{pred}_{g_b}(q^+)).
\]
A scale-$b$ net pair is an ordered pair of boundaries in
$\mathcal N_{g_b}$.
\end{definition}

\begin{lemma}\label{refine:lem:rounding}
For every state and every scale $b\in\mathcal D_R$,
$\dd(q,r_b(q))\le2b/P$.
\end{lemma}
\begin{proof}
If $g_b=1$, rounding is exact. Otherwise
Lemma~\ref{refine:lem:nets} bounds each endpoint error by
$2g_b=b/(8P)$. Their sum is at most the stated allowance.
\end{proof}

Fix a core-run namespace $\nu$, including the input and metric identities.
An invocation has a pass label $p\ge1$ and a copy label $c_{\rm copy}$.
Its local input $U$ is the supplied complete table with seed label $p-1$,
denoted by $U^{(p-1)}$. Its factor $A$ is fixed across the copies of this
invocation. These labels identify values; they do not prescribe a
procedure for constructing supplied seed tables. Within one invocation,
we suppress $\nu,p,c_{\rm copy}$.

\subsection{Initial cones and Bellman propagation}

\begin{definition}\label{refine:def:cones}
At a physical leaf set $w_0=Q$. At an internal node and a query $(I,q)$,
if $U(I,q)=0$, set $w_0(I,q)=0$. Otherwise enumerate scales
\[
 U(I,q)/(16A^2)\le b\le4AU(I,q)
\]
and scale-$b$ net pairs $r$ at mass sup-distance at most $U(I,q)$
from $q$. Retain exactly those satisfying
$b/2\le U(I,r)\le4Ab$. Define $w_0(I,q)$ as the minimum of their
cones
\[
 \widetilde U(I,r)+L_0\dd(q,r).
\]
If no center is eligible, set $w_0(I,q)=R$.
\end{definition}

\begin{lemma}\label{refine:lem:cones}
For a good input,
\[
 \max(c/a,|z|)\le w_0\le2Fc,\qquad\Lip(w_0)\le L_0.
\]
At an internal node, the local definition equals the global minimum of
these cones over all eligible scale-net centers, together with all
zero-seed centers. The local rule remains finite for every input table.
\end{lemma}
\begin{proof}
Lemma~\ref{refine:lem:geometry} makes every global cone dominate the two
lower bounds. A leaf has these bounds and Lipschitz constant at most
$3\le L_0$. For positive $c$, choose $b\le c<2b$ and $r=r_b(q)$.
Lemma~\ref{refine:lem:rounding} gives
$b/2\le c(I,r)\le4b$, so the center is eligible. Its cone is at most
$Fc+(F+L_0)2b/P\le2Fc$.

The cone's nonnegative base cost puts a minimizing center at endpoint distance at most
$2Fc/L_0=c/16$. Its true cost lies in $[15c/16,17c/16]$.
A zero center cannot minimize at a positive-cost query, since its
endpoint distance is at least $c$. Combining the cost interval with
eligibility and $c\le U(I,q)\le Ac$ gives the stated local scale range.
The local radius contains the minimizing center. At zero cost the
query returns zero, which is also the global minimum. Global cone
minima are $L_0$-Lipschitz.

For every input, the tested scale range bounds the ratio of the radius
to net spacing. Net separation gives $O((1+PA^2)^2)$ raw centers per
scale. The fixed scale set is finite, and an empty list has the explicit
default $R$.
\end{proof}

\begin{definition}\label{refine:def:local}
An eligible representative $(b,r)$ is a scale-$b$ net pair satisfying
\[
 b/(4a)\le w_0(I,r)\le8Fb.
\]
For a positive $w_0(I,q)$, its local scales satisfy
\[
 w_0(I,q)/(16F)\le b\le8a w_0(I,q).
\]
Partition the target endpoint-mass plane into half-open cells of side
$\theta b$. The cell containing $r$ has indices
\[
 e^\pm=\left\lfloor M_y(r^\pm)/(\theta b)\right\rfloor.
\]
This floor acts on bounded rational integer-core data. Let
$\mathcal R(I,q)$ contain all eligible representatives in every whole
cell meeting the mass sup-square of radius $w_0(I,q)$ around $q$, for
these local scales. It includes the rest of a meeting cell, even outside
the square. If $w_0(I,q)=0$, let $\mathcal R(I,q)$ be empty.

For a representative put $p_0=M_y(r^-)-M_x(l)$. The wide list for
physical child $i$ consists of the ordered scale-$b$ net states $s_i$
satisfying
\[
 |M_y(s_i^-)-M_x(l_{i-1})-p_0|\le10b,
 \qquad |M_y(s_i^+)-M_x(l_i)-p_0|\le10b.
\]
The wide action set $\cA_{b,r}$ is the Cartesian product of these lists.
Its child fragments may overlap or leave gaps. For a finite nonnegative child table $f$,
put
\[
 V_r(\alpha;f)=\conn_r(\alpha)+\sum_{i=1}^M f(I_i,s_i(\alpha)).
\]
At an internal node define
\[
 \cB_I(f)(q)=\min\left\{w_0(I,q),
 \min_{(b,r)\in\mathcal R(I,q)}
 \left[L\dd(q,r)+\min_{\alpha\in\cA_{b,r}}V_r(\alpha;f)\right]\right\}.
\]
An empty minimum is $+\infty$. At physical leaves, let $\cB_I(f)(q)=Q(I,q)$ and set every $w_j$, $m$,
and $d_0$ to $Q$. At internal nodes, set
\[
 w_j=\cB(w_{j-1})\quad(1\le j\le T),\qquad
 m=T^{-1}\sum_{j=0}^{T-1}w_j,\qquad d_0=w_T.
\]
\end{definition}

\begin{lemma}\label{refine:lem:localization}
Fix a good input table and an internal node $I$. For every parent state
$q$ and every eligible representative $(b,r)$,
\[
 L\dd(q,r)<w_0(I,q)\quad\Longrightarrow\quad
 (b,r)\in\mathcal R(I,q).
\]
\end{lemma}
\begin{proof}
Put $w_q=w_0(I,q)$ and $w_r=w_0(I,r)$. The hypothesis implies $w_q>0$.
By Lemma~\ref{refine:lem:cones} and $L_0/L=1/4$,
\[
 |w_r-w_q|\le L_0\dd(q,r)<w_q/4.
\]
Eligibility and $3w_q/4<w_r<5w_q/4$ give
\[
 b>\frac{3w_q}{32F}>\frac{w_q}{16F},\qquad
 b<5aw_q<8aw_q.
\]
These inequalities put $b$ in the local scale range. Each endpoint-mass
displacement is at most $\dd(q,r)<w_q/L\le w_q$, so $r$ lies in the
query's mass square. Its cell meets that square, so the whole-cell rule
includes the eligible representative $(b,r)$.
\end{proof}

\begin{lemma}\label{refine:lem:warmup}
The Bellman operator is monotone in its nonnegative child table.
For a good input, it is $L$-Lipschitz at the parent and preserves the
lower bounds $c/a$ and $|z|$ whenever the child table has them.
For such an input, the warmup tables satisfy
\[
 \max(c/a,|z|)\le d_0\le m\le w_0\le2Fc,
 \qquad \cB(d_0)\le d_0,
\]
and all $w_j$ and $m$ are $L$-Lipschitz.
\end{lemma}
\begin{proof}
Fix a nonnegative child table. A Bellman cone that improves the fallback
$w_0(I,q)$ has $L\dd(q,r)<w_0(I,q)$, because its connection charge and
child sum are nonnegative. Lemma~\ref{refine:lem:localization} puts
its center in the local representative list. A cone excluded
from that list cannot improve the fallback. Thus the local rule equals
the global minimum over all eligible representatives and wide actions,
together with $w_0$. This is a fixed family of $L$-Lipschitz functions,
so its minimum is $L$-Lipschitz.

Monotonicity follows directly from the minima. The connection
inequalities of Lemma~\ref{refine:lem:geometry} preserve the lower
bounds. Start with $\cB(w_0)\le w_0$ and apply monotonicity to obtain
$w_T\le\cdots\le w_0$ and $\cB(w_T)\le w_T$. Averaging gives the
stated bounds for $m$. Physical leaves retain $Q$ throughout.
\end{proof}

On arbitrary inputs, the local rule of Definition~\ref{refine:def:local}
is the definition. It always has the finite fallback $w_0(I,q)$.

\subsection{Complete bands and rounded canonical actions}

\begin{definition}\label{refine:def:bands}
For an internal node, scale $b$, and candidate pair $r$, reject unless
$r$ is an ordered scale-$b$ net pair and satisfies the eligibility test
in Definition~\ref{refine:def:local}. For a passing candidate, put $Z=\min_{\alpha\in\cA_{b,r}}V_r(\alpha;m)$.
Reject its band if $Z\le0$ or $Z=+\infty$.
Otherwise set $\widehat v=z(I,r)/Z$ and reject unless
$|\widehat v|\le1-\tau$.

For $1\le j<M$, define the prefix function
\[
\begin{split}
 \Phi_j(p)=\min_{s_1,\ldots,s_j}\Bigl\{&
 \sum_{i=1}^j m(I_i,s_i)
 +|M_y(s_1^-)-M_y(r^-)|\\
 &+\sum_{i<j}|M_y(s_{i+1}^-)-M_y(s_i^+)|
 +|M_y(s_j^+)-(M_x(l_j)+p)|\Bigr\},
\end{split}
\]
where the minimum uses the first $j$ wide child lists.
This is a finite minimum of $1$-Lipschitz functions of $p$.

The complete band $\cC_{b,r}$ consists of shared-endpoint actions
$s_i=(v_{i-1},v_i)$ with
$r^-=v_0\le\cdots\le v_M=r^+$ and each $v_j\in\mathcal N_{g_b}$.
Every cut, including the two outside cuts, must satisfy
$|M_y(v_j)-M_x(l_j)-p_0|\le6b$.
At each internal cut put $p=M_y(v_j)-M_x(l_j)$ and require
\[
 |p-p_0-\widehat v\Phi_j(p)|\le\tau\theta b/M.
\]
Reject an empty complete band. The algorithm evaluates $\Phi_j$ only
at the reported candidate shifts; it does not search a real interval.

For a nonempty band let
\[
 h_i=\operatorname{ceil}_2\!\left(
 \max\{b/M,\max_{\alpha\in\cC_{b,r}}w_0(I_i,s_i(\alpha))\}\right),
 \qquad h_\Sigma=\sum_i h_i.
\]
Here $\operatorname{ceil}_2(v)$ is the least power of two at least
$v>0$, allowing negative exponents. The maximum uses only states on
complete band actions. Reject unless $h_\Sigma\le64Fb$.

Fix an action order once. Take its first action in $\cC_{b,r}$ and
round both endpoints downward separately at child $i$ to
$\mathcal N_{D_{h_i}}$, where $D_h=\max(1,\theta h/16)$.
Call the resulting action $\beta(b,r)$. Reject unless
\[
 \beta(b,r)\in\cA_{b,r},\qquad
 |z(I_i,s_i(\beta(b,r)))|\le2h_i\quad(1\le i\le M).
\]
A representative passing every test is retained. Its band record
contains $\cC_{b,r}$, all $h_i$, $h_\Sigma$, and $\beta(b,r)$.
\end{definition}

The band constraints imply that two allowed shifts at a cut differ by
at most $2\theta b/M$. Subtract their inequalities and use
$|\widehat v|\le1-\tau$ and the Lipschitz bound for $\Phi_j$.
The physical cuts and original boundary indices remain part of every
action, including at repeated numeric source cuts.

Lemma~\ref{refine:lem:parameters} and $h_i\ge b/M$ give
$D_{h_i}\ge g_b$, so nesting preserves the parent-net membership.
Each rounded endpoint moves by at most $\theta h_i$ on every input.
For a good input, the label bound gives $|z|\le h_i$ before rounding.
The rounded imbalance is then at most $h_i+2\theta h_i\le2h_i$, and the
shift radius is at most $6b+64F\theta b<10b$. Thus the final guards
follow from the preceding ones on a good input.

Separate child rounding can make a formerly shared endpoint distinct.
We charge each rounded action its full connection cost.
A raw pair that cannot extend to a complete band action does not enter
any label maximum. Rejection and missing-candidate sentinels occur only
in auxiliary records, never in numerical tables.

\subsection{Entropy under monotone costs with fixed block masses}\label{analytic:sec:entropy}

For a probability vector define
$\HH(\rho)=-\sum_k\rho_k\ln\rho_k$, with $0\ln0=0$.
The following inequality applies to a complete sequence of child costs.

\begin{proposition}\label{analytic:prop:entropy}
Let $\QQ$ be a nonempty compact convex set of probability vectors in
$\mathbb R^d$. Suppose that $\QQ$ contains a vector positive in every
coordinate. Partition the coordinates into blocks and assume that every
vector in $\QQ$ has the same total mass in each block.

Let $S\ge1$ be an integer, let $\eta>0$, and let $0<\kappa<1/2$.
Let $f_0,\ldots,f_S\in[0,1]^d$ be coordinatewise nonincreasing. Put
\[
 \Delta_t=f_{t-1}-f_t,\qquad
 \widehat g_t=-f_{t-1},\qquad
 u_t=\widehat g_t+(1-2\kappa)\Delta_t.
\]
Suppose that the range of $\Delta_t$ within each block is at most
$\kappa/\eta$. Let $\rho_t$ maximize
\[
 \HH(\rho)+\eta\ip{\sum_{j<t}u_j+\widehat g_t}{\rho}
 \quad\text{over }\rho\in\QQ.
\]
Then, for every $\rho^*\in\QQ$ and every integer $0\le s\le S$,
\[
 \sum_{t=1}^s\ip{f_t+\kappa\Delta_t}{\rho_t}
 \le \sum_{t=1}^s\ip{f_t}{\rho^*}
       +\frac{\ln d}{\eta}+2\kappa.
\]
If each left-hand vector $\rho_t$ is replaced by a vector within
$\zeta\ge0$ in $\ell_1$ distance, the right side increases by at most
$s\zeta$.
\end{proposition}

\begin{proof}
Define the potential
\[
 \Psi(v)=\max_{\rho\in\QQ}\{\HH(\rho)+\eta\ip v\rho\}.
\]
Compactness gives a maximizer. Strict concavity of entropy gives its
uniqueness. To prove positivity, suppose a maximizing vector has a zero
coordinate. Move distance $a>0$ toward a positive feasible vector.
The entropy contribution from zero coordinates increases by a positive
multiple of $a\ln(1/a)$. Contributions from previously positive
coordinates and the linear objective change by $O(a)$.
For sufficiently small $a$, the objective increases, contradicting
optimality.

Let $p$ be the maximizer for $v$. First-order optimality along every
feasible segment gives
\[
 \ip{\nabla\HH(p)+\eta v}{\rho-p}\le0
 \qquad(\rho\in\QQ).
\]
For probability vectors,
\[
 \HH(\rho)-\HH(p)
 =\ip{\nabla\HH(p)}{\rho-p}-\KL(\rho\Vert p),
 \quad
 \KL(\rho\Vert p)=\sum_k\rho_k\ln(\rho_k/p_k).
\]
Consequently
\[
 \HH(\rho)+\eta\ip v\rho
 \le\Psi(v)-\KL(\rho\Vert p).
\]
This argument works on the constrained hull $\QQ$; it does not assume
that its optimizer is an unconstrained softmax.

Put $v_{t-1}=\sum_{j<t}u_j$. Evaluating the old objective at the new
maximizer $\rho_t$ gives
\[
 \Psi(v_{t-1}+\widehat g_t)-\Psi(v_{t-1})
 \le\eta\ip{\widehat g_t}{\rho_t}.
\]
This remains valid when the increment has negative coordinates.

Let $\mu_k$ equal the minimum of $\Delta_t$ in the block containing $k$,
and put $r=\Delta_t-\mu$. Then $\mu\ge0$ and
$0\le\eta r_k\le\kappa$. Fixed block masses imply that
$\ip\mu\rho$ is the same for every $\rho\in\QQ$.
Applying the displayed KL penalty bound at $v_{t-1}+\widehat g_t$, the potential
increment produced by adding $(1-2\kappa)\Delta_t$ is at most
\[
 \eta(1-2\kappa)\ip\mu{\rho_t}
 +\ln\!\left(\sum_k\rho_{t,k}e^{\eta(1-2\kappa)r_k}\right).
\]
Indeed, after removing the constant block contribution, enlarge the
KL-penalized optimization to the full probability simplex. If
$Z=\sum_k\rho_{t,k}e^{\eta(1-2\kappa)r_k}$ and
$q_k=\rho_{t,k}e^{\eta(1-2\kappa)r_k}/Z$, its objective is
$\ln Z-\KL(\rho\Vert q)\le\ln Z$.
Nonnegativity of KL follows, for example, from $\ln z\le z-1$,
with zero coordinates treated by continuity.

For $0\le x\le\kappa\le1/2$,
\[
 e^{(1-2\kappa)x}\le1+(1-\kappa)x.
\]
To verify this, differentiate
$\ln(1+(1-\kappa)x)-(1-2\kappa)x$. Its derivative is
\[
 \frac{\kappa-(1-2\kappa)(1-\kappa)x}
 {1+(1-\kappa)x}\ge0,
\]
and the expression is zero at $x=0$.

Apply the displayed scalar inequality with $x=\eta r_k$ and then use
$\ln(1+y)\le y$. The bound for the remaining increment is at most
\[
 \eta(1-2\kappa)\ip\mu{\rho_t}
 +\eta(1-\kappa)\ip r{\rho_t}
 \le\eta(1-\kappa)\ip{\Delta_t}{\rho_t},
\]
where the final inequality uses $\mu\ge0$. Combining this with
the displayed prediction bound yields
\[
 \Psi(v_t)-\Psi(v_{t-1})
 \le\eta\ip{\widehat g_t+(1-\kappa)\Delta_t}{\rho_t}
 =-\eta\ip{f_t+\kappa\Delta_t}{\rho_t}.
\]

Now $\Psi(0)\le\ln d$, and $\HH(\rho^*)\ge0$. Also,
\[
 u_t=-f_t-2\kappa\Delta_t,\qquad
 0\le\sum_{t=1}^s\Delta_t=f_0-f_s\le\boldsymbol1.
\]
Summing the potential inequalities and evaluating the last potential at
$\rho^*$ gives
\begin{align*}
 \eta\sum_{t=1}^s\ip{f_t+\kappa\Delta_t}{\rho_t}
 &\le\Psi(0)-\Psi(v_s)\\
 &\le\ln d+\eta\sum_{t=1}^s\ip{f_t}{\rho^*}
       +2\eta\kappa\ip{\sum_{t=1}^s\Delta_t}{\rho^*}\\
 &\le\ln d+\eta\sum_{t=1}^s\ip{f_t}{\rho^*}+2\eta\kappa.
\end{align*}
Division by $\eta$ proves the claimed inequality. The case $s=0$ is immediate.

Finally, $f_t+\kappa\Delta_t=\kappa f_{t-1}+(1-\kappa)f_t$
belongs to $[0,1]^d$. Its inner product changes by at most $\zeta$
under an $\ell_1$ perturbation of size $\zeta$. Sum over the prefix.
\end{proof}

\subsection{Finite rational entropy optimization and exact sampling}\label{analytic:sec:optimizer}

\begin{proposition}\label{analytic:prop:optimizer}
Let $\QQ$ be the convex hull of a nonempty finite list of rational
probability vectors $a_1,\ldots,a_L\in\mathbb Q^d$. Suppose that every
coordinate is positive in at least one listed vector. Given
$v\in\mathbb Q^d$ and an integer $h_{\rm opt}\ge4$, a finite rational algorithm
constructs a mixture of the listed vectors whose mean $\rho$ satisfies
\[
 \norm{\rho-\rho^*}_1\le h_{\rm opt}^{-3},
 \qquad
 \rho^*=\mathop{\rm argmax}_{\rho\in\QQ}
       \{\HH(\rho)+\ip v{\rho}\}.
\]
The maximizer is unique. The construction cost in bit operations and the
stored mixture size are polynomial in $G_0=h_{\rm opt}+d+2$, $L$, and the maximum
numerator or denominator bit length of an input rational. The returned
mixture has dyadic probabilities and can be sampled exactly using a
bounded number of fair random bits, with deterministic bit cost polynomial
in the same parameters and a fixed space bound.
\end{proposition}

\begin{proof}
Averaging the input list gives a strictly positive feasible vector.
The compactness, strict concavity, and positivity argument in the
proof of Proposition~\ref{analytic:prop:entropy} establishes existence,
uniqueness, and positivity of $\rho^*$.

Set
\[
 \sigma=G_0^{-50},\qquad \varepsilon_{\rm grad}=G_0^{-20},\qquad
 F(\rho)=\HH(\rho)+\ip v{\rho},
 \quad
 F_\sigma(\rho)=-\sum_k(\rho_k+\sigma)\ln(\rho_k+\sigma)+\ip v{\rho}.
\]
For $x\in[0,1]$, the increment
$-(x+\sigma)\ln(x+\sigma)+x\ln x$ decreases with $x$.
Its endpoint magnitudes are at most
$\sigma\ln(1/\sigma)$ and $(1+\sigma)\ln(1+\sigma)$.
Both are at most $4\sqrt{\sigma}$ for $0<\sigma\le1$.
Summing these coordinate bounds gives
\[
 \sup_{\rho\in\QQ}|F_\sigma(\rho)-F(\rho)|\le4d\sqrt{\sigma}.
\]
The Hessian of $F_\sigma$ is diagonal with entries
$-1/(\rho_k+\sigma)$, so its gradient is $1/\sigma$-Lipschitz
in Euclidean norm. The Euclidean diameter of $\QQ$ is at most $2$.

Start at any listed vertex $\rho^{(0)}$. At iteration $j\ge0$, compute each
gradient coordinate to absolute error at most $\varepsilon_{\rm grad}$. Enumerate the list
and choose a vertex $v_j$ maximizing its inner product with the
approximate gradient, resolving ties in a fixed order. Set
\[
 \xi_j=\frac2{j+2},\qquad
 \rho^{(j+1)}=(1-\xi_j)\rho^{(j)}+\xi_j v_j.
\]
The $\ell_1$ distance between two probability vectors is at most $2$.
Thus $v_j$ loses at most $2\varepsilon_{\rm grad}$ in the true linearized objective relative
to any feasible mixture.

Write $D_j=\max_{\QQ}F_\sigma-F_\sigma(\rho^{(j)})$.
Concavity and the Lipschitz gradient bound give
\[
 D_{j+1}\le(1-\xi_j)D_j+2\varepsilon_{\rm grad}\xi_j+\frac{2\xi_j^2}{\sigma}.
\]
In particular, $\xi_0=1$ gives $D_1\le2\varepsilon_{\rm grad}+2/\sigma$,
independently of the magnitude of the rational vector $v$.
Induction in the displayed recurrence yields
\[
 D_j\le2\varepsilon_{\rm grad}+\frac8{\sigma(j+1)}\qquad(j\ge1).
\]
For the induction step, the required scalar inequality is
$j/(j+1)+1/(j+2)\le1$.

Run $K_{\rm opt}=G_0^{80}$ steps to obtain the preliminary mixture mean $\rho^{(K_{\rm opt})}$.
Using the uniform smoothing bound twice bounds its true objective gap by
\[
 F(\rho^*)-F(\rho^{(K_{\rm opt})})
 \le2G_0^{-20}+8G_0^{-30}+8dG_0^{-25}
 \le12G_0^{-20}.
\]
The Hessian of $-\HH$ at a positive probability vector has entries
$1/\rho_k\ge1$. Combining this bound with constrained first-order
optimality at the positive maximizer gives
\[
 F(\rho^*)-F(\rho)\ge\tfrac12\norm{\rho-\rho^*}_2^2
 \qquad(\rho\in\QQ).
\]
For an endpoint with a zero coordinate this follows by continuity along
the segment from $\rho^*$. Hence
\[
 \norm{\rho^{(K_{\rm opt})}-\rho^*}_1
 \le\sqrt d\,\norm{\rho^{(K_{\rm opt})}-\rho^*}_2
 \le\sqrt{24d}\,G_0^{-10}\le \tfrac12 h_{\rm opt}^{-3}.
\]
For the last inequality use $d\le G_0$, $h_{\rm opt}\le G_0$, and $G_0\ge7$:
$\sqrt{24d}\,G_0^{-10}h_{\rm opt}^3\le\sqrt{24}G_0^{-13/2}<1/2$.
The remaining half of the error allowance will cover dyadic rounding.

We now implement every step with finite rational arithmetic.
After $K_{\rm opt}'\ge1$ iterations the weight of the vertex selected at
iteration $j$, where $0\le j<K_{\rm opt}'$, is
\[
 \frac2{j+2}\prod_{k=j+1}^{K_{\rm opt}'-1}\frac{k}{k+2}
 =\frac{2(j+1)}{K_{\rm opt}'(K_{\rm opt}'+1)}.
\]
The initial vertex has zero remaining weight because $\xi_0=1$.
The displayed weights sum to one. Store the sequence of selected
vertex indices, allowing repetitions.

Let $D$ be the product of all denominators in the input vertex list.
Its bit length is at most the sum of their bit lengths, hence polynomial
in the input parameters. All coordinates of $\rho^{(K_{\rm opt}')}$ have common
denominator $DK_{\rm opt}'(K_{\rm opt}'+1)$. They can be computed from the cumulative weighted
sum $\sum_{j<K_{\rm opt}'}(j+1)v_j$ and this common denominator.
This representation avoids multiplying denominators through successive
iterations. Its bit length is polynomial in the stated parameters.

It remains to approximate the gradient, whose $k$th coordinate is
$v_k-1-\ln(\rho_k+\sigma)$. For rational
$r=\rho_k+\sigma\in[\sigma,1+\sigma]$, set
\[
 b_0=\frac{r-1}{r+1}.
\]
Since $\sigma=G_0^{-50}<1/4$, we have $|b_0|\le1-\sigma$ and
$1-b_0^2\ge\sigma$. The identity
\[
 \ln r=2\sum_{j=0}^{\infty}\frac{b_0^{2j+1}}{2j+1}
\]
follows by integrating the geometric series for $2/(1-b_0^2)$ from
zero to $b_0$. Truncate after $K_{\log}=G_0^{60}$ terms. The absolute tail
is at most
\[
 \frac{2|b_0|^{2K_{\log}+1}}
 {(2K_{\log}+1)(1-b_0^2)}
 \le
 \frac{2e^{-2K_{\log}\sigma}}
 {(2K_{\log}+1)\sigma}
 <G_0^{-20}=\varepsilon_{\rm grad}.
\]
The final inequality follows from $K_{\log}\sigma=G_0^{10}$ and $G_0\ge7$.

Every summand is rational. The numerator and denominator lengths of
powers of $b_0$ grow at most linearly in $K_{\log}$ times the bit length
of $b_0$. Multiplying the odd denominators contributes
$O(K_{\log}\log K_{\log})$ bits.
Direct integer arithmetic computes the truncated sum with polynomial bit cost.
The rational linear
term $v_k$ and the constant $-1$ are added exactly.
Enumerating the input list and comparing rational scores now implements
the vertex choice. Repeating for $G_0^{80}$ iterations still has a
polynomial bit bound, with absolute constant exponents.
The approximate logarithms affect only the selected vertex; they are not
inserted into the exact mixture-coordinate representation.

Finally, let $\omega_j=2(j+1)/(K_{\rm opt}(K_{\rm opt}+1))$ for $0\le j<K_{\rm opt}$, and choose
the least integer $b_{\rm samp}$ with $2^{b_{\rm samp}}\ge4K_{\rm opt}h_{\rm opt}^3$. Compute
$q_j=\lfloor2^{b_{\rm samp}}\omega_j\rfloor$ by integer division. Add the integer
residual $2^{b_{\rm samp}}-\sum_jq_j$ to $q_0$. The probabilities
$\omega'_j=q_j/2^{b_{\rm samp}}$ are nonnegative and sum to one. If
$e_j=\omega_j-\lfloor2^{b_{\rm samp}}\omega_j\rfloor/2^{b_{\rm samp}}$, their total change is
\[
 \sum_j|\omega'_j-\omega_j|
 =2\sum_{j>0}e_j<2K_{\rm opt}2^{-b_{\rm samp}}\le\tfrac12h_{\rm opt}^{-3}.
\]
Each action vertex is a probability vector of $\ell_1$ norm one.
Consequently the new mixture mean $\rho=\sum_j\omega'_jv_j$ differs
from $\rho^{(K_{\rm opt})}$ in $\ell_1$ norm by at most $h_{\rm opt}^{-3}/2$. Together with
the preceding optimization bound this proves
$\|\rho-\rho^*\|_1\le h_{\rm opt}^{-3}$. The new mean stays in $\QQ$ and
preserves every block mass shared by the action vertices.

Store integer prefix sums of the $q_j$. Draw a uniform integer
$J_{\rm samp}\in\{0,\ldots,2^{b_{\rm samp}}-1\}$ using exactly $b_{\rm samp}$ independent fair bits.
Binary search for the first prefix sum strictly larger than $J_{\rm samp}$, and
return its vertex. This samples the returned dyadic mixture exactly;
zero weights and repeated vertices cause no difficulty. Here
$b_{\rm samp}=O(\log K_{\rm opt}+\log h_{\rm opt})$, and construction, storage, and sampling have
deterministic polynomial bit cost. There is no rejection loop. In
particular, this implementation terminates for every fair-bit sequence.
\end{proof}

\subsection{Group choices, online choices, and the clipped update}

Fix a pass, a copy, and its complete input $U$.
Give every potential group choice $(I,b,e,t)$ and every potential online
choice $(I,b,q,t)$ a separate private stream of independent fair bits.
These streams are mutually independent and independent of $U$.

\begin{definition}\label{refine:def:update}
For a node, scale $b$, whole cell $e$, and index $1\le t\le S$, let
$\mathcal R_e$ be all retained representatives in the cell.
If it is empty, the group has no candidate. Otherwise process dyadic
labels $h\in[b/M,64Fb]$ in increasing order, setting
\[
 k_h=\min\{M,\lceil Q_1Mh/b\rceil\}.
\]
If $k_h=M$, use the list $J_h=(1,\ldots,M)$. Otherwise draw $k_h$
independent uniform child labels, with replacement, from the group's
private stream. Use the same list for every representative in the whole
cell. Write $h_i(r)$ for the label in $r$'s band record and define
\[
 \sigma_t(r)=\conn_r(\beta(b,r))+
 \sum_h\frac{M}{k_h}
 \sum_{i\in J_h:\,h_i(r)=h}d_t(I_i,s_i(\beta(b,r))).
\]
The inner sum preserves sampled multiplicities. Choose the first
representative $r_*$ of minimum score in the fixed representative order,
and return its actual cost
$C_*=V_{r_*}(\beta(b,r_*);d_t)$, together with $r_*$.
The returned cost is not the sampled score. In the exact and sampled
sums, test $h_i(r)=h$ before requesting the corresponding scalar value.
A failing test creates no current-index child request.

For a retained band at $r$, introduce one coordinate $(i,s)$ for each
child state on a complete action. Map action $\alpha$ to the rational
probability vector $\rho(\alpha)$ with mass $h_i/h_\Sigma$ at
$(i,s_i(\alpha))$, and let $\cQ$ be their convex hull. Every vector has
fixed mass $h_i/h_\Sigma$ in child block $i$. Set
\[
 \widehat g_j(i,s)=-d_{j-1}(I_i,s)/h_i,\qquad
 \Delta_j(i,s)=(d_{j-1}(I_i,s)-d_j(I_i,s))/h_i,
 \qquad u_j=\widehat g_j+(1-2\kappa)\Delta_j.
\]
At index $t$, use Proposition~\ref{analytic:prop:optimizer} with
$h_{\rm opt}=H$ on the action list and the exact rational linear vector
\[
 v_t=\eta\left(\sum_{j=1}^{t-1}u_j+\widehat g_t\right).
\]
Its returned dyadic mixture has mean within $H^{-3}$ in $\ell_1$
distance of the unique maximizer of
\[
 \HH(\rho)+\eta\left\langle
 \sum_{j=1}^{t-1}u_j+\widehat g_t,\rho\right\rangle
 \quad\text{over }\rho\in\cQ.
\]
The online owner $(I,b,q,t)$ uses $r=r_b(q)$. If its band is rejected,
there is no candidate. Otherwise sample one whole action $\alpha_t$
exactly from this mixture using that owner's private stream. Return
its actual current child sum
$C=\sum_i d_t(I_i,s_i(\alpha_t))$ and retain the sampled action.
The action has shared endpoints and outside endpoints $r$, so its
connection charge is zero.

At a nonleaf, let $\Pi_t(I,q)$ be the minimum of
$\cB_I(d_{t-1})(q)$ and the selected actual group costs
$C_*+L\dd(q,r_*)$ for all distinct groups reached by
$\mathcal R(I,q)$. Let $Y_t(I,q)$ be the minimum of
$C+L\dd(q,r_b(q))$ over the online candidates at the local scales
of Definition~\ref{refine:def:local}. An empty online minimum is
$+\infty$. Define
\[
 d_t(I,q)=
 \begin{cases}
 0,&w_0(I,q)=0,\\
 \min\{\Pi_t(I,q),
       \max\{\Pi_t(I,q)/(1+\delta),Y_t(I,q)\}\},&w_0(I,q)>0.
 \end{cases}
\]
At physical leaves set $d_t=Q$ for every $t$.
\end{definition}

Every online coordinate appears in a complete band action. Averaging
the action list gives a positive feasible vector, as required by the
finite optimizer. Its coefficients use only child
indices $0,\ldots,t-1$. Numerical iterations reuse their finite arrays;
only the single sampled action requests current-index child values.
The Bellman candidate makes $\Pi_t$ finite, so the clipped update is a
finite rational even when no online candidate is available. Do not
clamp refinement values to $R$ within a pass.

\begin{definition}\label{refine:def:keys}
All source-node fields are physical labels. All target-index fields are
original indices in $[0,|y|]$. Scales belong to $\mathcal D_R$ and
indices $t$ satisfy $1\le t\le S$. Group keys use the cell indices from Definition~\ref{refine:def:local}.
Their cells intersect the target endpoint rectangle $[0,M_y(|y|)]^2$.
Leaf keys require physical height zero. Band, group, and online keys
require an internal physical node. Use fixed deterministic orders for
all finite enumerations.
The full keys are
\begin{description}
\item[Leaf.] $(\mathsf{leaf},\nu,I,q)$ specifies $Q(I,q)$.
It has no pass or copy label.
\item[Seed.] $(\mathsf{seed},\nu,p,I,q)$ specifies $U^{(p)}(I,q)$.
It names a supplied complete integer table in $[0,R]$.
When the supplied table is the coarse table, this key identifies its
corresponding standalone initial entry.
\item[Warmup.] $(\mathsf{warmup},\nu,p,c_{\rm copy},j,I,q)$
specifies $w_j$ for $0\le j\le T$.
\item[Mean.] $(\mathsf{mean},\nu,p,c_{\rm copy},I,q)$ specifies $m$.
\item[Refinement.] $(\mathsf{refinement},\nu,p,c_{\rm copy},t,I,q)$
specifies $d_t$ for $1\le t\le S$.
The index-zero key aliases the warmup key at index $T$.
\item[Band.] $(\mathsf{band},\nu,p,c_{\rm copy},I,b,r)$ specifies
rejection or a complete retained-band record.
\item[Group.] $(\mathsf{group},\nu,p,c_{\rm copy},I,b,e^-,e^+,t)$
specifies the label samples, winning representative, and actual cost
for the whole cell. It omits the requesting state.
\item[Online.] $(\mathsf{online},\nu,p,c_{\rm copy},I,b,q,t)$
specifies rejection or the sampled whole action and its actual child cost.
It includes the queried state.
\end{description}
The initial entries, group records, and online records are random owners.
Assign mutually independent fair-bit streams to every potential
owner, including owners that no evaluation reaches. Shared preprocessing for
the concrete coarse estimator is deterministic. Different core runs,
passes, and copies have disjoint owner namespaces.
\end{definition}

These streams define the mathematical random experiment. An implementation
can allocate bits when an owner is first reached and retain its decisions.
An unreached stream need not be stored. Initial entries use only their
own streams, identities, workspace, and completed read-only preprocessing.
They request no other initial entry. These properties are part of the
standalone coarse construction.

A scalar child key omits its caller, representative, group cell, sampled
position, and label $h$. These data describe requests and owners, not new
child values. Warmup, mean, and band keys retain the copy label even when
they are deterministic functions of the same input $U$. A displayed
$\wedkey{f}{I,q}$ below abbreviates the corresponding full key.



\begin{lemma}\label{refine:lem:group-cones}
Fix a good input table, an internal node $I$, and an index $t$.
Fix all complete child tables and the parent group streams, with
nonnegative child values. Let $\mathcal G_t$ contain every potential
whole-cell group at $I$ that returns a candidate, and write $(r_g,C_g)$
for that group's winning representative and actual cost.
Let $b_g$ be its scale.
For every parent state $q$,
\[
 \Pi_t(I,q)=\min\left\{
 \cB_I(d_{t-1})(q),
 \min_{g\in\mathcal G_t}\{C_g+L\dd(q,r_g)\}
 \right\}.
\]
Consequently, $\Pi_t(I,\cdot)$ is $L$-Lipschitz.
\end{lemma}
\begin{proof}
The group definition uses all retained representatives in its whole
cell and its own stream. Its canonical key omits the requesting state.
Once the child tables and streams are fixed, the payload $(r_g,C_g)$
is the same for every requester. Thus the displayed global cone family
does not depend on $q$. Each $r_g$ is eligible and $C_g\ge0$.

The global minimum is at most the local minimum that defines $\Pi_t$.
Conversely, consider a group cone below the Bellman fallback at $q$.
That fallback is at most $w_0(I,q)$, so
$L\dd(q,r_g)<w_0(I,q)$.
Lemma~\ref{refine:lem:localization} puts $(b_g,r_g)$ in
$\mathcal R(I,q)$. The local rule reaches its whole-cell group and
includes this same payload. Every omitted group cone
is at least the fallback. The two minima agree.

Lemma~\ref{refine:lem:warmup} makes the Bellman fallback $L$-Lipschitz.
Each selected group cone is also $L$-Lipschitz. Their global minimum
has the same Lipschitz bound, including at states outside the
comparison tree.
\end{proof}

\begin{lemma}\label{refine:lem:invariants}
For a good input, every outcome satisfies
\[
 \max(c/a,|z|)\le d_t\le d_{t-1}\le w_0,
\]
\[
 d_t(I,q)\le(1+\delta)d_t(I,q')+L\dd(q,q'),
 \qquad
 \cB(d_t)\le d_t\le\cB(d_{t-1})\quad\text{at nonleaves}.
\]
The spatial bound for $d_0$ holds without the factor $1+\delta$.
\end{lemma}
\begin{proof}
Induct first on physical height and then on $t$. Every extra candidate
is a current-child wide action with an endpoint cone. Earlier child
values dominate current ones. Thus both $\Pi_t$ and $Y_t$ dominate
$\cB(d_t)$, while $\Pi_t\le\cB(d_{t-1})$.
The clipped update of Definition~\ref{refine:def:update} gives the
Bellman inequalities. Monotonicity follows from
$\cB(d_0)\le d_0$ at $t=1$ and from the preceding lower Bellman
inequality thereafter. The connection inequalities of
Lemma~\ref{refine:lem:geometry} preserve the lower bounds.

Lemma~\ref{refine:lem:group-cones} makes $\Pi_t$ $L$-Lipschitz at every
parent state. The bounds
$\Pi_t/(1+\delta)\le d_t\le\Pi_t$ imply the displayed spatial
inequality. At leaves all assertions follow from the guarantee for $Q$
in Definition~\ref{refine:def:input}, since $3\le L$.
\end{proof}

\subsection{Functional table requests}

\begin{definition}\label{refine:def:requests}
Canonicalize a local $U$ key to its supplied seed-table key and a $d_0$
key to the warmup key at index $T$.
The functional operation $\operatorname{GET}(K)$ returns the complete
mathematical payload specified by the definitions above for the resulting
full key. For a finite list $\mathcal L$ of scalar keys,
$\operatorname{GATHER}(\mathcal L)$ canonicalizes, sorts, and
deduplicates them, and returns the local map from each distinct
canonical key to its $\operatorname{GET}$ value.
The scalar payload of an online record is its actual child-cost sum;
the sampled action is retained as that owner's decision.
The group payload is its winning representative and actual cost.
An auxiliary rejection or missing-candidate sentinel is not a numerical
value.
\end{definition}

These operations specify mathematical reads. Memoization can implement
them without changing values or owner decisions. After a gather, a
displayed table value in the code is a read from this local map; it does
not request the evaluator again. Repeated sampled occurrences retain
their multiplicity in arithmetic even when they share one canonical
scalar value. The call $\textsc{Leaf}(I,q)$ reads the shared $Q(I,q)$
value with key $(\mathsf{leaf},\nu,I,q)$.

\subsection{Exact procedural realization}

The routines below realize the complete definitions. Their inputs
include the fixed pass, copy, metric, and global parameters. Each state,
wide list, prefix minimum, and action retains the physical child label
and original target indices. The Bellman helper is a finite local
computation, not another mathematical table. Prefix minima, gathered
arrays, representative lists, scores, and optimizer iterates are scratch
data. The full scalar keys in Definition~\ref{refine:def:keys} fix every
recursive value independently of its requester.

\begin{wedalgorithm}{Warmup entries and local representatives}{alg:warmup-local}
\Require A fixed pass and copy, complete input $U$, physical node $I$, and state $q$; $0\le j\le T$ for \textsc{Warmup}, and internal $I$ for \textsc{Local}.
\Ensure The specified warmup value or the complete local representative list with whole-cell identifiers.
\Procedure{WARMUP}{$j,I,q$}
  \If{$I$ is a physical leaf} \State \Return \Call{LEAF}{$I,q$} \EndIf
  \If{$j>0$} \State \Return \Call{BELL}{$w_{j-1},I,q$} \EndIf
  \State $s\gets\Call{GET}{\wedkey{U}{I,q}}$.
  \If{$s=0$} \State \Return $0$ \EndIf
  \State Enumerate the initial scales and radius-$s$ centers of Definition~\ref{refine:def:cones}.
  \State Gather their $U$ keys, then retain the eligible centers.
  \If{no center is eligible} \State \Return $R$ \EndIf
  \State \Return the minimum initial cone, using the gathered values.
\EndProcedure
\Statex
\Procedure{LOCAL}{$I,q$}
  \State $w_{0,q}\gets\Call{GET}{\wedkey{w_0}{I,q}}$.
  \If{$w_{0,q}=0$} \State \Return the empty list \EndIf
  \State Initialize an empty representative list.
  \For{each local scale $b$ of Definition~\ref{refine:def:local}}
    \State Report all whole cells meeting the radius-$w_{0,q}$ mass square around $q$.
    \State Enumerate the ordered scale-$b$ net pairs $\mathcal R_b$ in these cells.
    \State Gather $\wedkey{w_0}{I,r}$ for every $r\in\mathcal R_b$.
    \State Retain exactly the eligible $(b,r)$, with their whole-cell identifiers.
  \EndFor
  \State \Return the representative list.
\EndProcedure
\end{wedalgorithm}

\begin{wedalgorithm}{Bellman propagation and the warmup mean}{alg:bell-mean}
\Require The fixed pass and copy, node $I$, and state $q$; for \textsc{Bell}, internal $I$ and a specified finite nonnegative child table $f$.
\Ensure The finite local Bellman value or the specified warmup mean.
\Procedure{BELL}{$f,I,q$}
  \State $v\gets\Call{GET}{\wedkey{w_0}{I,q}}$.
  \If{$v=0$} \State \Return $0$ \EndIf
  \State $\mathcal R\gets\Call{LOCAL}{I,q}$; generate its wide child-state lists.
  \State Gather the union of required keys $\wedkey{f}{I_i,s}$.
  \State Evaluate the local Bellman minimum using $v$ and the gathered array.
  \State \Return this minimum, with fallback $v$.
\EndProcedure
\Statex
\Procedure{MEAN}{$I,q$}
  \State Gather $\wedkey{w_j}{I,q}$ for $j=0,\ldots,T-1$.
  \State \Return their arithmetic mean.
\EndProcedure
\end{wedalgorithm}

\begin{wedalgorithm}{A complete retained band}{alg:band}
\Require Internal node $I$, scale $b$, candidate pair $r$, and the fixed pass and copy.
\Ensure Rejection or the complete band, child labels, and rounded canonical action.
\Procedure{BAND}{$I,b,r$}
  \If{$r$ is not an ordered scale-$b$ net pair} \State \Return rejection \EndIf
  \State $w_r\gets\Call{GET}{\wedkey{w_0}{I,r}}$.
  \If{$w_r<b/(4a)$ or $w_r>8Fb$} \State \Return rejection \EndIf
  \State Generate all wide child-state lists and gather their mean keys.
  \State Compute $Z$ from Definition~\ref{refine:def:bands}, using those means.
  \If{$Z\le0$ or $Z=+\infty$} \State \Return rejection \EndIf
  \State Compute $\widehat v$ and $p_0$ from that Definition.
  \If{$|\widehat v|>1-\tau$} \State \Return rejection \EndIf
  \State Report radius-$6b$ net candidates at every internal cut.
  \If{an outside endpoint fails the radius-$6b$ test} \State \Return rejection \EndIf
  \State Compute the finite prefix minima only at reported shifts, using the same means.
  \State Construct $\cC_{b,r}$ with all ordering, endpoint, radius, and band constraints of Definition~\ref{refine:def:bands}.
  \If{$\cC_{b,r}=\varnothing$} \State \Return rejection \EndIf
  \State Gather child $w_0$ only at states occurring on complete actions in $\cC_{b,r}$.
  \State Compute $(h_i)_i$ and $h_\Sigma$ from Definition~\ref{refine:def:bands}.
  \If{$h_\Sigma>64Fb$} \State \Return rejection \EndIf
  \State Take the first complete action in the fixed order and round both endpoints separately as specified there.
  \If{$\beta(b,r)\notin\cA_{b,r}$ or some $|z(I_i,s_i(\beta(b,r)))|>2h_i$}
    \State \Return rejection
  \EndIf
  \State \Return $(\cC_{b,r},(h_i)_i,h_\Sigma,\beta(b,r))$.
\EndProcedure
\end{wedalgorithm}

\begin{wedalgorithm}{One whole-cell group choice}{alg:group}
\Require Internal node $I$, scale $b$, whole cell $e$, index $1\le t\le S$, and this group's private stream.
\Ensure No candidate, or the first minimum-score representative and its actual current-index cost.
\Procedure{GROUP}{$I,b,e,t$}
  \State Enumerate all ordered scale-$b$ net pairs in the whole cell.
  \State Obtain each $\Call{GET}{\wedkey{\mathrm{band}}{I,b,r}}$; let $\mathcal R_e$ contain all passing representatives.
  \If{$\mathcal R_e=\varnothing$} \State \Return no candidate \EndIf
  \For{each dyadic label $h\in[b/M,64Fb]$, in increasing order}
    \State Compute $k_h$ from Definition~\ref{refine:def:update}.
    \If{$k_h<M$}
      \State Draw $J_h$ as $k_h$ independent uniform child labels with replacement from this group's stream.
    \Else
      \State $J_h\gets(1,\ldots,M)$.
    \EndIf
    \State Use the same $J_h$ for every representative in $\mathcal R_e$.
  \EndFor
  \State Initialize an empty scalar-key set $\mathcal K$.
  \For{each $r\in\mathcal R_e$, label $h$, and occurrence $i$ in $J_h$}
    \If{$h_i(r)=h$}
      \State Add $\wedkey{d_t}{I_i,s_i(\beta(b,r))}$ to $\mathcal K$.
    \EndIf
  \EndFor
  \State $\mathsf{val}\gets\Call{GATHER}{\mathcal K}$.
  \State Compute every score $\sigma_t(r)$ of Definition~\ref{refine:def:update} from $\mathsf{val}$, retaining sample multiplicities and the connection charge.
  \State Choose the first representative $r_*$ of minimum score.
  \State Gather all $M$ current child keys of $\beta(b,r_*)$.
  \State Retain the lists $J_h$ and the winning representative as the owner's decisions.
  \State \Return $r_*$ and its actual cost from Definition~\ref{refine:def:update}, using the gathered values.
\EndProcedure
\end{wedalgorithm}

\begin{wedalgorithm}{An online action from earlier child indices}{alg:online}
\Require Internal node $I$, local scale $b$, query $q$, index $1\le t\le S$, and this online owner's private stream.
\Ensure No candidate, or one sampled whole action's actual current child-cost sum.
\Procedure{ONLINE}{$I,b,q,t$}
  \State $r\gets r_b(q)$; $\mathsf{band}\gets\Call{GET}{\wedkey{\mathrm{band}}{I,b,r}}$.
  \If{$\mathsf{band}$ is rejected} \State \Return no candidate \EndIf
  \State List the distinct coordinates $(i,s)$ occurring on complete band actions.
  \State Gather $\wedkey{d_j}{I_i,s}$ at these coordinates for $j=0,\ldots,t-1$.
  \State Form the action vectors and $v_t$ from Definition~\ref{refine:def:update}.
  \State Run the finite optimizer of Proposition~\ref{analytic:prop:optimizer} with $h_{\rm opt}=H$, using only the gathered arrays.
  \State Draw one whole action $\alpha_t$ exactly from its returned dyadic mixture, using this owner's stream.
  \State Gather $\wedkey{d_t}{I_i,s_i(\alpha_t)}$ for $i=1,\ldots,M$.
  \State Retain $\alpha_t$ as the owner's decision.
  \State \Return the action's actual child-cost sum.
\EndProcedure
\end{wedalgorithm}

\begin{wedalgorithm}{The clipped refinement update}{alg:refine}
\Require Fixed pass and copy, node $I$, state $q$, and $1\le t\le S$; $d_0$ aliases $w_T$.
\Ensure The specified finite value $d_t(I,q)$, without an internal clamp to $R$.
\Procedure{REFINE}{$t,I,q$}
  \If{$I$ is a physical leaf} \State \Return \Call{LEAF}{$I,q$} \EndIf
  \State $w_{0,q}\gets\Call{GET}{\wedkey{w_0}{I,q}}$.
  \If{$w_{0,q}=0$} \State \Return $0$ \EndIf
  \State Compute the Bellman fallback with \Call{BELL}{$d_{t-1},I,q$}.
  \State $\mathcal R\gets\Call{LOCAL}{I,q}$.
  \For{each distinct whole-cell group $(b,e)$ reached by $\mathcal R$}
    \State Obtain $\Call{GET}{\wedkey{\mathrm{group}}{I,b,e,t}}$.
    \State Retain its representative and actual cost if it has a candidate.
  \EndFor
  \For{each local scale $b$ of Definition~\ref{refine:def:local}}
    \State Obtain $\Call{GET}{\wedkey{\mathrm{online}}{I,b,q,t}}$.
  \State Retain $(b,C)$ if it returns an actual cost $C$.
  \EndFor
  \State Compute $\Pi_t$ from the Bellman fallback and retained group cones, as in Definition~\ref{refine:def:update}.
  \State Compute $Y_t$ from retained online cones, taking $+\infty$ when none exists.
  \State \Return the clipped update of $\Pi_t$ and $Y_t$ in Definition~\ref{refine:def:update}.
\EndProcedure
\end{wedalgorithm}
\FloatBarrier



The initial cones use exactly the scale, radius, and eligibility rules in
Definition~\ref{refine:def:cones}. For local representatives, enumerate
the whole meeting cell before applying eligibility. Gather the $w_0$
values at all reported centers before that test. The whole-cell identity
makes a group independent of the state that first reaches it.

The band routine applies every rejection test in
Definition~\ref{refine:def:bands} on every input. It gathers label values
only for child states that lie on complete band actions. After separate
rounding, the selected group action keeps its full connection charge.

Each sampled group label uses exactly $\log_2M$ fair bits, since $M$
is a power of two. Repeated samples keep their multiplicity in the score
but share the same gathered scalar key. The label test precedes key
creation. The winner's actual current cost is evaluated after selection.
The online optimizer uses only gathered earlier-index values during its
iterations. The only current-index calls from this owner are the $M$
child values of its single sampled action.

The refinement update includes the Bellman fallback and has no internal
clamp to $R$. Repeated table reads use the same full key and owner
decisions.

\begin{lemma}\label{refine:lem:semantics}
Fix one supplied complete input table $U$, an invocation namespace, and
all its potential owner streams. The rules in this section determine a
unique complete scalar value and auxiliary payload at every key of that
invocation. Complete child histories depend only on $U$ and current-copy
owner streams at that child and its descendants. Conditional on $U$ and
these histories, current parent group and online streams retain their
joint independent fair-bit laws.
\end{lemma}
\begin{proof}
Order records lexicographically by within-pass stage, remaining
physical height, and phase. Supplied seed values and shared leaf values
are base values at stage zero. Warmup $w_j$ has stage $1+j$, a mean
has stage $T+1$, a band has stage $T+2$, and refinement, group, and
online records at index $t$ have stage $T+1+t$.
The alias $d_0=w_T$ has stage $T+1$.
Seed, warmup, and refinement records have phase $2$, group and online
records phase $1$, and mean, band, and leaf records phase $0$.

At a physical leaf, a warmup or refinement entry reads only its shared
leaf value. At an internal node, $w_0$ reads the supplied seed at its
query and at every reported initial center.
The local representative helper reads same-node $w_0$ values at the
query and every reported center before applying eligibility.
For $j>0$, a warmup entry uses this helper and reads earlier-index
warmup values at physical children.
A mean reads same-node warmup indices below $T$.
A band reads same-node $w_0$, child means, and child $w_0$ values on
complete band actions.

A scalar refinement entry reads same-node $w_0$, the local helper's
center values, and earlier-index child values for its Bellman fallback.
It also reads same-node group and online records with smaller phase.
These records read same-node bands at smaller stage, or at smaller
phase when $t=1$. An online record reads earlier child histories to
construct its distribution. A group or online record reads current-index
child values only for its tested score terms or its selected action.
These current-index reads decrease physical height, even when numerical
source intervals repeat. Thus every recursive read decreases the order.

All remaining operations are finite minima, rational calculations,
finite action enumeration, or the finite optimizer and bounded sampler
proved above. Induction on the finite order defines each
value uniquely and excludes a return to an unfinished caller, regardless
of input accuracy.

The same induction shows that a complete child table uses only its
input $U$ and owner streams in its physical subtree. It excludes all
current parent owners. Those streams are jointly independent of $U$
and of the descendant streams, by Definition~\ref{refine:def:keys}.
Conditioning on complete child histories, which are functions of these
other data, preserves their joint independent laws.
\end{proof}

A complete child history includes all mathematical states and all
indices at that child, whether evaluated or not. The realized reached-key
set or request trace can reveal a parent sample and cannot replace these
complete histories in the conditioning of Lemma~\ref{refine:lem:semantics}.

\FloatBarrier

\section{Accuracy of weighted refinement}\label{paper:sec:accuracy}

Fix a good input table $U$. The auxiliary comparison alignment follows true optimal costs. It lets us charge the refinement error across the source tree, including the bias of approximate leaf estimates.

\FloatBarrier

\subsection{Comparison tree and curvature}\label{refine:sec:comparison}

Start from a fixed state of true cost $D>0$.
The comparison tree stops at physical leaves and true zero-cost states.
At a zero-cost terminal, set $Q=0$ for this analysis.
At each nonterminal, choose $b\le c<2b$ and set $r=r_b(q)$ and
$e=\dd(q,r)$.
Fix an optimal weighted alignment for $(I,r)$ by a deterministic
tie rule. Choose ordered points on it at the child cuts, then round its
internal target anchors to the scale-$b$ net.
Let $q_i$ be the resulting child states and $\alpha^*$ their
shared-endpoint action.
This comparison tree depends on true costs and fixed rounding rules,
not on refinement randomness.

\begin{lemma}\label{refine:lem:comparison-costs}
The chosen $(b,r)$ is a representative in the local scale range, and $\alpha^*$ is a wide action with shift radius $6b$. Moreover,
\[
 c(I,r)\le\sum_i c(I_i,q_i)\le(1+4M/P)c,
 \qquad c\le Le+\sum_i c(I_i,q_i).
\]
Writing $C_{\rm term}$ for the sum of terminal true costs,
\[
 C_{\rm term}\le(1+8MJ/P)D\le2D,\quad
 \sum_{\rm nonterm}c\le2JD,\quad
 \sum_{\rm nonterm}e\le4JD/P.
\]
Terminal table values equal $Q$, and their sum $\widetilde C_{\rm term}$ lies between $C_{\rm term}$ and $(1+\varepsilon_{\rm leaf})C_{\rm term}$.
\end{lemma}
\begin{proof}
Endpoint regularity gives $c-e\le c(I,r)\le c+e$.
Lemma~\ref{refine:lem:cones} and $e\le2b/P$ imply representative
eligibility and the local scale bounds.
Along the optimal path, shift variation is at most $c(I,r)$ by
Lemma~\ref{refine:lem:geometry}.
Internal rounding adds at most $b/P$ per anchor, keeping the shift
radius below $6b$.

Before rounding, the subpath costs sum to $c(I,r)$.
Each internal anchor changes two child endpoints, adding at most
$2(M-1)b/P$ in total. This proves the upper child-cost bound.
The rounded children retain shared endpoints, so concatenation gives
$c(I,r)\le\sum_i c(I_i,q_i)$. Moving the parent's endpoints back to $q$
then gives $c\le Le+\sum_i c(I_i,q_i)$.

At each level, replace every nonterminal frontier state by its comparison
children and keep every terminal frontier state.
The frontier's total true cost grows by a factor of at most $1+4M/P$.
There are at most $J$ replacement levels.
Since $4MJ/P\le1/2$, the total multiplier is at most $1+8MJ/P\le2$.
Summing over levels and using $e\le2c/P$ proves the frontier bounds.
Lemma~\ref{refine:lem:leaves} gives the terminal approximation bound.
At a terminal zero-cost state, all table values are zero by the
good-input upper bound.
\end{proof}

At a nonterminal write
\[
 m_q=m(I,q),\quad m_i=m(I_i,q_i),\quad
 W=\sum_i m_i,\quad G=W-Z,\quad R_q=m_q-W,
\]
and $z_i=z(I_i,q_i)$, $z_r=z(I,r)$. Shared endpoints give $\sum_i z_i=z_r$. The preceding bounds yield
\[
 0<\max\{c(I,r)/a,|z_r|\}\le Z\le W\le4Fc<8Fb.
\]
Define
\[
 K=\sum_i\frac{z_i^2}{m_i}-\frac{z_r^2}{W}
   =\sum_{i:m_i>0}\frac{(z_i-vm_i)^2}{m_i},\qquad v=z_r/W,
\]
with $0/0=0$. The lower invariant makes this convention legitimate.

\begin{lemma}\label{refine:lem:curvature}
For $h_q=Le+w_0(I,q)/T$,
\begin{align*}
 \sum(K+G)&\le C_{\rm term}+3\widetilde C_{\rm term}
                   +3\sum e+4\sum h_q,\\
 \sum(K+G)&\le5000F(1+J/P+J/T)D.
\end{align*}
All sums are over nonterminals.
\end{lemma}
\begin{proof}
Averaging the warmup Bellman comparisons and minimizing their fixed action gives
\[
 m_q-(w_0-w_T)/T\le Le+Z.
\]
This gives $R_q+G\le h_q$, $R_q^+\le h_q$, and
$G\le R_q^-+h_q$.
The sum of $R_q=m_q-\sum_i m_i$ telescopes to
$m_{\rm root}-\widetilde C_{\rm term}$.
Combining these bounds gives
\[
 \sum R_q^-\le\sum h_q+\widetilde C_{\rm term},
 \qquad \sum G\le2\sum h_q+\widetilde C_{\rm term}.
\]

Let $E(I,q)=z(I,q)^2/m(I,q)$, with $E(I,q)=0$ when $m(I,q)=z(I,q)=0$.
These energy quantities are used only under the good-seed invariants
and are not evaluated by the algorithm.
The imbalance and endpoint bounds give
$|z_q|\le m_q$, $|z_r-z_q|\le e$, and $|z_r|\le W$.
Since $e\le2am_q/P\le m_q$, they also give $|z_r|\le2m_q$.
Expand the energy difference as
\[
 \frac{z_q^2}{m_q}-\frac{z_r^2}{W}
 =\frac{z_q^2-z_r^2}{m_q}+
   \frac{z_r^2(W-m_q)}{m_qW}
\]
The bounds $|z_q|\le m_q$ and $|z_r|\le2m_q$ make the first term
at most $3e$. If $W\le m_q$, the second term is nonpositive.
If $W>m_q$, the bounds $|z_r|\le2m_q$ and $|z_r|\le W$ give
$z_r^2/(m_qW)\le2$. The second term is then at most
$2(W-m_q)=2R_q^-$.
Consequently
\[
 K\le\sum_i E(I_i,q_i)-E(I,q)+3e+2R_q^-.
\]
At a terminal, $m=Q\ge c\ge|z|$, so its energy is at most $c$.
When we sum the energy inequality over the tree, internal energies
cancel. The remaining terminal energies sum to at most $C_{\rm term}$,
and the root energy is nonnegative.
Adding the bound for $\sum G$ proves
the first asserted inequality.

Finally, Lemma~\ref{refine:lem:comparison-costs} gives
\[
 \sum h_q\le512FJD/P+4FJD/T.
\]
Since $\widetilde C_{\rm term}\le(1+\varepsilon_{\rm leaf})C_{\rm term}\le4D$, the right side of the first asserted inequality is at most
\[
 14D+(2048F+12)JD/P+16FJD/T,
\]
which is at most $5000F(1+J/P+J/T)D$.
\end{proof}

Call a nonterminal bad if $K>bH^{-1/4}$ or $G>bH^{-1/4}$.

\begin{corollary}\label{refine:cor:badcost}
The comparison tree satisfies
\[
 \sum_{\rm bad}c\le
 10000FH^{1/4}(1+J/P+J/T)D.
\]
\end{corollary}
\begin{proof}
At a bad node, $c<2b<2H^{1/4}(K+G)$.
Sum this bound and apply Lemma~\ref{refine:lem:curvature}.
\end{proof}

In the accuracy analysis, this cost sum is multiplied by $4F/S$.
Substituting the parameters gives $O(\ell^4H^{-1/2}/\chi)D$.

\begin{lemma}\label{refine:lem:coverage}
At a non-bad comparison node with $|\widehat v|\le1-\tau$, the comparison action $\alpha^*$ belongs to the retained band.
\end{lemma}
\begin{proof}
Let $p_j^*$ be the comparison shift at cut $j$, and put
$W_j=\sum_{i\le j}m_i$.
The comparison prefix is feasible for $\Phi_j(p_j^*)$, so
$\Phi_j(p_j^*)\le W_j$.
Join a minimizing prefix to the comparison suffix at its actual target
boundary, whose mass is $M_x(l_j)+p_j^*$.
The resulting wide action has cost $\Phi_j(p_j^*)+W-W_j$.
Its cost is at least $Z$, giving
\[
 W_j-G\le\Phi_j(p_j^*)\le W_j.
\]

Shared endpoints give $p_j^*-p_0=\sum_{i\le j}z_i$.
Also $|v-\widehat v|W=|z_r|G/Z\le G$.
Weighted Cauchy--Schwarz gives
$\sum_i|z_i-vm_i|\le\sqrt{KW}$.
Write the deviation as
\[
\begin{split}
 p_j^*-p_0-\widehat v\Phi_j(p_j^*)
 ={}&\sum_{i\le j}(z_i-vm_i)+(v-\widehat v)W_j\\
    &+\widehat v\bigl(W_j-\Phi_j(p_j^*)\bigr).
\end{split}
\]
The three terms are bounded by $\sqrt{KW}$, $G$, and $G$, respectively.
Hence
\[
 |p_j^*-p_0-\widehat v\Phi_j(p_j^*)|\le\sqrt{KW}+2G.
\]
The band test allows deviation $\tau\theta b/M$.
After multiplying by $M/b$, the preceding bound is at most
\[
 \sqrt{8\ell}\,H^{-3/40}+2H^{-1/5}\le\ell^{-44}=\tau\theta,
\]
at the fixed cutoff. To check this, multiply the first term by
$\ell^{44}$ and use $\ell^{200}\le H^{1/100}$.
Its power of $H$ is at most $-3/40+44.5/20000<0$.
The second term has a smaller power of $H$.
Both terms are below half the required bound at $\chi\ge2^{1000}$.
Radius and shared endpoints follow from
Lemma~\ref{refine:lem:comparison-costs}.

Any band child state is within endpoint distance $4\theta b/M$ of $q_i$. By Lipschitzness of $w_0$ and dyadic ceiling,
\[
 h_\Sigma\le2b+2\sum_i w_0(I_i,q_i)+8L_0\theta b
 \le2b+4F\sum_i c(I_i,q_i)+8L_0\theta b\le64Fb.
\]
The rounding and imbalance guards then pass by the explicit estimates following their definition.
\end{proof}

\subsection{Sampling estimates}\label{refine:sec:imports}

For a band action put $x_t(\alpha)=\sum_i d_t(I_i,s_i(\alpha))$.

\begin{lemma}\label{refine:lem:roundedcost}
For every retained band and every $\alpha\in\cC_{b,r}$,
\[
 V_r(\beta(b,r);d_t)
 \le(1+\delta)x_t(\alpha)+(6L+2)\theta h_\Sigma,\qquad
 d_t(I_i,s_i(\beta(b,r)))\le2h_i.
\]
\end{lemma}
\begin{proof}
Two states of child $i$ in the band are within endpoint distance
$4\theta b/M$. Canonical rounding moves its endpoints by at most
$2\theta h_i$. Apply the spatial bound in
Lemma~\ref{refine:lem:invariants} and sum over children.
The connection cost introduced by rounding is at most $2\theta h_\Sigma$.
For the second bound, use the Lipschitz bound on $w_0$, the definition
of $h_i$, and $d_t\le w_0$.
\end{proof}

\begin{lemma}\label{refine:lem:group}
Condition on complete child tables. With probability at least $1-H^{-4}$, all sampled scores in one group differ from actual rounded-action costs by at most $\gamma b$. On that event, for every member $(b,r)$, every parent state $q$, and every $\alpha\in\cC_{b,r}$,
\[
 \Pi_t(q)\le L\dd(q,r)+(1+\delta)x_t(\alpha)+e_*b,
 \qquad e_*=2\gamma+512LF\theta\le\kappa\delta.
\]
\end{lemma}
\begin{proof}
For a sampled scale $h$, normalize a uniform-label summand as
\[
 Z_i=(M/b)\mathbf1_{h_i=h}d_t(I_i,s_i(\beta(b,r))).
\]
Lemma~\ref{refine:lem:roundedcost} gives
\[
 0\le Z_i\le2Mh/b,\qquad \E Z_i\le128F,\qquad
 \operatorname{Var}(Z_i)\le256FMh/b.
\]
With at least $Q_1Mh/b$ samples, Bernstein's inequality at normalized
error $\gamma/(2\ell)$ gives failure at most
$2\exp(-\ell^{37}/3000)$. Exact scales have no failure.

There are at most $\log_2(64FM)+2\le2\ell$ labels.
A cell side contains at most $1+16\theta P\le P+1$ net points, so a
group has at most $(P+1)^2$ members.
Lemma~\ref{refine:lem:parameters} bounds the full union failure by
\[
 4\ell(P+1)^2e^{-\ell^{37}/3000}\le H^{-4}.
\]
Summing the errors over labels gives score error at most $\gamma b$.

The winning actual cost is at most any member's actual cost plus
$2\gamma b$. At a member center $r$, the local rule reaches its whole cell.
Lemma~\ref{refine:lem:group-cones} transfers the bound
to every parent state $q$. Representatives in one cell differ by endpoint
distance at most $2\theta b$.
Add the resulting cone displacement to the rounded-cost bound in
Lemma~\ref{refine:lem:roundedcost}, using $h_\Sigma\le64Fb$.
Lemma~\ref{refine:lem:parameters} bounds the total error coefficient by
$2\gamma+512LF\theta\le\kappa\delta$.
The conditioning fixes all bands and child values. It leaves the fresh
group streams independent with their original laws.
\end{proof}

\begin{lemma}\label{refine:lem:online}
Fix a parent state $q$ and a local scale $b$ for which $r=r_b(q)$ has
a retained band. Condition on the complete child tables.
For every $\alpha^*\in\cC_{b,r}$ and every integer $0\le s\le S$,
the online actions satisfy
\[
 \E\sum_{t\le s}\left[
 x_t(\alpha_t)+\kappa(x_{t-1}(\alpha_t)-x_t(\alpha_t))\right]
 \le\sum_{t\le s}x_t(\alpha^*)+
 h_\Sigma\left(\frac{\ln d_{\rm act}}{\eta}+2\kappa+sH^{-3}\right).
\]
Here $d_{\rm act}$ is the number of action coordinates, and
\[
 d_{\rm act}\le M(321P)^2,\qquad \ln d_{\rm act}\le4(p_*+21)\chi.
\]
\end{lemma}
\begin{proof}
Put $f_t(i,s)=d_t(I_i,s)/h_i$.
Lemma~\ref{refine:lem:invariants} and the definition of $h_i$ put these
values in $[0,1]$ and make them nonincreasing.
Averaging all action vertices gives a positive vector in the action hull.
Its block masses equal $h_i/h_\Sigma$.
The spatial bound and band diameter give within-block cost variation
at most $\delta h_i+4L\theta b/M$.
The range of the normalized decrement is consequently at most
\[
 2\delta+8L\theta\le3\delta\le\kappa/\eta.
\]
Lemma~\ref{refine:lem:parameters} verifies the required decrement bound.
Apply Proposition~\ref{analytic:prop:entropy} with $d=d_{\rm act}$.
Proposition~\ref{analytic:prop:optimizer}, used with $h_{\rm opt}=H$,
supplies mean error at most
$H^{-3}$ at each iteration, which adds at most $sH^{-3}$ to the
normalized bound.
The conditioning fixes the child histories and leaves the parent
streams independent with their original laws.

Each wide endpoint lies in a mass interval of length $20b$.
Net separation gives at most $1+20b/g_b\le321P$ choices.
There are $M$ child blocks, so $d_{\rm act}\le M(321P)^2$.
The logarithmic estimate is Lemma~\ref{refine:lem:parameters}.
\end{proof}

\subsection{Clipped regret and refinement accuracy}

At comparison states define, for analysis only, $\overline d_t=\max(c,d_t)$ and
\[
 \mathcal R_t(I,q)=\overline d_t(I,q)-Le-\sum_i\overline d_t(I_i,q_i).
\]

\begin{lemma}\label{refine:lem:decrement}
Every nonterminal satisfies
\[
 \sum_{t=1}^S\mathcal R_t
 \le\sum_i(\overline d_0-\overline d_S)(I_i,q_i)
 \le\sum_i(m_i-|z_i|)\le W-|z_r|\le4Fc.
\]
At a non-bad node failing the velocity test, the bound is at most $(H^{-1/4}+4\tau F)c$.
\end{lemma}
\begin{proof}
The earlier-index Bellman action gives
$d_t(q)\le Le+\sum_i d_{t-1}(q_i)$.
After clipping the child values, this right side is at least $c$ by
Lemma~\ref{refine:lem:comparison-costs}. Since it bounds both $c$ and
$d_t(q)$, it also bounds
$\overline d_t(q)$. Subtracting the current clipped child sum gives
\[
 \mathcal R_t\le
 \sum_i(\overline d_{t-1}-\overline d_t)(I_i,q_i).
\]
Summing over $t$ proves the first inequality.

For a child with $d_0\le c_i$, monotonicity makes both clipped endpoints
equal to $c_i$. For any other child,
$\overline d_0=d_0\le m_i$ and $\overline d_S\ge|z_i|$.
Thus its total clipped decrement is at most $m_i-|z_i|$.
Sum over children and use $\sum_i z_i=z_r$ to obtain the remaining
general bounds.
At a non-bad node failing the velocity test,
$|z_r|/Z>1-\tau$. Since $Z=W-G\le W$, this gives
$W-|z_r|\le G+\tau W\le(H^{-1/4}+4\tau F)c$.
\end{proof}

\begin{lemma}\label{refine:lem:coveredregret}
At a non-bad comparison node passing the velocity test, conditional on complete child tables,
\[
 \E\sum_{t=1}^S\mathcal R_t
 \le h_\Sigma\left(\frac{\ln d_{\rm act}}{\eta}+2\kappa+SH^{-3}\right)
       +2\delta c+2SH^{-4}Fc.
\]
\end{lemma}
\begin{proof}
Coverage supplies $\alpha^*\in\cC_{b,r}$.
Write $A_t(\alpha)=Le+x_t(\alpha)$.
Let $t_\times$ be the first index with
$\min_{\alpha\in\cC}A_t(\alpha)<c$, or $S+1$ if there is no such index.
This crossing index is fixed by the conditioned child tables.
We bound regret before the crossing, at the crossing, and after it.

For $t<t_\times$, all action costs are at least $c$. Let $E_t$ be the group's failure event from Lemma~\ref{refine:lem:group}. We claim
\[
 \overline d_t(q)\le A_t(\alpha_t)+\kappa(x_{t-1}(\alpha_t)-x_t(\alpha_t))+2Fc\mathbf1_{E_t}.
\]
If $d_t(q)\le A_t(\alpha_t)$, clipping preserves this assertion because
$A_t(\alpha_t)\ge c$. Otherwise abbreviate $A_t=A_t(\alpha_t)$.
The sampled action is an online candidate, so $Y_t\le A_t$.
The clipped update then forces
$\Pi_t>(1+\delta)A_t$ and $d_t=\Pi_t/(1+\delta)$.
The earlier Bellman option gives $\Pi_t\le A_t+\Delta x$, where
$\Delta x=x_{t-1}(\alpha_t)-x_t(\alpha_t)$.
Consequently $\Delta x>\delta A_t\ge\delta b$.
On group success, Lemma~\ref{refine:lem:group} gives
\[
 d_t-A_t\le\frac{e_*b-\delta Le}{1+\delta}
 \le\frac{\kappa\delta b}{1+\delta}\le\kappa\Delta x.
\]
On failure the extra $2Fc$ suffices by $\overline d_t\le2Fc$.

If a crossing occurs, choose an action attaining it.
On group success, its score bound gives
$\Pi_{t_\times}<c+2\delta c$.
The comparison-child sum plus $Le$ is at least $c$, so the clipped
regret at this index is at most $2\delta c$.
Group failure adds at most $2Fc$.
At every later index, the crossing action and child monotonicity give
$\cB(d_{t-1})(q)<c$. Hence $\overline d_t(q)=c$, and its regret is
nonpositive.

Apply Lemma~\ref{refine:lem:online} to the fixed prefix before the crossing.
Its comparator cost $x_t(\alpha^*)$ is at most the clipped comparison-child
sum subtracted in $\mathcal R_t$.
The connection term $Le$ cancels.
There are at most $S$ failure indicators, each with conditional
probability at most $H^{-4}$.
Adding their expectations proves the bound.
This step uses no independence between the sampled action and group
failure.
\end{proof}

\begin{lemma}\label{refine:lem:accuracy}
Fix an integer table $U\in[0,R]$ satisfying $c\le U\le Ac$ at every
state, where $A\le A_{\rm init}$ is a power of two.
Put $a=\max(1,A/F)$.
The refinement table satisfies, at every state,
\[
 d_S\ge c/a\quad\text{on every outcome},\qquad
 \E\pos{d_S-c}\le c/\ell.
\]
The expectation fixes the entire input table $U$.
Every zero-cost state has value zero.
\end{lemma}
\begin{proof}

The lower bound is Lemma~\ref{refine:lem:invariants}.
At zero cost, $w_0=0$ gives the claimed zero value.
Now assume that the root cost $D$ is positive.
At a terminal, every clipped value equals $Q$.
Summing the definition of $\mathcal R_t$ over the comparison tree
cancels internal values and gives
\[
 \overline d_t({\rm root})=\widetilde C_{\rm term}+L\sum e+\sum\mathcal R_t.
\]
Average over $t$ and use monotonicity to bound the final root value.
The terminal term satisfies
$\widetilde C_{\rm term}-D\le(C_{\rm term}-D)+2\varepsilon_{\rm leaf}D$.
This is where the leaf approximation error enters.

For bad nodes, combine Lemma~\ref{refine:lem:decrement} with
Corollary~\ref{refine:cor:badcost}.
For other nodes failing the velocity test, use the sharper decrement bound.
For covered nodes, use Lemma~\ref{refine:lem:coveredregret},
$h_\Sigma\le64Fc$, and the action-coordinate bound in
Lemma~\ref{refine:lem:online}.
Taking conditional expectations at each node and then summing uses the
tower property. With $\sum c\le2JD$, these bounds give
{\small
\begin{align*}
 \frac{\E\pos{d_S-D}}D\le 2\varepsilon_{\rm leaf}+\mathcal E,\qquad
 \mathcal E={}&(8M+512F)J/P\\
 &+40000F^2H^{1/4}(1+J/P+J/T)/S\\
 &+\frac{2J}{S}\left[H^{-1/4}+4\tau F+
 64F\left(\frac{4(p_*+21)\chi}{\eta}+2\kappa\right)+2\delta\right]
 +132FJH^{-3}.
\end{align*}
\par}
Substitute $J\le100H/\chi$, $T\ge H^{3/4}$, $S=H/\ell^2$, and $F=\ell$.
Retaining all contributions gives the following nine-term bound:
\begin{align*}
 \mathcal E\le{}&\frac{100(8M+512\ell)H}{\chi P}
 +40000\ell^4H^{-3/4}\\
 &+\frac{4\cdot10^6\ell^4H^{1/4}}{\chi P}
 +\frac{4\cdot10^6\ell^4H^{-1/2}}{\chi}
 +\frac{200\ell^2H^{-1/4}}\chi\\
 &+\frac{26400}{\chi\ell}
 +\frac{51200(p_*+21)}{\ell^2}
 +\frac{400}{\chi\ell^8}
 +\frac{13200\ell}{\chi H^2}.
\end{align*}
The coefficient $26400$ combines $800$ from the velocity term and $25600$
from the $2\kappa$ term. The entropy contribution is
$51200(p_*+21)/\ell^2$, using the action-coordinate estimate in
Lemma~\ref{refine:lem:online}.

Each displayed term is at most $1/(16\ell)$ at the stated cutoff. For the sixth, seventh, and eighth terms, multiply by $16\ell$ and use respectively
\[
 \frac{422400}{\chi}<1,\qquad
 \frac{819200(\chi+21)}{\ell}\le
 \frac{819200(\chi+21)}{\chi^3}<1,\qquad
 \frac{6400}{\chi\ell^7}<1.
\]
For each remaining term, multiplying by $\ell$ gives at most $10^8H^{-1/5}$: use $\ell\le H^{1/20000}$, $M\le H^{1/20}$, and $P=H^{\chi+20}$. The cutoff makes this quantity smaller than $1/16$. Finally $2\varepsilon_{\rm leaf}=2\ell^{-3}\le1/(16\ell)$. Consequently
\[
 \mathcal E+2\varepsilon_{\rm leaf}\le\frac{10}{16\ell}<\frac1\ell.
\]

\end{proof}



\section{The bounded-integer approximation scheme}
\label{paper:sec:evaluation}\label{paper:sec:assembly}

Fix an integer metric $v$ whose positive values are at most $N^3$.
Apply the preceding constructions with $w=v$.
Write $R=M_v(x)+M_v(y)$ and $D_v=\WED_v(x,y)$.
Then $N\le R\le N^4$. The common parameters are
\[
 H=\operatorname{leastpow}_2\!\left(
       \left\lceil\log_2(N+R+2)\right\rceil\right),\qquad
 \chi=\log_2H,\qquad
 \ell=\operatorname{leastpow}_2(\max\{1,\chi^3\}).
\]
Use this same $H$ for every coarse entry, refinement copy, and pass.
On the large-input branch, $\chi\ge2^{1000}$ and
$\ell\varepsilon_{\rm core}\ge10$, where
$\varepsilon_{\rm core}\in(0,1)$ is the core accuracy.
Thus $p_*=\chi$, $B=H^{p_*}$, $A_{\rm init}=BH^4$, and $F=\ell$.
The remaining parameters and the labeled physical source tree are those
of the fixed refinement construction.

\subsection{The initial table on the complete state family}

\begin{lemma}\label{core:lem:initial-seed}
At the stated large-input cutoff, the coarse estimator defines a random
integer table $U^{(0)}\in[0,R]$ on the complete physical state family.
With probability at least $.99$, simultaneously at every state,
\[
 c(I,q)\le U^{(0)}(I,q)\le A_{\rm init}c(I,q).
\]
Every zero-cost state returns zero on every execution.
Shared preprocessing is deterministic and uses
$N\operatorname{poly}(\log N)$ work and polynomial space.
Conditional on the history before an entry's first request, its expected
work is at most
\[
 (1+|I|)\,2^{O(H/\chi+\chi^2)},
\]
uniformly in the target interval and without conditioning on accuracy.
Every entry computation terminates and uses polynomial space on every
fair-bit sequence.
\end{lemma}
\begin{proof}
Use Theorem~\ref{seed:thm:seed} with per-cost exponent $C=3$.
The complete source tree has at most
$2M^J\le2M\max(1,|x|)$ labeled nodes. There are at most
$(|y|+1)(|y|+2)/2$ target intervals.
Since $M\le N$, the complete state family has fewer than $N^6$ entries
at the stated cutoff. Different physical labels remain different
entries, even when their numerical source intervals agree.
This indexed family is fixed before any entry stream is exposed.

The concentration count in the proof of
Lemma~\ref{seed:lem:concentration}, with $C=3$ and $s=6$, bounds its failure
probability by
\[
 100HN^{11}e^{-H^3/768}.
\]
Here $H<8\log_2N+4$ for $N\ge2$, and $H\le N$ at the stated cutoff.
The failure probability is at most
\[
 100N^{12}e^{-H^3/768}<.01.
\]
Indeed, $N\le2^H$, and
$H^3/768>12H\ln2+\ln10000$ for $H\ge128$.
The stated cutoff is larger than this threshold.

Lemma~\ref{seed:lem:work} gives the entry overhead factor
$2^{O(H/p_*+p_*\log_2H)}=2^{O(H/\chi+\chi^2)}$.
Each entry reads only completed deterministic prefix arrays and reporting
structures. It uses its own private stream and workspace, and it does
not request another initial entry. Before its first request, that stream
is unexposed, so the conditional work bound follows from the same
calculation. Lemma~\ref{seed:lem:finite} and the fixed finite guess range
give termination and polynomial space on every execution.
\end{proof}

\subsection{The pass schedule and its accuracy}
\label{refine:sec:assembly-proof}

Set $A_0=A_{\rm init}$. While $A_{p-1}>F$, define
\[
 a_p=A_{p-1}/F,\qquad A_p=4a_p.
\]
Let $p_{\rm med}$ be the number of these median passes, and set
$p_{\rm fin}=p_{\rm med}+1$ and $s_{\rm pass}=p_{\rm fin}$.
This schedule depends only on the numerical parameters.
It is computed before any random table entry is requested.
Every $A_p$ is a power of two, and every $a_p$ used in a median pass is
an integer power of two.

In median pass $p$, apply the refinement to the preceding complete table
$U^{(p-1)}$ using $H^2$ fresh independent copies, indexed by
$c_{\rm copy}=1,\ldots,H^2$. At each state, let
\[
 U^{(p)}(I,q)=
 \min\!\left\{R,\max\!\left\{0,\
 \operatorname{os}_{H^2/2+1}
 \bigl(\lceil a_p d_S^{p,c_{\rm copy}}(I,q)\rceil:
          1\le c_{\rm copy}\le H^2\bigr)\right\}\right\},
\]
where $\operatorname{os}_k$ denotes the $k$th smallest value.
The copy streams are independent conditional on the complete preceding
table. After the last median pass, use one fresh refinement copy with input
$U=U^{(p_{\rm med})}$, factor $A=A_{p_{\rm med}}$, and $a=1$.
The mathematical core output is
\[
 Q_{\rm core}=
 \left\lfloor
   \min\{R,\max\{0,d_S^{p_{\rm fin},1}
                 (I_{\rm root},(0,|y|))\}\}
 \right\rfloor.
\]
These definitions specify complete tables. They do not require the
implementation to materialize them.

\begin{lemma}\label{core:lem:pass-count}
The deterministic schedule satisfies
\[
 s_{\rm pass}=O((p_*+1)\chi/\log\ell)
              =O(\chi^2/\log\chi)<H.
\]
\end{lemma}
\begin{proof}
Each median pass decreases $\log_2A$ by $\log_2F-2>0$.
Since $\log_2A_0=(p_*+4)\chi$, the total number of passes is at most
\[
 2+\frac{(p_*+4)\chi}{\log_2F-2}.
\]
Use $p_*=\chi$ and $\chi^3\le F<2\chi^3$.
The resulting bound is smaller than $H=2^\chi$ at the stated cutoff.
\end{proof}

\begin{lemma}\label{core:lem:scheme-accuracy}
On the large-input branch,
\[
 \Prb\{D_v\le Q_{\rm core}
                 \le(1+\varepsilon_{\rm core})D_v\}>4/5.
\]
\end{lemma}
\begin{proof}
If $D_v=0$, the initial table and every subsequent table return zero at
the root on every outcome. Assume $D_v>0$.
Condition on a complete preceding table satisfying
$c\le U^{(p-1)}\le A_{p-1}c$ at every state.
At a fixed state $(I,q)$, put $c_q=c(I,q)$.
Lemma~\ref{refine:lem:accuracy} gives
$d_S^{p,c_{\rm copy}}\ge c_q/a_p$ on every outcome.
Thus each scaled value $\lceil a_p d_S^{p,c_{\rm copy}}\rceil$
is at least $c_q$.
For a positive-cost state, the same lemma and Markov's inequality give
\[
 \Prb\{d_S^{p,c_{\rm copy}}>2c_q\}\le1/\ell.
\]
When $d_S^{p,c_{\rm copy}}\le2c_q$, the integrality of $a_p$ and $c_q$
implies $\lceil a_p d_S^{p,c_{\rm copy}}\rceil\le2a_pc_q$.
The probability that at least half the copies fail is at most
\[
 2^{H^2}\ell^{-H^2/2}\le e^{-H^2}.
\]
If fewer than half fail, the upper middle order statistic is at most
$2a_pc_q\le A_pc_q$.
Clamping preserves the lower bound because $c_q\le R$.
Zero-cost states stay zero.

The state count is at most $N^6\le e^{6H\ln2}<e^{10H}$, so one
transition fails with probability at most $e^{10H-H^2}$.
This union bound does not require independence across states.
Apply it to the first failed transition after a good preceding table.
Lemma~\ref{core:lem:initial-seed} and
Lemma~\ref{core:lem:pass-count} bound the probability of any required bad
seed by
\[
 .01+H e^{10H-H^2}<1/10.
\]

At the end, $A_{p_{\rm med}}\le F$.
Conditional on any fixed good final seed,
Lemma~\ref{refine:lem:accuracy} applies to the fresh final copy with $a=1$.
Its root value is at least the integer $D_v$, and
\[
 \Prb\{d_S^{p_{\rm fin},1}>
             (1+\varepsilon_{\rm core})D_v\}
 \le\frac1{\ell\varepsilon_{\rm core}}\le1/10.
\]
Clamping cannot increase the value and preserves its lower bound.
Flooring also preserves that bound because $D_v$ is an integer.
Combining the two failure bounds proves the assertion.
\end{proof}

\subsection{The coherent cache implementation}

Use the canonical identities and private owners of
Definition~\ref{refine:def:keys} in the fresh core-run namespace $\nu$.
A cache record has status absent, active, or complete.
Completed payloads and their stored owner decisions are immutable.
The program draws bits only when an owner is first reached.
Shared prefix arrays, nets, and reporting structures are deterministic.

Put $B_{\rm idx}=T+S$. Assign the scalar ranks
\begin{align*}
 \operatorname{rk}(U^{(0)})&=1,\\
 \operatorname{rk}(w_j^{p,c_{\rm copy}})
   &=(p-1)(B_{\rm idx}+2)+2+j,\\
 \operatorname{rk}(d_t^{p,c_{\rm copy}})
   &=(p-1)(B_{\rm idx}+2)+2+T+t,\\
 \operatorname{rk}(U^{(p)})&=p(B_{\rm idx}+2)+1.
\end{align*}
Here $c_{\rm copy}$ is the copy index, and $d_0=w_T$.
A mean has the rank of $w_T$, a band the rank of $d_1$, and a group
or online record at index $t$ the rank of $d_t$.
A leaf record has rank zero.
Give scalar seed, warmup, and refinement records phase $2$,
group and online records phase $1$, and mean, band, and leaf records
phase $0$.
Order keys lexicographically by rank, remaining physical height, and
phase.

A median seed reads same-node final refinement records of rank one less.
At a physical leaf, a warmup or refinement record reads its shared
rank-zero leaf record. At an internal node, $w_0$ reads the preceding
seed at its query and reported centers, all of rank one less.
The local helper reads same-node $w_0$ values. Warmup propagation and
the refinement's Bellman fallback read lower-rank child values.
Means read same-node warmup indices below $T$, and bands read lower-rank
child means and warmup values.
A scalar refinement record reads same-node group and online records
with smaller phase. These records read same-node bands with smaller
rank, or smaller phase at index one, and lower-rank earlier child
histories. Their current-index scalar reads have equal rank and smaller
physical height. Every persistent read decreases the displayed key order.

\begin{wedalgorithm}{Coherent lazy evaluation and batched scalar reads}{alg:get-gather}
\Require Complete keys in one core-run namespace and their defining routines.
\Ensure Immutable completed records and one \textsc{Get} per distinct canonical scalar key in a gather.
\Procedure{GET}{$K$}
  \State Canonicalize $K$, including the supplied-input identity and $d_0=w_T$.
  \If{$\mathsf{Cache}[K]$ is complete}
    \State \Return its stored payload.
  \EndIf
  \State \textbf{assert} that $\mathsf{Cache}[K]$ is absent.
  \State Mark $\mathsf{Cache}[K]$ active.
  \State Run its defining routine, using its private stream when needed.
  \State Store its payload and owner decisions immutably; mark it complete.
  \State \Return its stored payload.
\EndProcedure
\Statex
\Procedure{GATHER}{$\mathcal L$}
  \State Canonicalize, sort, and deduplicate the scalar keys in $\mathcal L$.
  \State Initialize an empty local map $\mathsf{val}$.
  \For{each distinct key $K$, in sorted order}
    \State $\mathsf{val}[K]\gets\Call{GET}{K}$.
  \EndFor
  \State \Return $\mathsf{val}$.
\EndProcedure
\end{wedalgorithm}
\FloatBarrier


Wide lists, prefix minima, representative lists, Bellman scores,
gathered arrays, and optimizer iterates are local scratch data.
Every repeated scratch read is charged.
An initial seed key runs the standalone initial-entry routine.
A median seed key gathers its $H^2$ final-copy values and applies the
order statistic and clamps already defined.

\begin{lemma}\label{core:lem:cache}
GET never encounters an active record. It terminates on every
fair-bit sequence and returns the complete-table value for its key.
The resulting output has the distribution of $Q_{\rm core}$.
\end{lemma}
\begin{proof}
Every recursive call decreases the displayed key order.
Its ranks, heights, and phases are nonnegative integers, so recursive
evaluation is well founded. The standalone initial routine is finite
by Lemma~\ref{seed:lem:finite}.
All other local enumerations have fixed finite domains.
Proposition~\ref{analytic:prop:optimizer} gives a finite optimizer and a
sampler with a bounded number of fair bits.
Thus every defining routine terminates after its dependencies terminate.
An active key cannot be reached from itself through decreasing keys.

Induction on the same order shows that each completed payload equals
its mathematical record with the same owner stream.
The program never redraws a completed owner or creates a new child
value for a new caller. Coupling its reached owner streams with the
potential streams of the complete construction proves the
distributional assertion.
\end{proof}

Accuracy conditions on complete child histories, including unreached
entries. Those histories exclude current parent group and online streams.
Work conditions instead on the exposed execution before a first request.
An absent initial entry has exposed none of its private stream.
The realized cache and request trace are not the conditioning objects
in the accuracy argument.

\subsection{Finite arithmetic and unconditional work}
\label{refine:sec:work}

\subsubsection{Finite returned values and local computation}

Every table used as a seed is clamped to $[0,R]$.
On any such seed, the local $w_0$ routine returns zero, a cone value,
or default $R$. Its tested centers are at mass sup-distance at most
$U(I,q)\le R$, and $\widetilde U\le R$. Consequently
\[
 0\le w_0\le(2L_0+1)R.
\]
All connection and cone charges are integers.
Warmup values are integers, and means have denominator $T$.
Every refinement update is nonnegative and at most its local $w_0$,
irrespective of seed accuracy.

Let $Q_{\rm den}=\ell^{10}+1$, and let $h(I)$ be remaining physical
height. Induction on height, for all indices together, shows that the
denominator of every returned $d_t(I,q)$ divides $Q_{\rm den}^{h(I)}$.
Leaves return the integer $Q$.
A parent candidate adds an integer charge to child values, or uses an
integer warmup value. Minima and maxima select arguments.
The only new denominator comes from dividing by
$1+\delta=Q_{\rm den}/\ell^{10}$.
Thus returned values have
$O(H+J\log\ell)=O(H)$ bits.
Sampling scores and optimizer mixtures choose an action; they do not
average recursive costs into a returned value.

All band and rounded-state guards are exact rational comparisons.
Their inputs, and the online linear coefficients, have
$H^{O(p_*+1)}$ bits.
The dyadic labels satisfy $b/M\le h_i\le h_\Sigma\le64Fb$.
Every $Mh_i$ is an integer, so the action vectors share the integer
denominator $Mh_\Sigma$, of $O(H)$ bits.
Each wide endpoint has at most $321P$ candidates.
Thus the optimizer dimension is at most $M(321P)^2$, and its input
list has at most $(321P)^{2M}$ actions.
Proposition~\ref{analytic:prop:optimizer} gives deterministic local bit
work and storage bounded by
\[
 2^{O(M(p_*+1)\chi)}H^{O(p_*+1)}
 =2^{O(H^{1/20}\chi^2+\chi^2)}
 =2^{o(H\log\chi/\chi)}.
\]
The numerical iterations reuse gathered earlier-index values.
Only the sampled whole action requests current-index child values.

\subsubsection{Distinct recursive dependencies and repeated reads}

The initial scale ratio is $O(A^3)$, and its radius-to-net-spacing
ratio is $O(PA^2)$. Later local scale ratios are $128aF$.
Both scale counts are $O((p_*+1)\chi)=O(\chi^2)=O(\ell)$.
In a later search, $w_0(I,q)/b\le16F$.
The query square meets only $O((1+F/\theta)^2)$ whole cells per scale.
Each cell contains $O(P^2)$ representatives.
Binary searches in the net arrays generate the lists without scanning
unrelated target positions. Cell indices have $O(H)$ bits.

Each child has $O(P^2)$ wide states.
Gathering their distinct earlier-index values and histories, over all
local representatives, cells, and scales, uses
$D_{\rm earlier}=H^{O(p_*+1)}$ scalar inputs.
Action enumeration has $(O(P))^{2M}$ local steps.
Prefix minima are evaluated only at the finite candidate boundary shifts
and use the same gathered values.

For current-index requests, fix a group, scale label $h$, and child.
Rounded states use the common net spacing
$D_h=\max(1,\theta h/16)$.
Their first endpoint ranges over mass length $O(b)$.
After fixing it, the guard $|z|\le2h$ confines the second endpoint
to mass length $4h$. Net separation bounds their number by
\[
 O((1+b/D_h)(1+h/D_h))
 =O(\theta^{-2}(1+b/h)).
\]
If $k_h<M$, then $k_h\le2Q_1Mh/b$ and $h/b<1/Q_1$.
Multiplying the state count by the number of sampled labels gives
$O(M\theta^{-2}Q_1)$ distinct inputs for this scale label.
If $k_h=M$, then $b/h<2Q_1$, and summing all labels gives the same bound.
There are $O(\ell)$ labels and $\ell^{O(1)}$ local groups.
Evaluating a group winner adds $M$ calls per group.
Each online draw adds $M$ calls per local scale.
Hence one refinement entry has at most
\[
 D_{\rm same}=M\ell^{O(1)}
\]
same-pass, same-index child dependencies.
The test $h_i=h$ precedes each group request, so a zero indicator does
not create a dependency.

These bounds count distinct dependencies.
A local dictionary associates each key with its completed scalar value,
including values first needed after a winner is selected.
Repeated scoring reads keep their multiplicities but reuse those values.
There are $O(P^2)$ representatives per group, at most $M$ sampled labels
per scale label, and $\ell^{O(1)}$ groups and labels.
Direct group scoring uses at most
$MP^2\ell^{O(1)}$ lookups, besides the earlier-value reads and selected
actions. Key construction and each dictionary lookup have
polynomial-in-$H$ bit cost.
Thus all repeated lookups fit within the local
$H^{O(p_*+1)}$ budget.
Full action enumeration and optimizer iterations have the separate local
bound already proved; they do not create one recursive call per action.

\subsubsection{The dependency count and source-width sum}

For this count, fix the potential owner streams and expand each auxiliary
record into its scalar requests, retaining its realized owner decisions.
The bounds above hold for every such decision.
In the resulting scalar request graph, every equal-rank edge descends
one physical child. Every other edge decreases rank, including the
$H^2$ calls of a median seed. Its fanout is absorbed in $D_{\rm earlier}$.
Put $K_{\rm dep}=s_{\rm pass}(T+S+2)+1$.
Lemma~\ref{core:lem:pass-count} and $T=o(S)$ give
\[
 K_{\rm dep}
 =O\!\left(\frac{H(p_*+1)\chi}{\ell^2\log\ell}\right)
 =O\!\left(\frac H{\chi^4\log\chi}\right).
\]
A path ending at physical depth $j$ has at most $j$ equal-rank edges
and at most $K_{\rm dep}$ rank-decreasing edges.
No edge moves to a parent or an unrelated source node.
Unfolding the dependency graph bounds its evaluations at depth
$j$ by
\[
 (j+1)(K_{\rm dep}+1)2^{j+K_{\rm dep}}
 M^j\ell^{C_1j}H^{C_2(p_*+1)K_{\rm dep}}
\]
for absolute constants $C_1,C_2$.
The logarithm beyond the factor $M^j$ is at most
\begin{align*}
 O\bigl(J\log\ell+J+K_{\rm dep}
       +(p_*+1)K_{\rm dep}\log H\bigr)
 &=O\!\left(\frac{H\log\chi}{\chi}
             +\frac H{\chi^2\log\chi}\right)\\
 &=O(H\log\chi/\chi).
\end{align*}
Thus at most $M^j2^{O(H\log\chi/\chi)}$ evaluations occur at that depth
on every execution, including inaccurate seeds.

A physical interval at depth $j$ has width at most
$\lceil|x|/M^j\rceil$.
At each first initial-entry request, condition on its exposed history
and charge Lemma~\ref{core:lem:initial-seed}.
The stream of that absent entry is still private.
Summing over the deterministic request-slot bound, and adding the local
work, gives the source-width factor
\[
 \sum_{j=0}^J M^j\bigl(1+\lceil|x|/M^j\rceil\bigr)
 =O((|x|+1)(J+M)).
\]
Here $M^J\le M|x|$ on the nontrivial source branch.
The factors $J+M$ and the number of depths contribute only
polynomial-in-$H$ overhead.

The logarithm of all remaining overhead is bounded by
\[
 O\!\left(
   \frac{H\log\chi}{\chi}
   +\frac H{\chi^2\log\chi}
   +\frac H\chi+\chi^2+H^{1/20}\chi^2
 \right)
 =O(H\log\chi/\chi).
\]
The middle terms include the concrete seed work.
The last term includes action enumeration and finite optimization.
Net preprocessing uses $O(NH\log R)$ binary work.
The leaf reporting structure and the seed preprocessing use
$N\operatorname{poly}(\log N)$ work.
Leaf queries use $O(\ell^3H\log N)$ binary work and $O(\ell^3)$ metric
queries. These costs, and a fixed uniform deterministic polylogarithmic
overhead per metric query, fit within the same bound.
The resulting unconditional expected work is
\[
 (N+1)\,2^{O(H\log\chi/\chi)}
 =N^{1+O(\log\log\log N/\log\log N)}.
\]
The conditioning used here is the exposed first-request history.
No charge is conditioned on a good complete seed.

There are polynomially many possible persistent keys.
Their descriptions have $O(H)$ bits, and balanced dictionaries add only
polynomial-in-$H$ lookup cost.
Each band list and optimizer workspace has the finite bounds already
proved. Together with Lemma~\ref{seed:lem:finite}, the stored records and
local scratch use polynomial space on every execution.

\subsection{The integer-core entry point}

\begin{wedalgorithm}{The bounded-integer entry point}{alg:integer-core}
\Require Strings $x,y$, integer metric oracle $v$ with positive costs at most $N^3$, and canonical accuracy encoding $\sigma_{\rm core}$ for $\varepsilon_{\rm core}\in(0,1)$.
\Ensure An integer output in $[0,R]$; every exact branch returns $D_v$.
\Procedure{INTEGER-CORE}{$x,y,v,\sigma_{\rm core}$}
  \State $N\gets|x|+|y|$; test equality by zero metric values at corresponding occurrences.
  \If{$x=y$} \State \Return $0$ \EndIf
  \If{either string is empty or $|x|\le1$}
    \State \Return \Call{EXACT-DP}{$x,y,v$}.
  \EndIf
  \State Query every occurrence's gap cost and form both prefix-mass arrays.
  \State $R\gets M_v(x)+M_v(y)$.
  \State $H\gets\operatorname{leastpow}_2(\lceil\log_2(N+R+2)\rceil)$.
  \State Inspect at most $H^2+12$ bits of $\sigma_{\rm core}$.
  \If{this prefix does not contain the complete encoding}
    \State \Return \Call{EXACT-DP}{$x,y,v$}.
  \EndIf
  \State Decode $\varepsilon_{\rm core}$ from the inspected prefix.
  \State $\chi\gets\log_2H$; $\ell\gets\operatorname{leastpow}_2(\max\{1,\chi^3\})$.
  \If{$\chi<2^{1000}$ or $\ell\varepsilon_{\rm core}<10$}
    \State \Return \Call{EXACT-DP}{$x,y,v$}.
  \EndIf
  \State Compute the fixed refinement parameters using this common $H$.
  \State Build the deterministic reporting structures and nested mass nets.
  \State Define the labeled physical source tree and state domain implicitly.
  \State Allocate a fresh core-run namespace $\nu$ and its implicit private-owner rule.
  \State Compute the median schedule and $p_{\rm fin}$.
  \State Initialize an empty cache in namespace $\nu$.
  \State $K_{\rm root}\gets\wedkey{d_S^{p_{\rm fin},1}}{I_{\rm root},(0,|y|)}$.
  \State $Q_{\rm root}\gets\Call{GET}{K_{\rm root}}$.
  \State \Return $\lfloor\min\{R,\max\{0,Q_{\rm root}\}\}\rfloor$.
\EndProcedure
\end{wedalgorithm}

\FloatBarrier

\begin{theorem}\label{integer:thm:scheme}

INTEGER-CORE is one uniform fair-bit algorithm for canonically encoded
rational accuracies $\varepsilon_{\rm core}\in(0,1)$.
For every such fixed accuracy, every pair of strings of total length
$N$, and every integer metric $v$ with positive costs at most $N^3$,
it returns an integer $Q\in[0,R]\subseteq[0,N^4]$ satisfying
\[
 \Prb\{D_v\le Q\le(1+\varepsilon_{\rm core})D_v\}>4/5.
\]
It returns zero on equal strings.
It terminates on every fair-bit sequence and uses polynomial workspace
on every execution.
Its large-input expected work, including deterministic preprocessing,
all random sampling, and finite binary computation, is
\[
 (N+1)2^{O(H\log\chi/\chi)}
 =N^{1+O(\log\log\log N/\log\log N)},
\]
uniformly over the allowed metrics, strings, and decoded accuracies
that pass the branch tests. The constant in this large-branch bound is
absolute and does not depend on the accuracy encoding.
In particular, for every fixed accuracy encoding there are a deterministic
function $r_{\varepsilon_{\rm core}}(N)\to0$ and a constant
$K_{\varepsilon_{\rm core}}$, independent of the metric and strings, such
that for $N\ge1$ the expected work is at most
$K_{\varepsilon_{\rm core}}N^{1+r_{\varepsilon_{\rm core}}(N)}$.
The same assertions permit a fixed uniform deterministic polylogarithmic
overhead per metric query.
Repeated complete runs use independent owner streams.
\end{theorem}
\begin{proof}
On the large-input branch, Lemma~\ref{core:lem:cache} identifies the
implemented output with the complete construction.
Lemma~\ref{core:lem:scheme-accuracy} gives its approximation probability.
The unconditional work and workspace bounds have been proved from the
concrete initial table and finite optimizer.
Since $N\le R\le N^4$, we have $H=\Theta(\log N)$ and
$\chi=\Theta(\log\log N)$, which gives the displayed exponent.
The bounds are uniform over all allowed metrics, including executions
with inaccurate initial or intermediate tables.

The remaining branches call exact dynamic programming or return the
exact zero answer. The bounded accuracy inspection reads only
$O(H^2)$ bits before any accuracy-dependent test.
If it succeeds, decoding and testing use polynomial-in-$H$ work.
For each fixed encoding, the inspection and numerical cutoff conditions
eventually cease as $N$ grows.
These inspection and numerical-cutoff branches occur only in a bounded
range of lengths for that accuracy and can be absorbed in
$K_{\varepsilon_{\rm core}}$.
Its final floor is finite rational division.
Fresh namespaces and independent owner streams give the stated
independence of repeated runs.
\end{proof}

\begin{remark}\label{integer:rem:unitcost}
For unit-cost edit distance on explicitly stored strings with
$O(\log N)$-bit integer symbols, INTEGER-CORE gives a
$(1+\varepsilon)$ approximation with probability greater than $4/5$
for every fixed rational $\varepsilon\in(0,1)$. Its expected running
time on the word RAM with $O(\log N)$-bit words is
\[
 N^{1+O(\log\log\log N/\log\log N)}.
\]
Indeed, the unit-cost metric $v(a,b)=\mathbf1_{\{a\ne b\}}$ on the
alphabet together with $\gap$ has positive costs equal to one, so
Theorem~\ref{integer:thm:scheme} applies. Metric queries reduce to
symbol equality tests. The core already charges all finite binary
computation and sampling, so the same bound holds on the word RAM
without normalization or exact real arithmetic.
\end{remark}

\FloatBarrier

\section{The approximation scheme for arbitrary metric costs}\label{paper:sec:general}

Return to the original metric $w$ and put $D=\WED_w(x,y)$ and
$U_0=M_w(x)+M_w(y)$. The normalized metric from
Section~\ref{paper:sec:normalization} has bounded integer costs.
We apply the bounded-integer scheme with fresh randomness and transfer
its answer to the original metric.

\subsection{Uniform algorithm, confidence, and expected work}

\begin{theorem}\label{wrapper:thm:wrapper}
There is one uniform fair-bit algorithm in the stated exact-real model
that takes a canonically encoded rational $\eps\in(0,1]$ as input.
For each fixed accuracy, its worst-case expected running time is
$N^{1+o(1)}$, uniformly over all original metric costs and their ratios. It uses
polynomial workspace on every outcome. It returns a value $\widehat D$
with
\[
\Prb\{D\le\widehat D\le(1+\eps)D\}
\ge\frac{99}{100}\frac{2944}{3125}>\frac56.
\]
It returns zero exactly for equal strings, on every outcome. Its value is
always in $[0,U_0]$. In addition,
\[
\Prb\{D\le\widehat D\}\ge
\frac{2944}{3125}>\frac{11}{12},
\]
including all bracket outcomes.
\end{theorem}
\begin{proof}
First test equality and handle empty inputs. Their answers are exact.
We also handle $|x|\le1$ by exact linear-time dynamic programming.
Define $H_0$ to be the least power of two at least
$\lceil\log_2(N+2)\rceil$. Inspect at most $H_0^2+8$ bits of the
accuracy encoding. If it is not complete within that prefix, compute
$D$ by exact dynamic programming without reading the remaining suffix.
This fallback is valid without knowing the accuracy value.
Otherwise the inspected numerator and denominator have
$O(H_0^2)$ bits. Test $\eps N<361$ by integer cross-multiplication,
and use exact dynamic programming if the test holds. These steps precede
the construction of $\eta_{\rm norm}$ or any accuracy-dependent integer search range.
For each fixed accuracy encoding, both fallback conditions eventually
cease. Rational control calculations and conversion of the inspected
integers to exact real registers use only $\operatorname{poly}(\log N)$
work. They can be implemented by bit scans and repeated doubling.

In the remaining branch run Lemma~\ref{wrapper:lem:bracket} once. Construct the
normalized oracle from Lemma~\ref{wrapper:lem:normalize}, and set
$\varepsilon_{\rm core}=\eps/3$. Canonically encode this rational as $\sigma_{\rm core}$. Run the complete core 5 times with fresh independent
randomness. Only deterministic preprocessing may be shared across these
runs without an additional conditional-independence argument. Let $Q$
be the median of the 5 outputs. Since the present strings are distinct,
$\bar D\ge1$. Clamp the median to $[1,\bar R]$, and return
\[
\widehat D=\min\{U_0,\eta_{\rm norm}\,\min\{\bar R,\max\{1,Q\}\}\}.
\]
Both clamps preserve a successful median: $1\le\bar D\le\bar R$, and
$D\le U_0$. They also ensure a strictly positive finite output on
distinct strings and an output at most $U_0$ on every outcome.

\emph{Confidence.} Condition on the complete bracket history
$\mathcal B$, including its revealed randomness. For every such history,
the normalized metric is fixed, satisfies the core promise, and is
independent of the fresh core streams. Each run succeeds with probability
at least $4/5$. Its success indicator is independent of the others under
this conditioning. The probability that at most two runs succeed is at
most
\[
\sum_{j=3}^{5}\binom{5}{j}(1/5)^j(4/5)^{5-j}
=\frac{181}{3125}<\frac1{12}.
\]
Thus, uniformly in $\mathcal B$, the median lies in
$[\bar D,(1+\eps/3)\bar D]$ with probability at least
$p_5=2944/3125$.
On this median event, Lemma~\ref{wrapper:lem:normalize} gives
$\widehat D\ge D$, even if the bracket's lower endpoint is invalid.
On its intersection with $L_{\rm br}\le D$, the good-bracket bound in Lemma~\ref{wrapper:lem:normalize} gives
\[
\widehat D\le(1+\eps/3)^2D\le(1+\eps)D,
\]
because $2\eps/3+\eps^2/9\le7\eps/9\le\eps$ for $\eps\le1$.
Integrating the uniform conditional median probability over the good
bracket histories gives success at least $(99/100)p_5>5/6$.
This product bound does not require unconditional independence of the
bracket event and the median event.

\emph{Uniform expected work, including failed brackets.}
On every bracket history, the normalized metric has positive integer
costs at most $N^3$. Theorem~\ref{integer:thm:scheme} gives the same
conditional expected-work bound
$K_{\eps/3}N^{1+r_{\eps/3}(N)}$ for every run, without conditioning on
either approximation event. A normalized query costs $O(\log N)$ original
oracle operations, and the remaining control arithmetic has only
polylogarithmic overhead. The tower property consequently bounds the
total expected core work, over all bracket outcomes, by
\[
5K_{\eps/3}N^{1+r_{\eps/3}(N)}\operatorname{polylog}N
=N^{1+o(1)}.
\]
The bracket costs deterministically $O(N\log^2N)$. Gap-prefix arrays,
reporting structures, and endpoint nets belong to the one normalized
metric. Their construction is included in Theorem~\ref{integer:thm:scheme} and can be
shared when deterministic. The inspection and numerical-cutoff fallbacks
occur only in a bounded range of lengths for each fixed accuracy, so
they preserve the pointwise asymptotic bound.

\emph{Workspace and finite randomness.}
The bracket has the deterministic register and bit bounds in
Lemma~\ref{wrapper:lem:bracket}. The normalized oracle stores finitely many
original real values and bounded integer control data. Core arrays contain
only bounded integers and the core's finite rationals. Execute the 5
runs sequentially, retaining only their scalar answers. The
polynomial core workspace and exact dynamic-programming fallbacks imply
polynomial workspace on every outcome. The bracket uses a bounded number
of fair bits; the core uses its own charged fair-bit implementation.
\end{proof}

\begin{wedalgorithm}{The original-metric wrapper}{alg:metric-scheme}
\Require Strings $x,y$, original metric oracle $w$, and canonical accuracy encoding $\sigma$.
\Ensure The scalar wrapper output, with the stated exact fallbacks.
\Procedure{METRIC-SCHEME}{$x,y,\sigma$}
  \State Test $x=y$ by oracle equality on the original occurrence handles.
  \If{$x=y$} \State \Return $0$ \EndIf
  \If{either string is empty or $|x|\le1$}
    \State \Return \Call{EXACT-DP}{$x,y,w$}
  \EndIf
  \State $N\gets|x|+|y|$.
  \State $H_0\gets\operatorname{leastpow}_2(\lceil\log_2(N+2)\rceil)$.
  \State Inspect at most $H_0^2+8$ bits of $\sigma$.
  \If{this prefix does not contain the complete encoding}
    \State \Return \Call{EXACT-DP}{$x,y,w$}
  \EndIf
  \State Decode $\eps$ from the inspected encoding; set $\eps\gets\min\{\eps,1\}$.
  \If{$\eps N<361$} \State \Return \Call{EXACT-DP}{$x,y,w$} \EndIf
  \State Query all original occurrence gap costs; set $U_0\gets M_w(x)+M_w(y)$.
  \State $(L_{\rm br},U_{\rm br})\gets\Call{BRACKET}{x,y,w}$.
  \State $\eta_{\rm norm}\gets\eps L_{\rm br}/(3N)$; $\varepsilon_{\rm core}\gets\eps/3$.
  \State Define the oracle $\bar w$ by the bounded integer search in Lemma~\ref{wrapper:lem:normalize}.
  \State Compute normalized gap costs and set $\bar R\gets M_{\bar w}(x)+M_{\bar w}(y)$.
  \State Canonically encode $\varepsilon_{\rm core}$ as $\sigma_{\rm core}$.
  \For{$\nu=1,\ldots,5$}
    \State Allocate a fresh random-owner namespace for run $\nu$.
    \State $Q_\nu\gets\Call{INTEGER-CORE}{x,y,\bar w,\sigma_{\rm core}}$.
    \State Release this run's private workspace and retain $Q_\nu$.
  \EndFor
  \State $Q\gets$ the third smallest of $Q_1,\ldots,Q_{5}$.
  \State \Return $\min\{U_0,\eta_{\rm norm}\min\{\bar R,\max\{1,Q\}\}\}$.
\EndProcedure
\end{wedalgorithm}


\subsection{Proof of the main theorem}
\begin{proof}
Theorem~\ref{integer:thm:scheme} supplies the bounded-integer algorithm and its quantitative expected-work bound. Theorem~\ref{wrapper:thm:wrapper} supplies the normalization, approximation guarantee, and output properties. A normalized query has polylogarithmic overhead, and the wrapper makes a constant number of core runs. Thus its expected running time is
\[
 N^{1+O(\log\log\log N/\log\log N)}.
\]
The workspace and termination guarantees follow from the same two theorems. The exact branches cover all remaining input lengths for each fixed accuracy.
\end{proof}

For a fixed rational accuracy $\varepsilon>1$, use accuracy $1$. The resulting factor $2$ is at most $1+\varepsilon$.


\subsection{Accuracy chosen from the input length}

\begin{proof}[Proof of Corollary~\ref{wrapper:cor:vanishing}]
Let $H_0$ be the wrapper parameter, which depends only on $N$, and put
\[
 \chi_0=\log_2H_0,\qquad
 \ell_0=\operatorname{leastpow}_2(\max\{1,\chi_0^3\}),\qquad
 \varepsilon_N=\min\{1,30/\ell_0\}.
\]
Run the wrapper with the canonical encoding of $\varepsilon_N$.
This choice precedes the bracket and normalization, so it has no
dependence on the normalized costs.

For sufficiently large $N$, we have $\ell_0\ge32$,
$\varepsilon_N=30/\ell_0$, and
$\varepsilon_{\rm core}=10/\ell_0$.
The normalized metric satisfies $N\le\bar R\le N^4$ on every bracket
outcome once $\varepsilon_NN\ge361$.
Since $\bar R\ge N$, its core parameter satisfies $H\ge H_0$.
Hence $\chi\ge\chi_0$ and $\ell\ge\ell_0$, giving
$\ell\varepsilon_{\rm core}\ge10$.

The accuracy encodings pass both bounded inspections for sufficiently
large $N$. To see this explicitly, write $\ell_0=2^k$.
The reduced fractions are $\varepsilon_N=15/2^{k-1}$ and
$\varepsilon_{\rm core}=5/2^{k-1}$.
Their code lengths are $2k+6$ and $2k+4$, respectively.
Since $k=O(\log\chi_0)$, these lengths fit within the wrapper's
$H_0^2+8$ and the core's $H^2+12$ inspection limits.
Also $\ell_0=O((\log\log N)^3)$, so
$\varepsilon_NN\ge361$ eventually holds.
Finally, $\chi\ge\chi_0$ implies that the absolute cutoff
$\chi\ge2^{1000}$ eventually holds on every bracket outcome.
Thus the inspection and numerical-cutoff fallbacks occur only in a
bounded range of input lengths,
with an absolute bound independent of the original metric.

On the large branch, accuracy affects only the normalization, the
encoding inspections, and the final Markov test. The cost cap leaves
the integer search range at $N^3$. The coarse estimator, refinement
parameters, pass schedule, and dependency count depend on $H$, not
on $\varepsilon_N$. The absolute large-branch bound in
Theorem~\ref{integer:thm:scheme} applies uniformly:
\[
 (N+1)2^{O(H\log\chi/\chi)}\operatorname{polylog}N
 =N^{1+O(\log\log\log N/\log\log N)}.
\]
The bracket work and accuracy construction fit within this bound.
The bounded range of these fallbacks is absorbed into one absolute
constant. This argument uses the large-branch estimate directly,
rather than the fixed-accuracy constants of that theorem.

The confidence calculation of Theorem~\ref{wrapper:thm:wrapper} applies
to this accuracy. On its success event the transferred factor is
$(1+\varepsilon_N/3)^2\le1+\varepsilon_N$.
For $N\ge4$, $\chi_0\ge\log_2\log_2N$, so
\[
 \varepsilon_N\le\frac{30}{(\log_2\log_2N)^3}.
\]
This proves the approximation bound. Inputs with $N<4$ are handled
exactly.
\end{proof}

\paragraph{AI assistance.}
OpenAI's Codex assisted with proof development, parameter calculations,
literature research, and drafting. Anthropic's Claude assisted with
reviewing proofs, proposing parameter refinements, comparing related
work, and refining the exposition.

\bibliographystyle{unsrturl}
\bibliography{ref}
\end{document}
